\documentclass[12pt,oneside]{book}
\providecommand{\blindreview}{0}

\usepackage[spanish,es-nodecimaldot]{babel}
\usepackage[utf8]{inputenc}
\usepackage[T1]{fontenc}
\usepackage{newtxtext,newtxmath}
\usepackage{microtype}
\setlength{\emergencystretch}{2em}
\usepackage{geometry}
\geometry{letterpaper,top=2.7cm,bottom=2.7cm,left=3.2cm,right=2.7cm}
\usepackage{setspace}
\onehalfspacing
\usepackage{csquotes}
\usepackage{amsmath,mathtools}
\usepackage{booktabs,tabularx,longtable,array,multirow}
\usepackage{graphicx}
\usepackage{xcolor}
\usepackage{enumitem}
\usepackage{tikz}
\usetikzlibrary{arrows.meta,positioning,shapes.geometric,fit,backgrounds}
\usepackage{pgfplots}
\pgfplotsset{compat=1.18}
\usepackage{caption}
\usepackage{subcaption}
\usepackage{float}
\usepackage{fancyhdr}
\usepackage{titlesec}
\usepackage{tocloft}
\setlength{\cftsecnumwidth}{3.2em}
\usepackage{hyperref}
\usepackage[nameinlink,noabbrev,spanish]{cleveref}
\usepackage[backend=biber,style=apa,sorting=nyt,maxcitenames=2]{biblatex}
\addbibresource{referencias.bib}
% Generado por sync-editorial.mjs; no editar.
\newcommand{\COBRIAAutor}{Billy Jak Cordero Porras}
\newcommand{\COBRIALugar}{Coto Brus, Costa Rica}
\newcommand{\COBRIAAnio}{2026}
\newcommand{\COBRIAVersionEcosistema}{1.0.0}
\newcommand{\COBRIAVersionCore}{1.0}


\definecolor{cobriaBlue}{gray}{0}
\definecolor{cobriaGreen}{gray}{0}
\definecolor{cobriaGold}{gray}{0}
\definecolor{cobriaLight}{gray}{0.96}

\hypersetup{
  hidelinks,
  pdftitle={COBRIA: diseño y evaluación de contratos reconstruibles para sistemas NoSQL documentales},
  pdfauthor={\ifnum\blindreview=1 Autor anonimizado\else \COBRIAAutor\fi},
  pdfsubject={Documento maestro de investigación de la propuesta independiente COBRIA},
  pdfkeywords={COBRIA, NoSQL documental, proyecciones reconstruibles, arquitectura de software, procedencia}
}

\pagestyle{fancy}
\setlength{\headheight}{14.5pt}
\fancyhf{}
\fancyhead[L]{\small COBRIA}
\fancyhead[R]{\small\nouppercase{\leftmark}}
\fancyfoot[C]{\thepage}
\renewcommand{\headrulewidth}{0.4pt}

\titleformat{\chapter}[display]
  {\normalfont\bfseries}
  {\filleft\Large\chaptername\ \thechapter}
  {1ex}
  {\titlerule[0.6pt]\vspace{1.5ex}\huge}

\newcommand{\COBRIA}{\textsf{COBRIA}}
\newcommand{\pendiente}[1]{\textit{[Pendiente: #1]}}

\begin{document}
\frontmatter
\hypersetup{pageanchor=false}
\begin{titlepage}
\centering
\vspace*{2.2cm}
{\Large Documento maestro de investigación\par}
\vspace{1.5cm}
{\Huge\bfseries COBRIA\par}
\vspace{0.6cm}
{\LARGE Diseño y evaluación de contratos reconstruibles\par}
\vspace{0.35cm}
{\Large para sistemas NoSQL documentales\par}
\vspace{1.2cm}
\rule{0.72\textwidth}{0.7pt}
\vfill
{\large \COBRIAAutor\par}
\vspace{0.35cm}
{\normalsize Investigador independiente\par}
{\normalsize \COBRIALugar\par}
\vspace{1cm}
{\normalsize Edición de investigación · \COBRIAAnio\par}
\end{titlepage}

\hypersetup{pageanchor=true}
\setcounter{page}{1}

\chapter*{Declaración de alcance y originalidad}
\addcontentsline{toc}{chapter}{Declaración de alcance y originalidad}
Este manuscrito propone, formaliza y evalúa una integración arquitectónica denominada
\COBRIA. No se atribuye como invención propia ninguno de los patrones consolidados que
la componen, entre ellos Repositorio, Mapeador de datos, Capa de servicios, casos de uso, vistas
materializadas, separación de modelos de lectura y escritura o generación dirigida por
metadatos. La contribución buscada reside en su articulación explícita alrededor de
entidades canónicas, proyecciones reconstruibles, aislamiento multiinquilino y preparación
trazable para minería de datos e inteligencia artificial, así como en su evaluación
mediante una implementación neutral y dos perfiles de caso anonimizados.

Los casos se describen de forma neutral como \emph{Sistema A} y \emph{Sistema B}. Los
experimentos emplearon datos sintéticos o semirrealistas y configuraciones reproducibles.
No se publicarán nombres comerciales, credenciales, datos personales ni detalles que
permitan identificar organizaciones o personas.

\chapter*{Resumen}
\addcontentsline{toc}{chapter}{Resumen}
Los sistemas NoSQL permiten construir aplicaciones flexibles y distribuidas, pero
trasladan al diseño de la aplicación decisiones complejas sobre desnormalización,
índices, consistencia, aislamiento entre organizaciones y evolución del esquema. Estas
dificultades aumentan cuando los mismos datos deben alimentar consultas operacionales,
procesamiento de eventos, minería de datos e inteligencia artificial. Este trabajo
presenta \COBRIA{} (\emph{Capa Orientada a Bases Reconstruibles para Inteligencia
Analítica}), una arquitectura dirigida por metadatos que distingue entidades
canónicas de representaciones físicas y proyecciones derivadas. La propuesta organiza
el acceso mediante repositorios, servicios y casos de uso; clasifica las proyecciones
según sus requisitos de consistencia; y exige que todo dato derivado declare su origen,
versión y mecanismo de reconstrucción. La investigación siguió ciencia del diseño y
examinó la propuesta en tres aplicaciones documentales anonimizadas. La evaluación
comprende inspección de contratos, pruebas NoSQL, reconstrucción de proyecciones,
aislamiento por ámbito, evolución de esquemas y recuperación autorizada para analítica e
inteligencia artificial. El repositorio conserva 270 corridas históricas distribuidas
sobre MongoDB, CouchDB y OpenSearch, además de una repetición preliminar del arnés
corregido en MongoDB y CouchDB. Las corridas históricas no se presentan como
certificación vigente porque una auditoría posterior fortaleció el oráculo. Las
mediciones SQL permanecen segregadas y no sostienen afirmaciones de rendimiento NoSQL.
Permanecen pendientes la repetición distribuida confirmatoria, la evaluación humana de
transferencia y la revisión independiente.

\noindent\textbf{Palabras clave:} NoSQL; datos canónicos; proyecciones reconstruibles;
multiinquilinato; vistas materializadas; minería de datos; inteligencia artificial;
arquitectura dirigida por metadatos.

\tableofcontents
\listoffigures
\listoftables

\mainmatter
\part{Fundamentos de la investigación}
\chapter{Introducción}
\label{ch:introduccion}

\section{Contexto}

Los sistemas de información contemporáneos deben atender simultáneamente necesidades
que antes solían resolverse en plataformas separadas. Una misma aplicación puede
registrar operaciones de negocio, reaccionar a eventos en tiempo real, conservar un
historial auditable, aislar información de varias organizaciones y producir datos para
analítica o inteligencia artificial. Esta convergencia ha aumentado el valor de las
bases de datos NoSQL, cuya flexibilidad permite representar estructuras heterogéneas y
distribuir grandes volúmenes de información. Sistemas influyentes como Bigtable
demostraron que un modelo distribuido orientado a columnas puede soportar cargas
muy diferentes cuando el diseño físico se alinea con los patrones de acceso
\parencite{chang2006bigtable}. Dynamo, por su parte, expuso de forma explícita los
compromisos entre disponibilidad, particionamiento, replicación y consistencia en un
almacenamiento distribuido orientado a clave \parencite{decandia2007dynamo}.

Esa flexibilidad no elimina la necesidad de modelar; la desplaza. En sistemas NoSQL,
la aplicación debe decidir qué datos se almacenan juntos, cuáles se duplican, qué
índices se mantienen y cómo se reparan las representaciones derivadas. El esquema puede
ser flexible en la base, pero sigue existiendo en el código, en las rutas, en las reglas
de acceso y en las expectativas de cada consumidor.

La dificultad aumenta en aplicaciones \emph{multiinquilino}, es decir, aquellas en las
que una misma plataforma sirve a varias organizaciones o contextos lógicos. La
consolidación puede reducir costos, pero introduce requisitos estrictos de aislamiento,
configuración, rendimiento y recuperación. La literatura ha mostrado que no existe un
único diseño físico que sea óptimo para todos los organizaciones y cargas; las alternativas
intercambian escalabilidad, espacio, flexibilidad y complejidad operativa
\parencite{aulbach2009flexibleschemas,yaish2014adaptiveschema,rico2016multitarget}.

Al mismo tiempo, los datos operacionales han dejado de ser únicamente un registro del
pasado. Se utilizan para detectar anomalías, descubrir secuencias, construir modelos
predictivos, recuperar evidencias y asistir decisiones humanas. Los sistemas de
aprendizaje automático necesitan características consistentes entre entrenamiento e
inferencia, versiones reproducibles y trazabilidad hacia las fuentes. Los almacenes de
características surgieron precisamente para organizar este flujo y reducir la
fragmentación entre preparación, entrenamiento y servicio de modelos
\parencite{orr2021featurestores,li2023featurestore}. Sin embargo, una capa analítica no
puede corregir por sí sola un origen operacional ambiguo o lleno de duplicaciones sin
propietario.

Este trabajo parte de una observación práctica: los sistemas NoSQL de alcance amplio
necesitan distinguir con claridad qué representación conserva el estado autorizado del
negocio y cuáles representaciones existen para acelerar consultas, integrar módulos o
alimentar procesos analíticos. A partir de esa separación se formula \COBRIA{}, una
arquitectura que organiza entidades canónicas, mapeos físicos, repositorios, servicios,
casos de uso y proyecciones reconstruibles dentro de una misma disciplina de diseño.

\section{Problema de investigación}

En un modelo desnormalizado, una modificación lógica puede afectar varias ubicaciones.
Por ejemplo, actualizar el estado de una entidad puede requerir modificar su registro
principal, retirar una referencia de un índice, añadirla a otro y actualizar un resumen.
Si estas ubicaciones se tratan como fuentes equivalentes, el sistema pierde una base
clara para decidir cuál valor es correcto cuando aparece una discrepancia. Si, por el
contrario, se reconoce una fuente canónica y las demás se definen como funciones
derivadas, la discrepancia puede detectarse y repararse.

La vista materializada es un antecedente directo de esta idea: una representación
derivada puede mejorar el acceso, pero debe mantenerse frente a cambios en los datos
base. La investigación clásica sobre mantenimiento incremental estudió cómo calcular
los cambios de una vista sin reconstruirla completamente
\parencite{gupta1993incremental,gupta1995materialized}. En aplicaciones NoSQL, el mismo
problema reaparece en índices manuales, resúmenes, relaciones inversas, contadores,
colas de salida y conjuntos de datos analíticos, aunque no siempre se describa con esa
terminología.

Se identifican seis manifestaciones principales del problema:

\begin{enumerate}[label=\textbf{P\arabic*.},leftmargin=*]
  \item \textbf{Ambigüedad de autoridad.} Varias copias contienen atributos similares,
  pero no se declara cuál es la fuente autorizada.

  \item \textbf{Consultas no acotadas.} Ante la ausencia de índices apropiados, el
  cliente recupera colecciones grandes y filtra localmente, aumentando transferencia,
  memoria, latencia y exposición innecesaria.

  \item \textbf{Derivaciones frágiles.} Los índices se actualizan mediante lógica
  dispersa y carecen de versión, propietario, verificador o procedimiento de
  reconstrucción.

  \item \textbf{Acoplamiento semántico--físico.} Nombres, claves compactas, rutas y
  reglas se duplican en múltiples capas; una evolución del esquema puede dejar
  consumidores incompatibles.

  \item \textbf{Aislamiento incompleto.} El organización se aplica como filtro de interfaz,
  pero no como dimensión obligatoria del modelo, la consulta y la autorización.

  \item \textbf{Preparación analítica no reproducible.} Los conjuntos de datos y características
  se crean con exportaciones o scripts independientes, sin linaje suficiente hacia los
  datos operacionales y sin corrección temporal verificable.
\end{enumerate}

Estas manifestaciones están relacionadas. Una consulta amplia suele motivar la creación
de un índice; el índice introduce duplicación; la duplicación exige mantenimiento; el
mantenimiento necesita autoridad y trazabilidad; y la trazabilidad se vuelve todavía
más importante cuando el resultado alimenta un modelo de IA. El problema de
investigación se formula, por tanto, de la siguiente manera:

\begin{quote}
¿Cómo diseñar y evaluar una arquitectura NoSQL multiinquilino que preserve una fuente
canónica comprensible, permita optimizar consultas mediante proyecciones derivadas,
pueda reconstruir esas proyecciones de forma verificable y ofrezca datos trazables para
minería de datos e inteligencia artificial?
\end{quote}

\section{Motivación}

\subsection{Motivación técnica}

El rendimiento de una aplicación conectada a una base remota no depende solamente del
tiempo de CPU. Importan también el número de solicitudes, los bytes transferidos, la
cantidad de registros procesados por el cliente y la variabilidad de la red. Una
proyección pequeña puede reducir el conjunto candidato de $N$ entidades a $K$
resultados relevantes, donde $K \ll N$. No obstante, esta mejora de lectura incrementa
el trabajo de escritura y almacenamiento. COBRIA busca hacer visible ese intercambio,
en vez de tratar cada índice como una optimización gratuita.

\subsection{Motivación de mantenibilidad}

En sistemas extensos, una entidad puede estar representada por un modelo de dominio, un
serializador, un repositorio, un servicio, reglas de validación, índices, catálogos y
documentación. Cuando esos artefactos evolucionan de manera independiente, aparecen
errores difíciles de detectar. Un enfoque dirigido por metadatos puede reducir la
duplicación de decisiones, siempre que no convierta los archivos generados en una falsa
medida de funcionalidad. COBRIA propone evaluar la generación por su capacidad para
preservar contratos y reducir defectos, no por la cantidad de código producido.

\subsection{Motivación analítica}

La minería de datos requiere transformar eventos y entidades en series, agregaciones,
secuencias, grafos o vectores de características. Cuando esas transformaciones se
consideran proyecciones versionadas, pueden reconstruirse desde datos canónicos,
compararse entre versiones y auditarse. Este enfoque conecta la arquitectura operacional
con prácticas de Almacén de características y MLOps, pero evita asumir que cualquier proyección
operacional es automáticamente adecuada para entrenar modelos.

\subsection{Motivación social y regional}

\ifnum\blindreview=1
El nombre COBRIA incorpora simbólicamente el origen regional de la investigación,
omitido durante la revisión ciega:
\else
El nombre COBRIA incorpora de manera simbólica la identidad de Coto Brus, Costa Rica:
\fi
\textbf{CO}--\textbf{BR}--\textbf{IA}. La referencia territorial no cambia el carácter
general del artefacto; expresa que la investigación tecnológica puede originarse fuera
de los centros tradicionales y producir conocimiento susceptible de evaluación y
reutilización internacional. El estudio evitará, sin embargo, convertir esa identidad
en una afirmación técnica o en evidencia de validez.

\section{La propuesta COBRIA en síntesis}

COBRIA significa \emph{Capa Orientada a Bases Reconstruibles para Inteligencia
Analítica}. La palabra \emph{objeto} se utiliza en un
sentido amplio: una entidad semántica identificable, no necesariamente una clase ligada
a un lenguaje de programación específico.

La arquitectura propone cinco movimientos principales:

\begin{enumerate}[leftmargin=*]
  \item representar el estado autorizado mediante entidades canónicas;
  \item separar nombres semánticos de la representación física almacenada;
  \item encapsular persistencia, reglas y operaciones mediante repositorios, servicios
  y casos de uso;
  \item definir índices, resúmenes, relaciones y conjuntos de datos como proyecciones con contrato
  de reconstrucción;
  \item conservar el linaje necesario para que minería de datos e IA utilicen
  representaciones versionadas, autorizadas y temporalmente correctas.
\end{enumerate}

\begin{figure}[H]
\centering
\resizebox{0.94\textwidth}{!}{%
\begin{tikzpicture}[
  node distance=7mm and 12mm,
  box/.style={draw=black,rounded corners=1pt,fill=white,
              minimum width=5.6cm,minimum height=9mm,align=center},
  derived/.style={draw=black,rounded corners=1pt,fill=black!6,
                  minimum width=4.4cm,minimum height=9mm,align=center},
  arrow/.style={-{Latex[length=2.4mm]},semithick,black}
]
  \node[box] (ui) {Interfaces y consumidores};
  \node[box,below=of ui] (uc) {Casos de uso y autorización};
  \node[box,below=of uc] (svc) {Servicios, políticas y composición};
  \node[box,below=of svc] (repo) {Repositorios / Mapeadores de datos};
  \node[box,below=of repo,fill=black!10] (canon) {Entidades canónicas};
  \node[derived,below left=of canon] (oper) {Proyecciones operacionales};
  \node[derived,below right=of canon] (anal) {Proyecciones analíticas y de IA};
  \draw[arrow] (ui) -- (uc);
  \draw[arrow] (uc) -- (svc);
  \draw[arrow] (svc) -- (repo);
  \draw[arrow] (repo) -- (canon);
  \draw[arrow] (canon) -- node[left,font=\scriptsize]{deriva} (oper);
  \draw[arrow] (canon) -- node[right,font=\scriptsize]{deriva} (anal);
  \draw[-{Latex[length=2.1mm]},dashed,black] (oper) to[bend right=24]
    node[below,font=\scriptsize]{verificación y reconstrucción} (canon);
  \draw[-{Latex[length=2.1mm]},dashed,black] (anal) to[bend left=24]
    node[below,font=\scriptsize]{linaje y reconstrucción} (canon);
\end{tikzpicture}
}
\caption{Vista conceptual inicial de COBRIA. Las proyecciones optimizan consumidores
específicos, pero su autoridad procede de las entidades canónicas.}
\label{fig:cobria-general}
\end{figure}

La flecha de reconstrucción en la \cref{fig:cobria-general} no implica que una
proyección sobrescriba libremente la entidad canónica. Indica que la proyección puede
compararse con el resultado esperado de su función de derivación y regenerarse cuando
sea necesario. Las correcciones del estado de negocio deben realizarse por medio de un
caso de uso autorizado.

\section{Vacío de investigación}

COBRIA se apoya en conceptos conocidos. Repositorio y Mapeador de datos separan el dominio de
la persistencia; las vistas materializadas aceleran consultas; CQRS distingue modelos
de lectura y escritura; las arquitecturas multiinquilino estudian la consolidación y el
aislamiento; los Almacenes de características organizan características para aprendizaje automático;
y la ciencia del diseño ofrece un método para construir y evaluar artefactos
\parencite{hevner2004designscience,peffers2007dsrm}.

El vacío que motiva este trabajo no es la inexistencia de esas ideas, sino la falta de
una articulación y evaluación conjunta para aplicaciones NoSQL basadas en estructuras
JSON que:

\begin{itemize}[leftmargin=*]
  \item trate la reconstruibilidad como contrato de toda proyección, no solo como una
  tarea de mantenimiento ocasional;
  \item conecte el modelo semántico y la representación física compacta mediante un
  mapeo versionado y reversible;
  \item incorpore el organización en identidad, rutas, consultas, autorización y pruebas;
  \item incluya proyecciones operacionales y analíticas bajo una taxonomía común;
  \item mida simultáneamente beneficios de lectura y amplificaciones de escritura y
  almacenamiento;
  \item preserve linaje desde la entidad operacional hasta características,
  predicciones y retroalimentación humana;
  \item sea contrastada en más de un tipo de sistema y exponga también los casos en que
  su complejidad no se justifica.
\end{itemize}

Por tanto, la originalidad pretendida es \emph{composicional y evaluativa}. El estudio
no afirmará que COBRIA sustituye todas las alternativas documentales ni que su
complejidad sea preferible para cualquier problema.

\section{Síntesis ampliada de trabajos relacionados}

La procedencia de datos antecede a COBRIA: la literatura distingue el origen de un
resultado y las partes de la fuente responsables de él \parencite{buneman2001provenance}.
COBRIA adopta esa preocupación, pero la lleva al contrato operacional de cada proyección:
fuente, versión, transformación, alcance y procedimiento verificable de reconstrucción.
Por ello, el linaje no se presenta como una invención, sino como una condición para que
una copia optimizada pueda auditarse y regenerarse.

La evolución de esquemas NoSQL también es un problema conocido. La ausencia de un
esquema global puede trasladar heterogeneidad y migraciones al código de aplicación, y
se han propuesto lenguajes declarativos para administrar esos cambios sin modificar el
motor \parencite{scherzinger2013schema}. COBRIA coincide en explicitar versiones, pero
añade la separación entre semántica canónica, mapeo físico y derivados reconstruibles.

En analítica, la arquitectura lakehouse busca reunir almacenamiento abierto, gobierno,
inteligencia de negocios y aprendizaje automático \parencite{armbrust2021lakehouse}.
Su escala es la plataforma de datos; COBRIA opera en otra capa: contratos de dominio y
proyecciones que pueden alimentar una plataforma analítica sin convertirla en autoridad
del estado transaccional. Ambas ideas pueden coexistir.

Finalmente, los sistemas de datos autoajustables estudian el uso del historial de carga
para planificar cambios físicos \parencite{pavlo2017selfdriving}. Esto delimita una línea
futura para COBRIA: recomendar o retirar proyecciones con un presupuesto observado. La
presente evaluación no implementa un optimizador autónomo; mide primero los beneficios y
costos que una política de ese tipo necesitaría comparar.

En conjunto, la revisión sitúa la novedad posible en la integración y en sus contratos,
no en renombrar vistas materializadas, procedencia, CQRS, repositorios o evolución de
esquema. La contribución es refutable porque puede fallar si la reconstrucción no es
exacta, si el aislamiento se rompe o si la amplificación supera el ahorro de lectura.

\section{Objetivos}

\subsection{Objetivo general}

Diseñar, formalizar y evaluar COBRIA como una arquitectura dirigida por metadatos,
basada en entidades canónicas y proyecciones reconstruibles, para determinar su efecto
sobre la eficiencia, integridad, mantenibilidad, recuperabilidad, aislamiento
multiinquilino y preparación para minería de datos e inteligencia artificial en sistemas
NoSQL.

\subsection{Objetivos específicos}

\begin{enumerate}[label=\textbf{OE\arabic*.},leftmargin=*]
  \item Definir los conceptos, componentes, responsabilidades y reglas de dependencia
  de COBRIA con terminología independiente de lenguajes y proveedores.

  \item Formalizar entidades canónicas, mapeos semántico--físicos, proyecciones,
  contratos de reconstrucción, alcance multiinquilino y linaje analítico.

  \item Clasificar las proyecciones en críticas, operacionales, analíticas y efímeras,
  y asociar a cada clase requisitos de consistencia y recuperación.

  \item Analizar la complejidad temporal y espacial de lectura, escritura, hidratación,
  mantenimiento y reconstrucción.

  \item Cuantificar latencia, transferencia de datos y amplificación de lectura,
  escritura y almacenamiento frente a alternativas sin proyecciones explícitas.

  \item Evaluar la exactitud y el costo de reconstrucciones totales e incrementales,
  incluida la detección de referencias huérfanas y proyecciones faltantes.

  \item Medir el efecto de metadatos compartidos sobre la evolución del esquema, la
  consistencia entre artefactos y el esfuerzo de mantenimiento.

  \item Evaluar el aislamiento multiinquilino mediante rutas acotadas, reglas de acceso y
  pruebas negativas con datos sintéticos.

  \item Analizar la capacidad de COBRIA para producir conjuntos de datos, series temporales y
  características versionadas con linaje y corrección temporal.

  \item Examinar la integración segura de modelos y agentes de IA mediante casos de
  uso autorizados, evidencia, consentimiento, auditoría y retroalimentación humana.

  \item Validar la aplicabilidad del artefacto en dos casos de estudio anónimos de
  diferente escala y naturaleza operacional.

  \item Identificar límites, antipatrones y condiciones bajo las cuales COBRIA resultaría
  innecesaria o excesivamente compleja.
\end{enumerate}

\section{Preguntas de investigación}

\begin{table}[H]
\centering
\caption{Preguntas de investigación de COBRIA.}
\label{tab:preguntas}
\begin{tabularx}{\textwidth}{>{\bfseries}p{1.6cm}X}
\toprule
Código & Pregunta \\
\midrule
RQ1 & ¿Cómo afecta COBRIA la complejidad y el comportamiento observado de las
operaciones de lectura, escritura y reconstrucción? \\
RQ2 & ¿En qué medida las proyecciones reducen la latencia, los registros procesados y
los bytes transferidos en consultas frecuentes? \\
RQ3 & ¿Qué amplificación de escritura y almacenamiento introduce el mantenimiento de
proyecciones y cuándo se justifica ese costo? \\
RQ4 & ¿Con qué exactitud, tiempo y consumo de recursos pueden reconstruirse las
proyecciones a partir de entidades canónicas? \\
RQ5 & ¿Cómo influye un mapeo semántico--físico dirigido por metadatos en la evolución
del esquema y en la consistencia entre modelos, validadores, reglas e índices? \\
RQ6 & ¿En qué medida el ámbito obligatorio mejora el aislamiento y limita las consultas
en un sistema multiinquilino? \\
RQ7 & ¿Puede la misma arquitectura adaptarse a un sistema orientado a eventos y
multimedia y a una plataforma empresarial con un dominio mucho más extenso? \\
RQ8 & ¿En qué medida las proyecciones analíticas versionadas mejoran la
reproducibilidad, el linaje y la corrección temporal de conjuntos de datos y características? \\
RQ9 & ¿Qué mecanismos arquitectónicos son necesarios para que modelos y agentes de IA
utilicen datos autorizados, auditables y acompañados de evidencia? \\
RQ10 & ¿Cuáles son los límites prácticos y las señales de sobrearquitectura de COBRIA? \\
\bottomrule
\end{tabularx}
\end{table}

Las preguntas combinan propiedades técnicas y de ingeniería. RQ1--RQ4 se responderán
principalmente mediante mediciones; RQ5 y RQ6 combinarán análisis de artefactos con
pruebas; RQ7 utilizará comparación entre casos; RQ8 y RQ9 abordarán preparación y
gobernanza de IA; y RQ10 sintetizará resultados negativos y costos de adopción.

\section{Hipótesis}

\begin{longtable}{>{\bfseries}p{1.6cm}p{5.2cm}p{7.1cm}}
\caption{Hipótesis y evidencia prevista.}\label{tab:hipotesis}\\
\toprule
Código & Hipótesis & Indicadores principales \\
\midrule
\endfirsthead
\toprule
Código & Hipótesis & Indicadores principales \\
\midrule
\endhead
H1 & Las consultas apoyadas en proyecciones procesan una cantidad de datos más cercana
al número de resultados que al tamaño total de la colección. & Registros leídos, bytes,
latencia p50/p95/p99 y amplificación de lectura. \\
H2 & La reducción del costo de lectura se obtiene a cambio de mayor amplificación de
escritura y almacenamiento. & Escrituras físicas por operación lógica, tamaño de
canónicos y derivados, duración de escritura. \\
H3 & Un contrato explícito permite reconstruir proyecciones deterministas con alta
exactitud a partir del estado canónico. & Cobertura de proyección, referencias
huérfanas, igualdad esperada--observada y tiempo de reconstrucción. \\
H4 & La reconstrucción incremental resulta más eficiente que la total cuando el número
de entidades modificadas es pequeño respecto al conjunto canónico. & Tiempo, lecturas,
escrituras y memoria para $\Delta N$ frente a $N$. \\
H5 & La definición compartida de metadatos reduce discrepancias entre modelo lógico,
representación física, validación, reglas e índices. & Defectos detectados, artefactos
modificados manualmente, tiempo de evolución y cobertura de contrato. \\
H6 & La incorporación obligatoria del organización en rutas, consultas y autorización reduce
la superficie de exposición cruzada. & Pruebas negativas, consultas fuera de ámbito,
registros innecesarios recuperados y violaciones detectadas. \\
H7 & Las proyecciones analíticas versionadas mejoran la reproducibilidad y el linaje de
conjuntos de datos frente a exportaciones ad hoc. & Tiempo de preparación, repetibilidad,
porcentaje de campos con procedencia y diferencias entre ejecuciones. \\
H8 & COBRIA puede aplicarse a los dos casos anónimos sin cambiar sus principios
centrales, aunque requiera proyecciones y políticas diferentes. & Correspondencia de
componentes, adaptaciones necesarias y resultados por caso. \\
H9 & El beneficio neto disminuye cuando una entidad tiene pocos registros, pocas
consultas o demasiadas proyecciones respecto a su valor funcional. & Función de costo,
presupuesto de proyección y análisis de sensibilidad. \\
\bottomrule
\end{longtable}

Las hipótesis son deliberadamente refutables. H2 y H9 reconocen costos y evitan evaluar
la arquitectura únicamente con métricas favorables. La investigación considerará
valioso identificar escenarios donde COBRIA no produzca una mejora neta.

\section{Contribuciones esperadas}

\subsection{Contribución conceptual}

Una definición neutral de COBRIA y un vocabulario que distinga entidad canónica,
representación física, proyección, índice, resumen, conjunto de datos y artefacto efímero. Esta
distinción busca facilitar la discusión entre desarrollo operacional, ingeniería de
datos y aprendizaje automático.

\subsection{Contribución formal}

Un modelo de entidades, ámbitos, mapeos y proyecciones; propiedades de autoridad,
reconstruibilidad, idempotencia y linaje; y condiciones bajo las cuales una proyección
puede considerarse verificable.

\subsection{Contribución arquitectónica}

Una integración reproducible de modelos de dominio, Repositorio/Mapeador de datos, Servicio
Capa, casos de uso, metadatos, proyecciones reconstruibles y componentes de analítica e
IA. El aporte arquitectónico incluye reglas de dependencia y una clasificación de
proyecciones por consistencia.

\subsection{Contribución algorítmica}

Pseudocódigo y posteriormente implementaciones de referencia para sincronización,
reconstrucción total e incremental, detección de inconsistencias, migración, generación
de conjuntos de datos y verificación de linaje.

\subsection{Contribución empírica}

Un conjunto de experimentos sobre dos casos anónimos que mida tanto beneficios como
costos. Se publicarán configuraciones y datos sintéticos suficientes para permitir la
replicación sin revelar información privada.

\subsection{Contribución práctica}

Una guía de adopción, un modelo de madurez, un presupuesto de proyecciones y criterios
para decidir cuándo utilizar COBRIA, cuándo simplificarla y cuándo preferir otra
organización documental.

\section{Casos de estudio y neutralidad del análisis}

La investigación se apoya en dos implementaciones existentes que comparten una familia
de decisiones arquitectónicas, pero difieren en dominio, escala y comportamiento. Para
evitar que el estudio se convierta en documentación de productos particulares, se
adoptan descripciones neutrales.

\begin{table}[H]
\centering
\caption{Caracterización inicial y anónima de los casos.}
\label{tab:casos}
\begin{tabularx}{\textwidth}{>{\bfseries}p{3.2cm}XX}
\toprule
Dimensión & Sistema A & Sistema B \\
\midrule
Naturaleza & Orientado a eventos, relaciones contextuales y artefactos multimodales. &
Plataforma empresarial multiinquilino con dominio extenso y operaciones transaccionales.\\
Escala estructural & Decenas de entidades y servicios especializados. & Cientos de
entidades, catálogos, schemas y proyecciones.\\
Datos destacados & Eventos temporales, secuencias, texto, métricas y multimedia. &
Entidades operacionales, financieras, configurables, de auditoría e integración.\\
Procesamiento & Servicios servicios internos, colas, procesos de trabajo y generación de artefactos. &
Operaciones autorizadas, índices, procesos de trabajo, políticas y generación por metadatos.\\
Valor para el estudio & Adaptabilidad multimodal, temporal, analítica y de IA. &
Escalabilidad estructural, multiinquilinato, gobierno del esquema y reconstrucción.\\
\bottomrule
\end{tabularx}
\end{table}

Los conteos exactos se fijarán en el protocolo experimental mediante scripts
reproducibles y una fecha de corte. No se utilizará el tamaño del código como sustituto
de calidad o completitud. Cuando una capacidad solo esté modelada o documentada, pero no
verificada mediante pruebas, se declarará como tal.

\section{Alcance}

El estudio incluye:

\begin{itemize}[leftmargin=*]
  \item sistemas NoSQL cuyo modelo de aplicación pueda representarse mediante
  documentos, árboles JSON o estructuras clave--valor jerárquicas;
  \item entidades canónicas y representaciones derivadas para consulta;
  \item repositorios, servicios y casos de uso como fronteras de responsabilidad;
  \item aislamiento lógico multiinquilino;
  \item mantenimiento y reconstrucción de índices y proyecciones;
  \item metadatos para modelos, esquemas, reglas, validación y documentación;
  \item complejidad computacional y mediciones de rendimiento;
  \item proyecciones analíticas, series temporales, características y linaje;
  \item integración gobernada de modelos y agentes de IA;
  \item dos casos de estudio anonimizados y una implementación neutral de referencia.
\end{itemize}

\section{Delimitaciones}

El estudio no pretende:

\begin{itemize}[leftmargin=*]
  \item demostrar que un motor NoSQL específico es superior a todos los demás;
  \item reemplazar los patrones Repositorio, CQRS, vistas materializadas, Almacén de características
  o arquitectura por capas;
  \item evaluar todos los motores NoSQL disponibles;
  \item probar que una cantidad elevada de entidades o archivos implica mayor madurez;
  \item presentar una implementación de IA como evidencia suficiente de exactitud,
  equidad o seguridad;
  \item entrenar modelos con información personal o sensible;
  \item revelar nombres, reglas comerciales, credenciales o datos de los sistemas
  utilizados como casos;
  \item afirmar causalidad más allá de lo que permita el diseño experimental.
\end{itemize}

Los experimentos iniciales se realizarán en entornos controlados y con datos sintéticos.
Las diferencias entre emuladores y producción se tratarán como amenaza a la validez
externa. Las integraciones externas no se considerarán funcionales únicamente porque
existan modelos o contratos en el código.

\section{Unidad de análisis y niveles de observación}

La unidad principal de análisis es una \emph{operación de datos COBRIA}: una acción que
lee o modifica una entidad canónica y, cuando corresponde, consulta o mantiene sus
proyecciones. El estudio observará cuatro niveles:

\begin{enumerate}[leftmargin=*]
  \item \textbf{Nivel de entidad:} identidad, campos, versión, organización e invariantes.
  \item \textbf{Nivel de operación:} lecturas, escrituras, transacciones, reintentos y
  efectos sobre proyecciones.
  \item \textbf{Nivel de sistema:} latencia, transferencia, almacenamiento,
  concurrencia, aislamiento y recuperación.
  \item \textbf{Nivel analítico:} conjunto de datos, característica, modelo, predicción, evidencia
  y retroalimentación.
\end{enumerate}

Esta separación evita mezclar conclusiones. Por ejemplo, que un acceso a caché tenga
costo promedio constante no implica que la operación distribuida completa sea
instantánea; y que un conjunto de datos sea reconstruible no implica que sea adecuado o ético
para entrenar un modelo.

\section{Consideraciones éticas iniciales}

La anonimización de los sistemas no es suficiente si los datos conservan combinaciones
que permitan reidentificar personas. Por ello, los experimentos utilizarán datos
sintéticos y no exportarán información de salud, identidad, pagos, conversaciones,
ubicaciones precisas ni contenido multimedia privado. Las pruebas de aislamiento se
realizarán únicamente en entornos autorizados.

En los componentes de IA se aplicarán, desde el diseño, los principios de trazabilidad,
supervisión humana, control de acceso y evaluación de riesgo. Marcos oficiales como el
marco de gestión de riesgos de inteligencia artificial de NIST enfatizan que la gestión del riesgo debe acompañar
el ciclo de vida completo del sistema, no limitarse al modelo entrenado
\parencite{nist2023airmf}. COBRIA no considerará una respuesta generada como una fuente
canónica de negocio mientras no pase por el caso de uso y las confirmaciones requeridas.

\section{Organización del documento maestro}

Después de esta introducción, el documento se organizará en siete bloques:

\begin{enumerate}[leftmargin=*]
  \item \textbf{Fundamentos y literatura:} conceptos necesarios y comparación con
  trabajos relacionados.
  \item \textbf{Diseño de COBRIA:} componentes, modelo formal, algoritmos, consistencia
  y reconstrucción.
  \item \textbf{Rendimiento:} complejidad, amplificación, modelo de costos y
  optimización.
  \item \textbf{Analítica e IA:} minería, Almacén de características, multimodalidad, RAG, agentes y
  gobernanza.
  \item \textbf{Implementación y casos:} implementación neutral y caracterización
  anónima de los dos sistemas.
  \item \textbf{Evaluación:} protocolo, resultados, análisis estadístico, discusión y
  amenazas a la validez.
  \item \textbf{Síntesis:} guía de adopción, conclusiones y líneas futuras.
\end{enumerate}

La estructura seguirá la lógica de la ciencia del diseño: el problema conduce a un
artefacto; el artefacto se formaliza e implementa; y su utilidad se juzga mediante
evaluación explícita, no por afirmación del autor \parencite{hevner2004designscience,
peffers2007dsrm}.

\section{Comparación sistemática con arquitecturas relacionadas}

La revisión ampliada distingue coincidencia de mecanismos y coincidencia de objetivos.
CQRS separa modelos de lectura y escritura, pero no obliga por sí mismo a declarar
linaje, ámbito o reconstrucción. Las vistas materializadas anticipan consultas y pueden
regenerarse desde sus fuentes, aunque su gobierno suele quedar dentro del motor
\parencite{microsoft2025cqrs,microsoft2025materialized}.
El abastecimiento por eventos conserva una secuencia autoritativa capaz de rehidratar
estado y proyecciones, pero COBRIA permite que la autoridad sea una entidad actual
versionada y no exige un registro completo de eventos \parencite{fowler2005eventsourcing}.

Los almacenes de características gobiernan definiciones para aprendizaje automático;
los sistemas de linaje explican procedencia; y los entornos tipo \emph{lakehouse}
unifican operación analítica. COBRIA conecta esos mecanismos con la frontera de casos de
uso, sin pretender sustituirlos \parencite{orr2021featurestores,li2023featurestore,
buneman2001provenance,armbrust2021lakehouse}. COBRIA trata las semánticas particulares
de persistencia, aislamiento e índices como capacidades explícitas de cada adaptador
NoSQL, no como supuestos universales.

\begin{table}[H]
\centering
\caption{Comparación entre COBRIA y enfoques relacionados.}
\label{tab:comparacion-literatura-ampliada}
\small
\begin{tabularx}{\textwidth}{p{2.7cm}p{2.2cm}p{2.1cm}X}
\toprule
Enfoque & Autoridad & Reconstrucción & Gobierno y diferencia principal \\
\midrule
CQRS & Parcial & Parcial & Separa lectura y escritura; ámbito y linaje no son inherentes.\\
Vista materializada & Externa & Sí & Optimiza consultas; el ámbito depende del motor y del diseño.\\
Abastecimiento por eventos & Eventos & Sí & El registro de eventos es autoridad y aporta historia.\\
Almacén de características & Fuentes analíticas & Parcial & Gobierna variables de modelos y su procedencia.\\
Arquitectura tipo \emph{lakehouse} & Archivos o tablas & Parcial & Unifica almacenamiento y procesamiento analítico.\\
COBRIA & Declarada & Obligatoria & Exige ámbito, linaje y gobierno desde operación hasta IA.\\
\bottomrule
\end{tabularx}
\end{table}

La matriz no establece superioridad. Muestra que la contribución propuesta reside en
hacer simultáneamente verificables autoridad, ámbito, reconstrucción, costo y continuidad
hacia inteligencia artificial. Si una solución existente satisface esos contratos con
menor complejidad, COBRIA recomienda conservar la solución más sencilla.

\section{Síntesis de la Parte I}

Esta primera parte estableció el problema que COBRIA pretende investigar. Los sistemas
NoSQL multiinquilino necesitan optimizar consultas sin perder autoridad, consistencia ni
capacidad de recuperación. Las necesidades analíticas añaden el requisito de conservar
linaje, versiones y corrección temporal desde la operación hasta la característica y la
predicción. COBRIA se propone como una integración de ideas existentes alrededor de dos
conceptos organizadores: \emph{datos canónicos} y \emph{proyecciones reconstruibles}.

La investigación no presupone que dicha integración sea siempre beneficiosa. Las
preguntas e hipótesis obligan a medir reducción de lectura junto con amplificación de
escritura y almacenamiento, exactitud de reconstrucción junto con su costo, y
reproducibilidad analítica junto con complejidad de adopción. Esta posición crítica será
la base del marco teórico, el modelo formal y la evaluación posterior.


\part{Diseño y formalización de COBRIA}
\chapter{Metodología para el diseño de COBRIA}
\label{ch:metodologia-diseno}

\section{Enfoque de ciencia del diseño}

COBRIA se estudia como un artefacto de ingeniería y no únicamente como una descripción
de prácticas observadas. La ciencia del diseño resulta apropiada cuando la investigación
busca construir y evaluar un artefacto destinado a resolver un problema relevante. Ese
artefacto puede adoptar la forma de constructos, modelos, métodos o instanciaciones
\parencite{hevner2004designscience}. COBRIA comprende las cuatro formas: propone un
vocabulario, formaliza relaciones, prescribe un método de diseño y se instancia en
sistemas ejecutables.

El proceso sigue seis actividades adaptadas de la metodología de
\textcite{peffers2007dsrm}:

\begin{enumerate}[leftmargin=*]
  \item identificar el problema de autoridad, duplicación, consulta y trazabilidad;
  \item definir objetivos verificables para una solución;
  \item diseñar y desarrollar los componentes de COBRIA;
  \item demostrar su aplicación en dos casos anonimizados;
  \item evaluar beneficios, costos y fallos mediante experimentos;
  \item comunicar resultados, límites y material de replicación.
\end{enumerate}

El desarrollo no es estrictamente lineal. Los hallazgos de evaluación pueden exigir
reformular un contrato, retirar una proyección o reducir una abstracción. Este ciclo
evita presentar como principio general una decisión que solo funciona en un caso.

\section{Proceso de abstracción desde los casos}

Los dos sistemas de estudio contienen implementaciones con diferentes grados de
uniformidad. El Sistema A combina componentes heredados y especializados; el Sistema B
utiliza metadatos y generación con mayor regularidad. COBRIA no copia literalmente
ninguna de las dos estructuras. Su derivación sigue cuatro pasos:

\begin{enumerate}[leftmargin=*]
  \item \textbf{Inventario:} identificar entidades, accesos, servicios, operaciones,
  índices, catálogos y procesos asíncronos.
  \item \textbf{Normalización conceptual:} traducir nombres internos a términos
  reconocibles, como Entidad, Repositorio, Mapeador de datos, Capa de servicios y Proyección.
  \item \textbf{Comparación:} separar mecanismos comunes de decisiones exclusivas de
  un dominio o de deuda histórica.
  \item \textbf{Generalización cautelosa:} conservar únicamente propiedades que puedan
  definirse, implementarse y posteriormente evaluarse.
\end{enumerate}

La correspondencia terminológica inicial se muestra en la
\cref{tab:correspondencia-componentes}. El término interno \emph{Dater}, útil para
describir los casos, se sustituye en la propuesta por Repositorio/Mapeador de datos, una
terminología más próxima a la literatura de arquitectura empresarial
\parencite{fowler2002poeaa}.

\begin{table}[H]
\centering
\caption{Normalización conceptual de componentes observados.}
\label{tab:correspondencia-componentes}
\begin{tabularx}{\textwidth}{p{3.2cm}p{4.2cm}X}
\toprule
Nombre observado & Término COBRIA & Responsabilidad principal \\
\midrule
Objeto & Canónico Entidad & Representar identidad, estado semántico, versión y ciclo de vida.\\
Dater & Repositorio/Mapeador de datos & Traducir representaciones y encapsular persistencia y consulta.\\
Servicio & Dominio/Aplicación Servicio & Aplicar políticas, componer datos y coordinar capacidades.\\
CasoDeUso & Caso de uso de aplicación & Ejecutar una intención autorizada y completa.\\
Servicio indexador & Gestor de proyecciones & Mantener, verificar y reconstruir proyecciones.\\
thingsData & Registro de catálogos & Administrar vocabularios finitos y datos de referencia.\\
DataConnect & Pasarela de almacenamiento & Adaptar el motor físico y sus primitivas.\\
\bottomrule
\end{tabularx}
\end{table}

\section{Criterios de diseño}

Cada componente propuesto debe satisfacer cinco criterios:

\begin{description}[leftmargin=3.4cm,style=nextline]
  \item[Comprensibilidad] Su propósito debe poder explicarse sin conocer una ruta física.
  \item[Verificabilidad] Deben existir invariantes o resultados observables.
  \item[Reemplazabilidad] El dominio no debe depender innecesariamente de un proveedor.
  \item[Proporcionalidad] La complejidad añadida debe responder a una necesidad medible.
  \item[Trazabilidad] Debe conocerse el origen y la versión de toda derivación relevante.
\end{description}

Estos criterios impiden equiparar uniformidad con calidad. Una clase generada puede ser
estructuralmente correcta y, aun así, carecer de comportamiento funcional. Del mismo
modo, un componente especializado puede ser adecuado aunque no herede de una base
genérica.

\section{Evidencia prevista}

La evaluación combinará cuatro clases de evidencia:

\begin{itemize}[leftmargin=*]
  \item \textbf{estructural:} dependencias, contratos, metadatos y cobertura;
  \item \textbf{funcional:} resultados de operaciones y pruebas de invariantes;
  \item \textbf{experimental:} latencia, bytes, escrituras, almacenamiento y tiempo de
  reconstrucción;
  \item \textbf{analítica:} linaje, repetibilidad, corrección temporal y calidad de los
  datos derivados.
\end{itemize}

La presencia de archivos o documentación se considerará evidencia estructural, no prueba
de funcionamiento. Una capacidad solo se declarará validada si existe una ejecución
observable bajo condiciones documentadas.

\chapter{Arquitectura COBRIA}
\label{ch:arquitectura-cobria}

\section{Definición}

\begin{quote}
COBRIA es una arquitectura dirigida por metadatos que representa el estado autorizado
mediante entidades canónicas, separa el modelo semántico de su representación física y
produce proyecciones reconstruibles para operación, integración, minería de datos e
inteligencia artificial dentro de un ámbito explícito y auditable.
\end{quote}

La definición contiene cinco compromisos. Primero, existe una fuente de autoridad
identificable. Segundo, almacenar no equivale a modelar: la representación física puede
cambiar sin alterar el significado. Tercero, los datos derivados poseen contrato y no
son copias accidentales. Cuarto, la pertenencia a un organización o contexto forma parte del
diseño. Quinto, analítica e IA consumen proyecciones trazables en lugar de exportaciones
opacas.

\section{Capas y dirección de dependencias}

\begin{figure}[H]
\centering
\resizebox{0.92\textwidth}{!}{%
\begin{tikzpicture}[
  layer/.style={draw=black,minimum width=12.8cm,minimum height=1cm,align=center},
  side/.style={draw=black,dashed,minimum width=3.2cm,minimum height=1cm,align=center},
  arr/.style={-{Latex[length=2.4mm]},semithick},
  node distance=5mm and 8mm
]
\node[layer] (present) {Presentación, API y consumidores};
\node[layer,below=of present] (use) {Casos de uso: intención, autorización e idempotencia};
\node[layer,below=of use] (service) {Servicios de dominio y aplicación};
\node[layer,below=of service] (repo) {Repositorios / Mapeadores de datos};
\node[layer,below=of repo,fill=black!8] (canonical) {Entidades canónicas};
\node[layer,below=of canonical] (storage) {Pasarela de almacenamiento y motor NoSQL};
\node[side,right=of use] (meta) {COBRIA Metadatos};
\node[side,right=of canonical] (proj) {Gestores de proyecciones};
\node[side,left=of canonical] (gov) {Auditoría y gobierno};
\draw[arr] (present)--(use);
\draw[arr] (use)--(service);
\draw[arr] (service)--(repo);
\draw[arr] (repo)--(canonical);
\draw[arr] (canonical)--(storage);
\draw[dashed] (meta)--(use);
\draw[dashed] (meta)--(proj);
\draw[dashed] (proj)--(canonical);
\draw[dashed] (gov)--(canonical);
\end{tikzpicture}%
}
\caption{Capas y preocupaciones transversales de COBRIA.}
\label{fig:capas-cobria}
\end{figure}

La dirección descendente de la \cref{fig:capas-cobria} representa conocimiento de
abstracciones, no necesariamente el orden temporal de cada llamada. Un escucha del
almacenamiento puede iniciar una notificación ascendente, pero su dato debe cruzar el
Mapeador de datos antes de presentarse como entidad. La interfaz no debe construir rutas
físicas ni actualizar proyecciones por cuenta propia.

\section{Módulos internos}

\subsection{COBRIA Núcleo}

Incluye entidades, repositorios, servicios, casos de uso, catálogos y contexto. Es el
subconjunto mínimo. Un sistema pequeño puede adoptar Núcleo sin Analítica o Inteligencia.

\subsection{COBRIA Metadatos}

Contiene descripciones de campos, tipos, mapeos, ámbito, relaciones, sensibilidad,
índices y ciclo de vida. Los metadatos pueden alimentar validadores y generadores, pero
no sustituyen pruebas funcionales.

\subsection{COBRIA Proyecciones}

Define funciones de derivación, destinos, niveles de consistencia, verificadores y
procedimientos de reconstrucción. La proyección puede residir en el mismo motor o en uno
diferente.

\subsection{COBRIA Analítica}

Transforma datos autorizados en eventos, agregaciones, series, grafos, conjuntos de datos y
características versionadas. Su diseño debe preservar linaje y tiempo de observación.

\subsection{COBRIA Inteligencia}

Administra versiones de modelos, inferencias, evidencia, retroalimentación, consentimiento y
agentes. Ninguna predicción adquiere por sí sola autoridad sobre el estado canónico.

\subsection{COBRIA Gobierno}

Reúne auditoría, retención, protección de datos, observabilidad, reconciliación y
políticas de acceso. La gobernanza no se considera un módulo opcional cuando se manejan
datos sensibles o decisiones asistidas por IA.

\section{COBRIA frente a patrones relacionados}

COBRIA puede usar separación de lectura y escritura, pero no exige dos bases ni
mensajería. Por ello no es sinónimo de CQRS. CQRS permite optimizar modelos de lectura y
escritura de forma independiente, aunque añade sincronización y consistencia eventual
\parencite{microsoft2025cqrs}. COBRIA adopta esa separación cuando existe una razón
medible y añade el contrato de reconstrucción y el linaje analítico.

Tampoco exige abastecimiento por eventos. La fuente canónica puede ser el estado actual versionado,
un historial de eventos o una combinación. En un sistema basado en estado, la auditoría
no convierte automáticamente cada cambio en un evento capaz de reconstruir todo el
agregado.

Las proyecciones COBRIA incluyen vistas materializadas, pero su alcance es mayor: pueden
ser índices booleanos, relaciones inversas, resúmenes, artefactos multimedia o conjuntos de datos.
Comparten la propiedad esencial de ser desechables y regenerables desde una fuente
\parencite{microsoft2025materialized,gupta1995materialized}.

\chapter{Modelo formal de datos canónicos}
\label{ch:modelo-formal}

\section{Conjuntos fundamentales}

Sea $\mathcal{E}$ el conjunto de entidades canónicas, $\mathcal{P}$ el conjunto de
proyecciones, $\mathcal{S}$ el conjunto de ámbitos y $\mathcal{M}$ el conjunto de
versiones de metadatos. Una entidad se define como:

\begin{equation}
E=(I,A,S,V,L,R),
\label{eq:entidad}
\end{equation}

donde $I$ es la identidad, $A$ los atributos semánticos, $S\in\mathcal{S}$ el ámbito,
$V$ la versión del esquema, $L$ el estado de ciclo de vida y $R$ la revisión de
concurrencia.

La identidad global lógica puede expresarse como:

\begin{equation}
I_g=(\tau,S,I_l),
\end{equation}

donde $\tau$ es el tipo de entidad e $I_l$ el identificador local. Dos entidades con el
mismo identificador local pero ámbitos distintos no son la misma entidad.

\section{Atributos y ausencia de valor}

Cada atributo $a_j\in A$ declara tipo, nulabilidad, valor inicial, sensibilidad y reglas
de validación. COBRIA distingue entre ausencia, valor nulo y valor vacío cuando el
dominio lo requiere. Convertir indiscriminadamente los tres estados puede destruir
información necesaria para migraciones o analítica.

Los valores monetarios deben conservar moneda y unidad menor; las fechas de negocio no
deben confundirse con instantes; y los identificadores externos deben diferenciarse de
las claves internas. Estas decisiones pertenecen al modelo semántico, no al nombre corto
utilizado en almacenamiento.

\section{Ámbito y multiinquilinato}

Un ámbito puede contener varios niveles:

\begin{equation}
S=(\mathrm{organizacion},\mathrm{contexto}_1,\ldots,\mathrm{contexto}_q).
\end{equation}

Para toda operación autorizada $o$ realizada por un actor $u$ sobre una entidad $E$, se
exige:

\begin{equation}
\operatorname{autorizado}(u,o,E)\Rightarrow \operatorname{ambito}(E)\in
\operatorname{ambitos}(u,o).
\label{eq:ambito}
\end{equation}

La condición de la \cref{eq:ambito} debe verificarse en una frontera confiable. Ocultar
una entidad en la interfaz no constituye aislamiento. La ruta física puede incluir el
ámbito para acotar consultas, pero la seguridad no debe depender únicamente de que el
cliente construya correctamente esa ruta.

\section{Ciclo de vida}

El estado $L$ se modela como:

\begin{equation}
L\in\{active,archived,deleted\},
\end{equation}

con extensiones específicas del dominio. \emph{Eliminado} puede representar una marca
lógica cuando existen obligaciones de auditoría. Una política de retención define si y
cuándo procede la eliminación física. Las proyecciones deben excluir o representar
entidades archivadas de acuerdo con su contrato, no según una convención implícita.

\section{Revisión y concurrencia}

La revisión $R\in\mathbb{N}_0$ aumenta con cada modificación aceptada. Una actualización
optimista exige:

\begin{equation}
R_{request}=R_{stored}.
\end{equation}

Si la igualdad no se cumple, el caso de uso rechaza, reintenta desde un estado nuevo o
ejecuta una reconciliación explícita. El control de revisión no reemplaza transacciones
atómicas cuando una invariante involucra varias ubicaciones.

\section{Invariantes canónicas}

Se proponen cinco invariantes generales:

\begin{align}
&\forall E\in\mathcal{E}: I(E)\neq\varnothing,\\
&\forall E\in\mathcal{E}: S(E)\neq\varnothing\quad\text{si el tipo es ámbitod},\\
&\forall E_1,E_2: I_g(E_1)=I_g(E_2)\Rightarrow E_1=E_2,\\
&\forall E: deserialize(serialize(E,M),M)=E,\\
&\forall p\in\mathcal{P}: source(p)\subseteq\mathcal{E}\cup\mathcal{P}.
\end{align}

La cuarta igualdad se entiende sobre campos persistidos y admite normalizaciones
declaradas. Un timestamp generado por servidor, por ejemplo, debe modelarse como tal en
vez de ocultarse como una diferencia accidental.

\chapter{Representación física y persistencia}
\label{ch:persistencia}

\section{Mapeo semántico--físico}

Sea $M_v$ un mapeo versionado y $F$ una función de serialización:

\begin{equation}
F_{M_v}:\mathcal{E}\rightarrow\mathcal{R},
\end{equation}

donde $\mathcal{R}$ es el conjunto de representaciones físicas. La operación inversa
$G_{M_v}$ debe satisfacer:

\begin{equation}
G_{M_v}(F_{M_v}(E))\equiv E.
\label{eq:roundtrip}
\end{equation}

La equivalencia de la \cref{eq:roundtrip} incluye identidad y semántica, no propiedades
transitorias de memoria. Un ejemplo neutral es:

\begin{table}[H]
\centering
\caption{Ejemplo de mapeo semántico--físico versionado.}
\label{tab:fieldmap}
\begin{tabular}{lll}
\toprule
Campo semántico & Clave física & Regla \\
\midrule
organizaciónId & a & cadena no vacía \\
contextId & b & cadena no vacía \\
status & c & código de catálogo versionado \\
createdAt & d & instante UTC \\
revision & e & entero no negativo \\
\bottomrule
\end{tabular}
\end{table}

Las claves compactas pueden reducir repetición, pero disminuyen legibilidad y elevan el
costo de diagnóstico. COBRIA no las prescribe; exige medir su beneficio y conservar un
diccionario versionado. Una clave física nunca debe reutilizarse con otro significado
dentro de la misma versión.

\section{Repositorio y Mapeador de datos}

El Repositorio ofrece una colección conceptual de entidades; el Mapeador de datos traduce
entre objeto y almacenamiento sin introducir dependencias físicas en el dominio
\parencite{fowler2002poeaa}. En una implementación compacta ambas responsabilidades
pueden residir en una clase, siempre que su interfaz permanezca explícita.

El contrato mínimo es:

\begin{verbatim}
get(id, ámbito) -> Entidad | null
list(query, ámbito, page) -> Page<Entidad>
save(entity, expectedRevision) -> Entidad
remove(id, ámbito, policy) -> Result
subscribe(id, ámbito, observer) -> Unsubscribe
\end{verbatim}

El Repositorio no decide por sí solo una transición de negocio. Puede validar forma y
ámbito, pero una operación como aprobar, cancelar o transferir pertenece al caso de uso.

\section{Mapa de identidad y caché}

Una caché $C:I_g\rightarrow E$ evita cargar repetidamente la misma identidad durante una
unidad de trabajo. Su acceso promedio puede aproximarse a $O(1)$ mediante tabla hash,
pero introduce expiración, invalidación y riesgo de servir datos obsoletos. COBRIA exige
declarar el alcance de la caché:

\begin{itemize}[leftmargin=*]
  \item por operación, con vida corta y consistencia sencilla;
  \item por sesión, con invalidación mediante eventos o TTL;
  \item compartida, con políticas explícitas de coherencia y aislamiento.
\end{itemize}

Una caché nunca debe omitir la dimensión de organización en su clave.

\section{Pasarela de almacenamiento}

El Pasarela adapta primitivas como lectura, transacción, actualización multipath,
escucha, paginación y generación de identidad. Esta frontera permite contrastar el
modelo en emuladores y sustituir el motor, pero no promete portabilidad automática: las
garantías varían entre almacenes. Sistemas como Dynamo priorizan disponibilidad y
resolución de versiones \parencite{decandia2007dynamo}, mientras Spanner demuestra un
modelo distribuido con transacciones externamente consistentes bajo otras decisiones de
infraestructura \parencite{corbett2012spanner}.

\section{Registro de rutas semánticas}

Cuando el almacenamiento utiliza nombres físicos compactos, el desacoplamiento no se
logra solo con un Mapeador. También se necesita un registro central que traduzca nombres
semánticos a nodos y construya rutas con ámbito. Sea
\(\rho_v:S\times K^*\rightarrow Path\) la función versionada que recibe un símbolo
semántico \(S\) y segmentos de clave \(K\). Los servicios de aplicación dependen de
\(\rho_v\), no de literales físicos dispersos.

Este registro aporta cuatro propiedades comprobables:

\begin{enumerate}[leftmargin=*]
  \item una ruta física cambia en un punto controlado;
  \item una búsqueda estática puede prohibir literales fuera de infraestructura;
  \item las pruebas verifican normalización y presencia obligatoria de ámbito;
  \item telemetría, reglas y migraciones comparten el mismo vocabulario.
\end{enumerate}

El registro no convierte automáticamente un esquema compacto en modelo canónico. Su
función es establecer una frontera semántico--física y hacer localizable la deuda.

\section{Capa de servicios y casos de uso}

El Capa de servicios establece capacidades y coordina la respuesta de la aplicación
\parencite{fowler2002poeaa}. COBRIA distingue:

\begin{itemize}[leftmargin=*]
  \item \textbf{servicio de dominio:} regla que no pertenece naturalmente a una sola
  entidad;
  \item \textbf{servicio de aplicación:} composición técnica de repositorios, políticas
  y proveedores;
  \item \textbf{caso de uso:} intención completa iniciada por un actor o proceso.
\end{itemize}

Un caso de uso recibe comando, actor, ámbito, clave de idempotencia y revisión esperada.
Produce un resultado funcional, una auditoría y los cambios derivados requeridos. La
interfaz no debe convertir un cambio de estado complejo en un CRUD genérico.

\section{Contexto autorizado como dato derivado confiable}

El ámbito solicitado por cabeceras, query string o cuerpo identifica una intención, no
demuestra autorización. COBRIA distingue \(S_{solicitado}\) de
\(S_{autorizado}\). El segundo debe derivarse en servidor a partir de identidad,
declaraciones y relación comprobada con el recurso:

\[
S_{autorizado}=\operatorname{autorizar}(actor,recurso,politica).
\]

Servicios, auditoría, facturación, FinOps y proyecciones posteriores consumen únicamente
el contexto autorizado. Recalcularlo de manera distinta en cada capa crea divergencia;
aceptarlo del cliente permite atribución falsa. Por ello, el contexto validado es una
proyección efímera de seguridad: tiene fuente, política, alcance y vida limitada, pero
no reemplaza la autorización de cada operación sensible.

\chapter{Teoría de proyecciones reconstruibles}
\label{ch:proyecciones}

\section{Definición}

Una proyección $P_i$ se define mediante una función versionada:

\begin{equation}
P_i=g_{i,v}(D_i),
\label{eq:projection}
\end{equation}

donde $D_i$ es un conjunto de entidades canónicas o proyecciones previamente
autorizadas. Para evitar ciclos no controlados, el grafo de dependencias de
reconstrucción debe ser acíclico o declarar un punto fijo y un algoritmo de convergencia.

Una copia es proyección COBRIA solo si declara:

\begin{enumerate}[leftmargin=*]
  \item fuentes y ámbito;
  \item función y versión;
  \item destino y clave;
  \item propietario de escritura;
  \item nivel de consistencia;
  \item procedimiento de reconstrucción;
  \item verificador de integridad;
  \item política de privacidad y retención.
\end{enumerate}

\section{Taxonomía}

\begin{table}[H]
\centering
\caption{Taxonomía de proyecciones COBRIA.}
\label{tab:taxonomia-proyecciones}
\begin{tabularx}{\textwidth}{p{2.5cm}p{3.3cm}X}
\toprule
Clase & Consistencia esperada & Ejemplos y tratamiento \\
\midrule
Crítica & Atómica con la operación o verificada antes de confirmar & Reserva de recurso,
saldo operativo, unicidad. Si no puede mantenerse, falla el caso de uso.\\
Operacional & Eventual acotada & Índices de consulta, resúmenes de lista, relaciones
inversas. Debe existir reconciliación.\\
Analítica & Eventual y versionada & Series, características, conjuntos de datos, agregados históricos.
Prioriza linaje y corrección temporal.\\
Efímera & Sin garantía durable & Caché, vista previa o cálculo temporal. Puede descartarse sin
reconstrucción persistente.\\
\bottomrule
\end{tabularx}
\end{table}

Clasificar no convierte una proyección en correcta. La clase define qué retrasos y
fallos son aceptables y qué pruebas se requieren.

\section{Propiedades}

\subsection{Determinismo}

Con fuentes y versión iguales, una proyección determinista produce el mismo resultado:

\begin{equation}
D_i^{(1)}=D_i^{(2)}\Rightarrow g_{i,v}(D_i^{(1)})=g_{i,v}(D_i^{(2)}).
\end{equation}

Si utiliza tiempo actual, aleatoriedad o un proveedor externo, esos valores deben
capturarse como entradas versionadas.

\subsection{Idempotencia}

Aplicar dos veces el mismo cambio no debe duplicar resultados:

\begin{equation}
apply(apply(P,c),c)=apply(P,c).
\end{equation}

\subsection{Reconstructibilidad}

Si $P_i$ se pierde, debe poder calcularse $P_i'$ tal que:

\begin{equation}
verify(P_i',g_{i,v}(D_i))=true.
\end{equation}

\subsection{No autoridad}

Una proyección no se edita para corregir negocio. La corrección se ejecuta sobre el
canónico mediante un caso de uso y luego se propaga.

\section{Grafo de dependencias}

Sea $G=(V,E_d)$ un grafo donde $V=\mathcal{E}\cup\mathcal{P}$ y una arista $(x,p)$
indica que $p$ depende de $x$. El orden de reconstrucción se obtiene mediante orden
topológico. Si se detecta un ciclo no declarado, el plan se rechaza.

\begin{figure}[H]
\centering
\begin{tikzpicture}[
  n/.style={draw=black,rounded corners=1pt,minimum width=3.2cm,minimum height=8mm,align=center},
  p/.style={draw=black,fill=black!7,minimum width=3.2cm,minimum height=8mm,align=center},
  a/.style={-{Latex[length=2.2mm]},semithick},node distance=9mm and 14mm]
\node[n] (e1) {Entidad canónica A};
\node[n,below=of e1] (e2) {Entidad canónica B};
\node[p,right=of e1] (idx) {Índice operacional};
\node[p,right=of e2] (sum) {Resumen compuesto};
\node[p,right=of sum] (data) {Conjunto de datos versionado};
\draw[a] (e1)--(idx);
\draw[a] (e1)--(sum);
\draw[a] (e2)--(sum);
\draw[a] (sum)--(data);
\end{tikzpicture}
\caption{Ejemplo de grafo acíclico de reconstrucción.}
\label{fig:grafo-proyecciones}
\end{figure}

\section{Contrato de proyección}

Una representación neutral puede adoptar la siguiente forma:

\begin{verbatim}
ProjectionContract {
  name, version, class,
  sources[], ámbito,
  destination, keyFunction,
  build, verify,
  consistency, responsable,
  privacy, retention
}
\end{verbatim}

El contrato es declarativo respecto al propósito, aunque `build` y `verify` pueden ser
funciones. Debe poder responder quién escribe la proyección y cómo se repara. Si esas
preguntas no tienen respuesta, la duplicación permanece fuera de gobierno.

\chapter{Algoritmos, consistencia y recuperación}
\label{ch:algoritmos-consistencia}

\section{Creación de una entidad}

\begin{verbatim}
algorithm Create(command, actor, idempotencyKey):
  require authorize(actor, command.ámbito, command.action)
  if resultExists(idempotencyKey): return previousResult
  metadata <- registry.forType(command.type)
  entity <- metadata.construct(command.carga útil)
  validate(entity)
  critical <- projections.criticalFor(entity.type)
  atomicWrite(entity, buildAll(critical, entity))
  enqueue(projections.nonCriticalFor(entity.type), entity.revision)
  audit(actor, command, entity)
  return remember(idempotencyKey, entity)
\end{verbatim}

Las proyecciones críticas forman parte de la confirmación. Las demás se envían a un
proceso idempotente. Si el motor no ofrece una escritura atómica suficiente, el diseño
debe reducir la frontera crítica o implementar una reserva y compensación explícitas.

\section{Actualización optimista}

\begin{verbatim}
algorithm Update(command, actor, expectedRevision):
  current <- repository.get(command.id, command.ámbito)
  require current.revision = expectedRevision
  next <- useCase.transition(current, command)
  next.revision <- current.revision + 1
  validateTransition(current, next)
  atomicWrite(next, criticalDelta(current, next))
  enqueue(nonCriticalDelta(current, next))
  return next
\end{verbatim}

El delta debe eliminar claves antiguas además de añadir nuevas. Cambiar el estado o el
ámbito sin retirar índices previos produce referencias fantasma.

\section{Reconstrucción total}

\begin{verbatim}
algorithm RebuildAll(contract, ámbito):
  acquireLease(contract.name, ámbito)
  targetVersion <- createShadowDestination()
  for page in sourcePages(contract.sources, ámbito):
    records <- contract.build(page)
    writeIdempotently(targetVersion, records)
    checkpoint(page.cursor)
  report <- contract.verify(targetVersion)
  require report.errors = 0
  atomicCutover(targetVersion)
  releaseLease()
  return report
\end{verbatim}

El destino sombra evita dejar una proyección parcialmente reconstruida como si estuviera
completa. Para volúmenes pequeños puede reconstruirse en sitio, pero debe exponerse el
estado `rebuilding` y tolerar reanudación.

\section{Reconstrucción incremental}

\begin{verbatim}
algorithm RebuildChanged(contract, marca de corte):
  changes <- canonicalChangesAfter(marca de corte)
  for change in ordered(changes):
    require notProcessed(change.id, contract.version)
    delta <- contract.buildDelta(change.before, change.after)
    applyIdempotently(delta)
    markProcessed(change.id, contract.version)
  return contract.verifyAffected(changes)
\end{verbatim}

La reconstrucción incremental necesita un registro confiable de cambios o una marca de
actualización consultable. Si no existe, escanear el conjunto canónico puede ser la única
forma correcta de detectar diferencias.

\section{Verificación de integridad}

Se definen cuatro clases de anomalía:

\begin{align}
orphan &= \{p\in P: sourceId(p)\notin I(\mathcal{E})\},\\
missing &= \{e\in E: expected(e)\setminus actual(e)\neq\varnothing\},\\
stale &= \{p\in P: sourceRevision(p)<revision(source(p))\},\\
divergent &= \{p\in P:p\neq g_{v}(source(p))\}.
\end{align}

Un reporte debe separar conteos, ejemplos anonimizados y severidad. La reparación
automática es apropiada para derivados deterministas; una divergencia canónica requiere
revisión de negocio.

\section{Consistencia por clase}

\begin{table}[H]
\centering
\caption{Mecanismos de consistencia recomendados.}
\label{tab:consistencia-mecanismos}
\begin{tabularx}{\textwidth}{p{3cm}p{4cm}X}
\toprule
Necesidad & Mecanismo & Consideración \\
\midrule
Invariante local & Transacción o compare-and-set & Rechazar revisiones obsoletas.\\
Varias rutas del mismo motor & Actualización multipath atómica & Validar límites reales
del proveedor.\\
Derivado operacional & Evento/cola de salida y proceso de trabajo idempotente & Medir retraso y reconciliar.\\
Integración externa & Saga o compensación & No asumir transacción distribuida.\\
Conjunto de datos analítico & Lote versionado con marca de corte & Preservar tiempo y linaje.\\
\bottomrule
\end{tabularx}
\end{table}

La consistencia fuerte global no debe afirmarse por analogía. Las garantías dependen del
motor, partición y frontera transaccional. COBRIA documenta esas garantías en el contrato
del caso de uso.

\section{Estado materializado: completitud, frescura y procedencia}

Una vista no debe considerarse válida por el solo hecho de existir. El contrato mínimo
incluye metadatos de control:

\begin{verbatim}
_meta:
  complete: boolean
  updatedAt: timestamp
  schemaVersion: integer
  sourceRevision|marca de corte: value
  count|checksum: optional
\end{verbatim}

La lectura acepta la vista únicamente si está completa, dentro de su presupuesto de
frescura y en una versión compatible. Una mutación puede actualizarla, invalidarla o
marcarla incompleta. Si la proyección es mejor esfuerzo, su fallo no revierte la escritura
canónica; la lectura cae a la fuente autorizada y programa reparación.

Cuando varias solicitudes encuentran una vista fría, la reconstrucción debe
\emph{coalescerse}: una promesa, arrendamiento o bloqueo por ámbito evita que todas repitan el mismo
escaneo. Esta protección contra \emph{cache stampede} forma parte del contrato de
reconstrucción y debe probarse con solicitudes concurrentes.

\section{Arquitectura ejecutable y funciones de aptitud}

Una regla escrita puede degradarse silenciosamente. COBRIA incorpora funciones de
aptitud arquitectónica: pruebas o análisis estáticos que convierten una propiedad
estructural en una condición de integración continua. Ejemplos:

\begin{itemize}[leftmargin=*]
  \item presentación no importa infraestructura de datos;
  \item rutas físicas compactas no aparecen en servicios de aplicación;
  \item todo modelo de lectura declara metadatos de completitud y frescura;
  \item una operación sensible usa contexto autorizado por servidor;
  \item el número de consumidores de una fachada heredado no puede crecer;
  \item rutas críticas respetan presupuestos de bytes, LCP y bloqueo.
\end{itemize}

Para una deuda medible \(d_t\), el control automatizado de migración impone
\(d_{t+1}\leq d_t\), y el objetivo final puede fijarse en cero. Este presupuesto
monótono permite migrar sin una reescritura total y evita que una fachada temporal se
convierta en arquitectura permanente.

\section{Cola de salida, reintentos y duplicados}

Cuando guardar el canónico y publicar un mensaje no forman una transacción única, existe
una ventana de pérdida. El patrón cola de salida conserva la intención de publicación junto al
cambio autorizado; un proceso de trabajo la entrega y registra el resultado. La entrega puede ser
al menos una vez, por lo que consumidores y Gestores de proyecciones deben ser idempotentes.

Una clave de idempotencia se asocia con actor, ámbito, operación y ventana temporal. No
debe reutilizarse entre organizaciones. El resultado previo puede devolverse siempre que los
parámetros relevantes coincidan; en caso contrario, la reutilización se rechaza.

\section{Recuperación y observabilidad}

Cada proceso de reconstrucción registra:

\begin{itemize}[leftmargin=*]
  \item contrato y versión;
  \item ámbito y rango temporal;
  \item cursor o checkpoint;
  \item entidades leídas y registros escritos;
  \item anomalías antes y después;
  \item duración y consumo aproximado;
  \item identidad del operador o proceso de trabajo;
  \item resultado del cambio de versión.
\end{itemize}

El objetivo de recuperación (RTO) y el punto de recuperación (RPO) deben definirse por
clase de proyección. Una caché puede perderse sin RPO; un saldo derivado utilizado en una
operación exige un tratamiento diferente.

\section{Antipatrones}

COBRIA identifica como antipatrones:

\begin{enumerate}[leftmargin=*]
  \item editar una proyección para corregir la fuente;
  \item construir rutas físicas en la presentación;
  \item mantener el mismo índice desde múltiples servicios sin propietario;
  \item usar un ID sin ámbito como clave de caché multiinquilino;
  \item generar una clase por cada estructura sin justificar identidad o ciclo de vida;
  \item llamar transacción a una secuencia de escrituras independientes;
  \item entrenar con una exportación sin versión, marca de corte ni linaje;
  \item permitir que un agente de IA escriba directamente en el almacenamiento;
  \item declarar una integración completa por la sola presencia de modelos o contratos;
  \item añadir proyecciones sin medir amplificación y frecuencia de consulta;
  \item confiar en el ámbito enviado por el cliente para auditoría o autorización;
  \item servir una vista existente sin comprobar completitud, versión y frescura;
  \item crear una capa de compatibilidad sin inventario ni presupuesto decreciente;
  \item documentar límites de capas sin un control automatizado que detecte regresiones.
\end{enumerate}

\section{Criterio mínimo de adopción}

COBRIA Núcleo puede justificarse cuando existe al menos una de estas condiciones: dominio
con reglas relevantes, multiinquilinato, duplicación derivada, necesidad de reconstrucción,
procesamiento asíncrono o uso analítico trazable. Si el sistema es un CRUD pequeño, con
pocos registros y consultas directas, Repositorio más un modelo sencillo puede ser
suficiente. La arquitectura se considera exitosa cuando hace explícitos los compromisos,
incluso si la decisión resultante es no crear una proyección.

\section{Síntesis de la Parte II}

Esta parte transformó la intuición inicial en un artefacto definido. COBRIA organiza el
estado mediante entidades con identidad, ámbito, versión, ciclo de vida y revisión;
separa semántica de persistencia mediante un mapeo reversible; y asigna a Repositorio,
Servicio y Caso de uso fronteras distintas. Su elemento diferenciador es el contrato de
proyección, que convierte índices, resúmenes y conjuntos de datos en derivaciones gobernadas y
reconstruibles.

También se establecieron algoritmos de creación, actualización, reconstrucción y
verificación, junto con mecanismos de consistencia proporcionales a la criticidad. Estas
definiciones preparan la siguiente parte del estudio: analizar complejidad, amplificación
y costo, y comprobar mediante experimentos si la disciplina propuesta produce un
beneficio neto.


\part{Rendimiento, complejidad y optimización}
\chapter{Modelo de rendimiento y notación}

La utilidad de una arquitectura de datos no puede deducirse únicamente de la claridad
de sus clases o diagramas. Una separación correcta de responsabilidades puede reducir
errores y, al mismo tiempo, introducir lecturas, escrituras, almacenamiento y tareas de
mantenimiento adicionales. Esta parte establece un modelo para razonar sobre esos
compromisos antes de medirlos y un protocolo para comprobarlos sin atribuir a COBRIA
resultados que dependan del motor, la red o una configuración particular.

El análisis distingue dos preguntas. La primera es asintótica: cómo crece el trabajo
cuando aumenta el volumen de datos. La segunda es empírica: cuánto tarda y cuánto cuesta
una operación bajo una carga y una plataforma concretas. La notación $O(\cdot)$ responde
a la primera pregunta; no representa milisegundos, precio monetario ni garantía de
rendimiento. Dos operaciones de orden $O(1)$ pueden tener costos reales muy distintos
si una requiere una lectura local y otra varias comunicaciones de red
\parencite{cormen2022algorithms}.

\section{Variables del modelo}

Sea $N$ el número de entidades canónicas dentro de un ámbito; $R$ el número de resultados
de una consulta; $P$ el número de proyecciones asociadas con una entidad; $P_w\leq P$ el
número de proyecciones afectadas por una escritura; $\Delta N$ el número de entidades
modificadas desde un marca de corte; $B_c$ el tamaño medio de un canónico y $B_{p_i}$ el
tamaño medio de su contribución a la proyección $i$. Se utilizan además:

\begin{table}[H]
\centering
\caption{Símbolos del modelo de rendimiento.}
\label{tab:simbolos-rendimiento}
\begin{tabularx}{\textwidth}{>{\bfseries}p{2cm}X>{\bfseries}p{2cm}X}
\toprule
Símbolo & Significado & Símbolo & Significado \\
\midrule
$L$ & latencia observada & $Q$ & operaciones por unidad de tiempo \\
$S$ & bytes almacenados & $T$ & bytes transferidos \\
$W_f$ & escrituras físicas & $R_f$ & lecturas físicas \\
$C_x$ & costo unitario del recurso $x$ & $h$ & proporción de aciertos de caché \\
$f_q$ & frecuencia de la consulta & $f_w$ & frecuencia de escritura \\
$d$ & número de dependencias & $k$ & tamaño del lote \\
\bottomrule
\end{tabularx}
\end{table}

Estas variables se calculan por organización y también de manera agregada. El promedio global
puede ocultar que un organización pequeño recibe peor servicio o que uno grande domina el
consumo. Por ello, cada resultado debe conservar al menos identificador anonimizado de
ámbito, tipo de operación, versión del contrato y configuración experimental.

\section{Frontera de medición}

Una medición debe indicar dónde comienza y termina. COBRIA propone tres fronteras:

\begin{enumerate}[leftmargin=*]
  \item \textbf{Algorítmica:} trabajo efectuado por la transformación, sin red ni
  persistencia.
  \item \textbf{Componente:} Repositorio, Gestor de proyecciones o Constructor de conjuntos de datos junto
  con sus dependencias simuladas o reales.
  \item \textbf{Extremo a extremo:} desde el caso de uso hasta la respuesta observable,
  incluyendo autorización, serialización, red, almacenamiento y reintentos.
\end{enumerate}

No deben compararse cifras de fronteras distintas como si midieran el mismo fenómeno.
La instrumentación se incorpora fuera del dominio siempre que sea posible para evitar
que el modelo canónico dependa de un proveedor de telemetría.

\section{Métricas principales}

La latencia se reportará mediante mediana y percentiles p95 y p99, no solo con el
promedio. El rendimiento efectivo se expresa como operaciones completadas por segundo bajo una
concurrencia declarada. También se registrarán bytes leídos y escritos, cantidad de
operaciones físicas, errores, reintentos, uso de memoria, tiempo de CPU y retraso de
actualización de proyecciones. El enfoque sigue la práctica de separar carga de trabajo,
métricas y configuración al comparar sistemas de almacenamiento
\parencite{cooper2010ycsb,jain1991performance}.

\chapter{Complejidad computacional de las operaciones}

\section{Acceso canónico}

Cuando una clave física se deriva directamente de $(S,id)$ y el motor ofrece acceso
indexado por clave, la búsqueda lógica de una entidad se modela como $O(1)$ promedio.
Esta expresión presupone una clave conocida; no cubre resolución DNS, autenticación,
red, deserialización ni contención. Si se busca por un atributo no indexado, el costo
puede crecer hasta $O(N)$ porque deben examinarse los registros del ámbito.

La recuperación de una página de $R$ resultados desde una proyección ordenada puede
modelarse como $O(\log N+R)$ cuando existe una estructura de búsqueda balanceada, o como
$O(R)$ después de localizar un cursor. El uso de offset puede degradarse al aumentar la
profundidad; COBRIA prefiere cursores estables compuestos por orden, ámbito e identidad.

\section{Creación y actualización}

Validar un objeto con $m$ campos tiene costo $O(m)$ si cada regla es local. Las reglas
que consultan otras entidades agregan $d$ accesos y deben expresarse como
$O(m+d)$ más el costo de esos accesos. Persistir solamente el canónico requiere una
escritura lógica. Mantener síncronamente $P_w$ proyecciones agrega, en el caso simple,
$O(P_w)$ transformaciones y escrituras:

\begin{equation}
T_{write}=T_{validate}+T_{canonical}+\sum_{i=1}^{P_w}
  (T_{g_i}+T_{persist_i}).
\label{eq:tiempo-escritura}
\end{equation}

La suma no implica que todos los términos sean iguales. Una proyección de clave puede
ser constante, mientras un agregado que inspecciona una colección puede depender de su
tamaño. La ejecución paralela reduce tiempo de pared, pero no elimina trabajo total ni
garantiza atomicidad.

\section{Reconstrucción total e incremental}

Una reconstrucción total que aplica $P$ transformaciones independientes a $N$ entidades
tiene costo superior $O(NP)$. Si cada entidad alimenta solamente un subconjunto fijo,
el costo efectivo es $O(N\bar{P})$, donde $\bar{P}$ es el promedio de proyecciones
aplicables. La memoria puede mantenerse en $O(k)$ mediante procesamiento por lotes, en
vez de cargar $O(N)$ objetos.

Una reconstrucción incremental basada en un índice confiable de cambios procesa
$\Delta N$ entidades:

\begin{equation}
T_{incremental}\approx O(\Delta N\bar{P})+T_{discover}+T_{verify}.
\end{equation}

Es ventajosa cuando $\Delta N\ll N$ y el costo de descubrir cambios no domina. Si el
marca de corte es incompleto, una reconstrucción incremental rápida puede producir un estado
incorrecto. La optimización, por tanto, está subordinada a la cobertura del registro de
cambios.

\section{Verificación de proyecciones}

Comparar cada canónico con sus derivados cuesta $O(N\bar{P})$. Un muestreo de $s$
entidades reduce el trabajo a $O(s\bar{P})$, pero solo estima la tasa de anomalías. Los
resúmenes criptográficos por partición permiten localizar diferencias sin transferir todos los objetos;
su construcción inicial sigue siendo lineal y su actualización requiere disciplina de
versionado.

\begin{table}[H]
\centering
\caption{Complejidad esperada de operaciones COBRIA.}
\label{tab:complejidad-operaciones}
\begin{tabularx}{\textwidth}{p{4cm}p{3.2cm}X}
\toprule
Operación & Orden esperado & Condición principal \\
\midrule
Buscar canónico por clave & $O(1)$ promedio & Ruta derivable de ámbito e ID.\\
Escanear por atributo & $O(N)$ & No existe proyección o índice.\\
Consultar proyección & $O(\log N+R)$ & Índice ordenado; depende del motor.\\
Validar entidad & $O(m+d)$ & Reglas locales y dependencias explícitas.\\
Actualizar derivados & $O(P_w)$ & Transformaciones acotadas por proyección.\\
Reconstrucción total & $O(N\bar{P})$ & Recorrido completo por lotes.\\
Reconstrucción incremental & $O(\Delta N\bar{P})$ & Registro de cambios completo.\\
Verificación completa & $O(N\bar{P})$ & Comparación fuente--derivado.\\
Generar conjunto de datos & $O(NF)$ & $F$ características por entidad, sin uniones externos.\\
\bottomrule
\end{tabularx}
\end{table}

La tabla representa modelos de crecimiento y no contratos del proveedor. Cada
implementación deberá confirmar sus supuestos con planes de consulta, contadores y
pruebas de carga.

\chapter{Amplificación y modelo de costos}

\section{Amplificación de lectura}

La amplificación de lectura relaciona trabajo físico y resultado útil. Para una consulta
$q$ se define:

\begin{equation}
A_r(q)=\frac{B_{read}(q)}{\max(B_{result}(q),\epsilon)},
\end{equation}

donde $\epsilon$ evita división por cero en consultas sin resultados. También puede
calcularse por registros: $A_r^{reg}=R_f/\max(R,1)$. Una consulta que escanea diez mil
registros para devolver diez tiene amplificación aproximada de mil. Una proyección bien
diseñada acerca el numerador al volumen útil, aunque añade mantenimiento y almacenamiento.

\section{Amplificación de escritura}

Para una operación lógica $o$:

\begin{equation}
A_w(o)=\frac{W_f(o)}{W_l(o)},
\end{equation}

donde normalmente $W_l=1$. Si se persiste el canónico, una entrada cola de salida y tres
proyecciones, $A_w=5$ sin contar réplicas internas del motor. Deben reportarse por
separado la amplificación visible para la aplicación y la amplificación interna cuando
el proveedor la exponga.

\section{Amplificación de almacenamiento}

El factor de almacenamiento es:

\begin{equation}
A_s=\frac{S_{canonical}+S_{projections}+S_{metadata}+S_{registros}}
{S_{canonical}}.
\end{equation}

Una proyección no se justifica por ser pequeña de forma aislada. Su costo acumulado
incluye versiones durante reconstrucciones, índices del motor, historial, respaldos y
transferencia. COBRIA asigna a cada proyección un presupuesto y un propietario.

\section{Costo total por ventana}

Para una ventana temporal $\tau$, el costo observable se modela como:

\begin{align}
C_{total}(\tau)=&\ C_{read}R_f + C_{write}W_f + C_{store}S\tau \\
&+ C_{transfer}T + C_{compute}U + C_{ops}H,
\label{eq:costo-total}
\end{align}

donde $U$ representa consumo computacional y $H$ esfuerzo operativo estimado. El último
término evita declarar óptima una solución barata para el proveedor pero costosa de
operar, depurar o reconstruir. Los precios cambian; por ello el experimento publicará
cantidades físicas y dejará los precios como parámetros reemplazables.

\section{Punto de equilibrio de una proyección}

Sea $C_s$ el costo de una consulta sin proyección, $C_p$ el costo de consultarla con
proyección y $C_m$ el costo de mantenerla por escritura. La proyección produce beneficio
en una ventana cuando:

\begin{equation}
f_q(C_s-C_p)>f_wC_m+C_{storage}+C_{rebuild}.
\label{eq:break-even}
\end{equation}

Esta desigualdad no pretende reducir toda decisión a dinero. Latencia máxima,
disponibilidad, aislamiento y corrección pueden actuar como restricciones obligatorias.
Sí permite identificar proyecciones con uso insuficiente o mantenimiento excesivo.

\section{Presupuesto de proyecciones}

Cada contrato debe declarar límites de bytes, escrituras por cambio, retraso tolerable,
tiempo máximo de reconstrucción y frecuencia mínima que justifica su existencia. Un
reporte periódico compara presupuesto y observación. Si una proyección excede el límite,
las opciones son optimizarla, reducir campos, agrupar eventos, cambiar consistencia o
retirarla; no se oculta el exceso mediante promedios.

\chapter{Estrategias de optimización}

\section{Optimizar desde la consulta}

La secuencia recomendada parte de una consulta observable: frecuencia, selectividad,
orden, tamaño de respuesta y objetivo de latencia. Después se decide si basta una clave
canónica, un índice del motor, una proyección específica o una caché. Crear estructuras
derivadas antes de conocer la carga aumenta $P$ y puede empeorar H9.

Una proyección debe almacenar solo campos necesarios para filtrar, ordenar y responder.
Duplicar el documento completo facilita algunas lecturas, pero aumenta transferencia,
amplificación y riesgo de exposición. La paginación por cursor evita recorrer desplazamientos
crecientes y conserva una frontera consistente de lectura cuando el orden incluye una
clave de desempate.

\section{Lotes y paralelismo controlado}

El tamaño de lote $k$ equilibra memoria, viajes de red y duración de transacciones. Un
lote demasiado pequeño multiplica comunicaciones; uno demasiado grande presiona memoria
y extiende la recuperación después de un fallo. COBRIA no fija un $k$ universal: lo
selecciona experimentalmente y registra la configuración.

El paralelismo se limita por partición, cuota y capacidad de escritura. Aumentar procesos de trabajo
puede reducir tiempo hasta alcanzar contención o limitación de capacidad; después, el rendimiento efectivo se
estanca y la latencia de cola aumenta. El controlador debe aplicar control de contrapresión, variación aleatoria
y reintentos acotados, conservando idempotencia.

\section{Caché con ámbito y revisión}

Con una tasa de aciertos $h$, el costo medio simplificado de lectura es:

\begin{equation}
E[T]=hT_{hit}+(1-h)T_{miss}.
\end{equation}

La clave incluye organización, tipo, ID y, cuando corresponda, versión. La caché no sustituye
una proyección si se necesita consulta por atributos, ni se considera autoridad. TTL
reduce permanencia, pero no garantiza invalidación inmediata. Para datos sensibles se
deben medir además exposición, borrado y segregación.

\section{Reconstrucción eficiente}

Una reconstrucción segura y optimizada sigue cinco decisiones:

\begin{enumerate}[leftmargin=*]
  \item fijar versión y marca de corte de entrada;
  \item dividir por ámbito y rango estable;
  \item procesar lotes idempotentes con checkpoints;
  \item escribir en una versión paralela de la proyección;
  \item verificar y efectuar un cambio de versión atómico o reversible.
\end{enumerate}

La versión paralela requiere almacenamiento temporal, pero evita exponer una mezcla de
contratos. Para reconstrucciones incrementales se conserva una operación total como
mecanismo de recuperación y como prueba periódica de reconstruibilidad.

\section{Optimización adaptativa}

La arquitectura puede recomendar cambios a partir de telemetría, sin permitir que una
heurística modifique por sí sola el estado canónico. Un ciclo adaptativo observa
frecuencia, amplificación, latencia y retraso; propone crear, rediseñar o retirar una
proyección; estima el punto de equilibrio; ejecuta una comparación controlada; y somete
el cambio a aprobación. Esto hace posible usar minería de datos para descubrir patrones
de acceso y aprendizaje automático para predecir carga, manteniendo una frontera de
gobierno.

\begin{figure}[H]
\centering
\begin{tikzpicture}[
  node distance=1.15cm and 1.25cm,
  box/.style={draw,rounded corners,align=center,minimum width=2.55cm,
    minimum height=.9cm,fill=cobriaLight},
  line/.style={-{Latex[length=2mm]},thick}
]
\node[box] (observe) {Observar\\telemetría};
\node[box,right=of observe] (analyze) {Analizar\\costo y carga};
\node[box,right=of analyze] (propose) {Proponer\\cambio};
\node[box,below=of propose] (evaluate) {Evaluar\\en paralelo};
\node[box,left=of evaluate] (approve) {Aprobar y\\versionar};
\node[box,left=of approve] (verify) {Verificar\\resultado};
\draw[line] (observe) -- (analyze);
\draw[line] (analyze) -- (propose);
\draw[line] (propose) -- (evaluate);
\draw[line] (evaluate) -- (approve);
\draw[line] (approve) -- (verify);
\draw[line] (verify) -- (observe);
\end{tikzpicture}
\caption{Ciclo gobernado de optimización adaptativa.}
\label{fig:ciclo-optimizacion}
\end{figure}

\chapter{Protocolo experimental de rendimiento}

\section{Objetivo y comparaciones}

El protocolo evaluará si COBRIA reduce el costo de lectura y mejora reconstrucción y
trazabilidad a cambio de costos cuantificables. No se presentarán resultados hasta
ejecutar el protocolo. Se compararán, como mínimo, estas configuraciones:

\begin{enumerate}[leftmargin=*]
  \item \textbf{B0, acceso directo:} canónicos sin proyección específica;
  \item \textbf{B1, optimización nativa:} índices disponibles en el motor;
  \item \textbf{C1, COBRIA síncrona:} canónico y proyecciones críticas en la operación;
  \item \textbf{C2, COBRIA asíncrona:} canónico, cola de salida y proyecciones actualizadas por
  procesos de trabajo;
  \item \textbf{C3, COBRIA incremental:} C2 con registro de cambios y reconstrucción
  incremental.
\end{enumerate}

No todas las configuraciones aplican a toda operación. La comparación se declara antes
de observar resultados y mantiene la misma semántica, conjunto de datos y política de consistencia.

\section{Factores y niveles}

\begin{table}[H]
\centering
\caption{Factores propuestos para los experimentos.}
\label{tab:factores-experimento}
\begin{tabularx}{\textwidth}{p{3.6cm}X X}
\toprule
Factor & Niveles iniciales & Propósito \\
\midrule
Volumen $N$ & $10^3$, $10^4$, $10^5$ y, si es viable, $10^6$ & Estimar tendencia de
crecimiento.\\
Mezcla lectura/escritura & 95/5, 80/20, 50/50 & Representar cargas distintas.\\
Concurrencia & 1, 8, 32 y 128 clientes & Detectar saturación y contención.\\
Número de proyecciones & 0, 1, 3, 5 y 10 & Medir amplificación marginal.\\
Distribución & uniforme y sesgada & Evaluar claves calientes.\\
Tamaño de entidad & pequeño, medio y grande & Medir serialización y transferencia.\\
Cambios $\Delta N/N$ & 0.1\%, 1\%, 10\% y 100\% & Comparar reconstrucción incremental
y total.\\
\bottomrule
\end{tabularx}
\end{table}

Los valores son puntos iniciales y se ajustarán a la capacidad del entorno antes de
registrar la especificación definitiva. El Sistema A aportará cargas temporales,
multimodales y de eventos; el Sistema B, operaciones empresariales, catálogos e
aislamiento multiinquilino. Los nombres y datos originales no aparecerán en scripts ni
resultados.

\section{Datos y cargas reproducibles}

Los generadores usarán semilla publicada, distribuciones declaradas y relaciones
válidas. Para evitar que datos triviales favorezcan una estrategia, incluirán cardinalidad,
sesgo, valores ausentes, revisiones y tamaños variables. Los artefactos multimedia se
representarán mediante metadatos y archivos sintéticos; no se copiará contenido privado.

Cada ejecución incluye calentamiento, ventana de medición y enfriamiento. Se aleatoriza
el orden de configuraciones cuando sea posible y se repiten ejecuciones independientes.
Se registra hardware, sistema operativo, entorno de ejecución, versión del motor, región, límites,
configuración de índices y revisión del código. Un prueba comparativa como YCSB sirve de referencia
para estructurar cargas, pero las operaciones de dominio y reconstrucción requieren
escenarios COBRIA adicionales \parencite{cooper2010ycsb}.

\section{Instrumentación y control}

Cada operación recibe un identificador de correlación. Los contadores se toman en el
cliente y, cuando exista acceso, en el almacenamiento. Se sincronizan relojes para medir
retraso asíncrono, pero la latencia local utiliza un reloj monotónico. Se evita registrar
cargas útiles sensibles. La telemetría debe distinguir éxito, rechazo de negocio, conflicto,
timeout y reintento.

Los ensayos controlarán cachés, calentamiento, concurrencia y tareas de fondo. Si un
servicio administrado introduce ruido no controlable, se reportará y se aumentarán
repeticiones, sin eliminar observaciones solo por ser desfavorables.

\section{Análisis estadístico}

Para cada combinación se reportarán tamaño muestral, mediana, p95, p99, intervalo de
confianza y distribución gráfica. Las comparaciones incluirán diferencia absoluta,
razón y tamaño de efecto. Antes de usar pruebas paramétricas se examinarán independencia,
forma de distribución y varianza. Cuando no se satisfagan supuestos, se preferirán
intervalos remuestreo con reemplazo o pruebas no paramétricas. Las múltiples comparaciones se
identificarán y corregirán cuando corresponda.

La significancia estadística no reemplaza relevancia práctica. Una mejora pequeña puede
carecer de valor si aumenta mucho $A_w$, $A_s$ o el retraso. El análisis de sensibilidad
variará precios, frecuencias y objetivos de latencia en la ecuación
\ref{eq:costo-total}.

\section{Criterios de aceptación}

Las hipótesis se juzgarán con reglas previas:

\begin{itemize}[leftmargin=*]
  \item H1 se apoya si el crecimiento de lecturas se aproxima a $R$ y la latencia de
  cola mejora sin cambiar la semántica;
  \item H2 se cuantifica, no se considera un fallo: toda reducción se acompaña de
  $A_w$ y $A_s$;
  \item H3 requiere igualdad según el contrato, ausencia de huérfanos y cobertura
  completa después de reconstruir;
  \item H4 compara total e incremental para distintos $\Delta N/N$ y busca su punto de
  cruce;
  \item H9 se apoya si existen regiones donde el costo marginal supera el beneficio de
  lectura.
\end{itemize}

No se modificará un umbral después de conocer resultados sin declarar el cambio como
exploratorio.

\section{Plantilla de resultados}

\begin{table}[H]
\centering
\caption{Matriz mínima de resultados por escenario.}
\label{tab:plantilla-resultados}
\begin{tabularx}{\textwidth}{p{3.7cm}X X}
\toprule
Dimensión & Indicadores & Interpretación requerida \\
\midrule
Respuesta & p50, p95, p99, rendimiento efectivo & Saturación y variabilidad.\\
Lectura & registros, bytes, $A_r$ & Trabajo evitado frente a resultado.\\
Escritura & operaciones, bytes, $A_w$ & Precio de mantener derivados.\\
Almacenamiento & bytes y $A_s$ & Incluye versiones e índices.\\
Reconstrucción & duración, cobertura, anomalías & Exactitud antes que velocidad.\\
Consistencia & retraso, conflictos, reintentos & Ventana de obsolescencia.\\
Operación & CPU, memoria, horas de trabajo & Costo total y mantenibilidad.\\
Equidad de ámbito & percentiles por organización & Evitar ocultar degradación.\\
\bottomrule
\end{tabularx}
\end{table}

\chapter{Rendimiento para minería de datos e inteligencia artificial}

\section{Preparación analítica como proyección}

Un conjunto de datos se trata como una proyección versionada del estado canónico y de eventos
permitidos. Su costo de construcción depende de entidades, ventanas temporales y número
de características. Para $N$ entidades y $F$ características locales, el caso simple es
$O(NF)$. Uniones, agregados temporales, texto o embeddings añaden términos que deben
medirse por separado.

Esta estructura favorece minería de datos porque conserva el significado de la entidad,
el ámbito, el tiempo del evento, la versión de transformación y el linaje. Permite volver
a producir un conjunto de datos, comparar versiones y detectar qué cambios de datos explican una
diferencia. No garantiza que las características sean útiles ni que el modelo sea justo.

\section{Corrección temporal y prevención de fuga}

Para una observación con tiempo de corte $t$, solo se utilizan hechos disponibles en o
antes de $t$. Una característica calculada con datos futuros produce fuga y métricas de
entrenamiento engañosas. El Constructor de conjuntos de datos recibe marca de corte, ventana y versión; el
resultado registra tiempo de evento, tiempo de ingestión y tiempo de cálculo.

El acceso punto-en-tiempo puede requerir historial canónico, log de eventos o snapshots.
Si el sistema conserva únicamente el estado actual, COBRIA debe declarar que ciertas
reconstrucciones históricas no son posibles en vez de inferirlas.

\section{Características, embeddings y multimodalidad}

Las características tabulares, series, texto y referencias multimedia pueden compartir
un contrato de procedencia. Los archivos grandes permanecen fuera del objeto canónico;
este conserva ubicación autorizada, hash, tipo, versión y política. Una extracción de
audio, texto o imagen genera un derivado con modelo, parámetros y fecha.

Los embeddings son proyecciones reconstruibles si se conserva la entrada autorizada,
el identificador exacto del modelo y la transformación. Cambiar de modelo produce una
nueva versión; no se mezclan vectores incompatibles en el mismo índice sin una estrategia
explícita. La reconstrucción puede ser costosa, por lo que se presupuestan cómputo,
transferencia y ventana de coexistencia.

\section{RAG y agentes}

En recuperación aumentada, el índice vectorial no es autoridad: sirve para localizar
candidatos. La respuesta debe enlazar fragmento, documento canónico, versión y permisos.
La autorización se aplica antes de incorporar contexto al instrucción. Un agente puede leer
proyecciones adaptadas a su tarea, pero cualquier cambio de negocio vuelve a un caso de
uso que valida identidad, ámbito, invariantes e idempotencia.

Esta frontera limita el daño de una inferencia incorrecta y permite auditar propuesta,
evidencia, decisión humana y efecto final. El rendimiento del modelo se mide separado
del rendimiento de datos: recuperación, construcción de contexto, inferencia y
persistencia tienen latencias y costos distintos.

\section{Optimización de cadenas de procesamiento de IA}

COBRIA permite reutilizar transformaciones deterministas, materializar características
de alto uso y reconstruirlas en lotes. La decisión se rige por el mismo punto de
equilibrio de la ecuación \ref{eq:break-even}. Para entrenamiento se favorecen lotes
grandes y reproducibilidad; para inferencia en línea, baja latencia y frescura. Una sola
representación física rara vez optimiza ambos caminos, por lo que los contratos deben
declarar equivalencia semántica entre versión fuera de línea y en línea, idea central de los
Almacenes de características \parencite{orr2021featurestores,li2023featurestore}.

\begin{figure}[H]
\centering
\begin{tikzpicture}[
  node distance=.95cm and 1.05cm,
  box/.style={draw,align=center,minimum width=2.45cm,minimum height=.9cm},
  derived/.style={box,fill=cobriaLight},
  line/.style={-{Latex[length=2mm]},thick}
]
\node[box] (canonical) {Datos canónicos\\y eventos};
\node[derived,right=of canonical] (builder) {Constructor de conjuntos de datos\\versionado};
\node[derived,right=of builder] (datosNodo) {Conjunto de datos /\\Conjunto de características};
\node[derived,below=of datosNodo] (model) {Modelo y\\evaluación};
\node[derived,left=of model] (prediction) {Predicción con\\evidencia};
\node[box,left=of prediction] (usecase) {Caso de uso\\autorizado};
\draw[line] (canonical) -- (builder);
\draw[line] (builder) -- (datosNodo);
\draw[line] (datosNodo) -- (model);
\draw[line] (model) -- (prediction);
\draw[line] (prediction) -- (usecase);
\draw[line,dashed] (usecase) -- node[left,font=\small]{cambio validado} (canonical);
\end{tikzpicture}
\caption{Flujo gobernado desde datos canónicos hasta decisiones asistidas por IA.}
\label{fig:flujo-ia}
\end{figure}

\section{Métricas analíticas y de IA}

Además de exactitud, precisión, exhaustividad u otras métricas propias del problema, se
registrarán cobertura de linaje, reproducibilidad entre ejecuciones, frescura,
divergencia en línea--fuera de línea, proporción de campos desconocidos, costo por ejemplo y
latencia por etapa. Para RAG se separan recuperación y generación; para agentes se miden
también acciones rechazadas, confirmaciones y efectos reversibles. La gestión de riesgo
acompaña el ciclo completo y no se reduce a una puntuación del modelo
\parencite{nist2023airmf}.

\chapter{Amenazas a la validez y síntesis de la Parte III}

\section{Validez interna}

Cambios de caché, calentamiento, throttling, tareas de fondo o distinta configuración
pueden explicar una diferencia atribuida a COBRIA. Se mitigarán mediante automatización,
orden aleatorio, repetición, configuración registrada y comparación sobre la misma
semántica. La instrumentación también tiene costo; se medirá una configuración de
control cuando su efecto pueda ser relevante.

\section{Validez de constructo}

Latencia no equivale a calidad arquitectónica, cantidad de clases no equivale a
mantenibilidad y reconstrucción exitosa no equivale a corrección del negocio. Por ello
cada concepto utiliza varios indicadores. La reconstruibilidad exige fuente suficiente,
contrato versionado, resultado verificable y ausencia de autoridad en el derivado.

\section{Validez externa}

Dos sistemas anonimizados permiten variación de dominio, pero no representan todos los
motores ni escalas. Los emuladores pueden diferir de servicios administrados. Los
resultados se generalizarán como condiciones y mecanismos ---por ejemplo, relación entre
frecuencia de lectura y amplificación---, no como una promesa universal de milisegundos.

\section{Validez de conclusión}

Muestras pequeñas, colas largas y múltiples comparaciones pueden producir conclusiones
inestables. Se publicarán distribuciones, intervalos, tamaños de efecto y resultados
negativos. Una hipótesis no apoyada no se ocultará; puede delimitar el dominio donde la
arquitectura deja de ser conveniente.

\section{Síntesis de la Parte III}

Esta parte convirtió la afirmación general de ``optimizar'' en magnitudes observables.
El beneficio esperado de las proyecciones es reducir el trabajo de lectura desde una
dependencia de $N$ hacia una dependencia más cercana a $R$, a cambio de trabajo de
escritura $O(P_w)$, almacenamiento adicional y reconstrucción. Las ecuaciones de
amplificación y costo permiten comparar ese intercambio sin confundir complejidad
asintótica con tiempo real.

También se definieron estrategias de optimización, un protocolo reproducible y reglas
de análisis que impiden inventar resultados. Finalmente, el mismo principio de
proyección reconstruible se extendió a conjuntos de datos, características, embeddings, RAG y
agentes, incorporando corrección temporal, linaje y autorización. Con ello queda
preparada la siguiente parte: desarrollar en profundidad la adaptabilidad de COBRIA a
analítica e inteligencia artificial y formalizar sus mecanismos de gobierno.


\part{Analítica, minería de datos e inteligencia artificial}
\chapter{Extensión analítica de COBRIA}

La analítica y la inteligencia artificial suelen incorporarse después de que un sistema
operacional ha acumulado datos. En ese momento aparecen exportaciones manuales,
significados incompatibles, atributos sin procedencia y modelos difíciles de reproducir.
COBRIA aborda este problema desde la arquitectura: considera conjuntos de datos, características,
embeddings, índices de recuperación y predicciones como proyecciones gobernadas de datos
canónicos y eventos autorizados.

La propuesta no convierte el dominio en un sistema de aprendizaje automático. Mantiene
una frontera explícita entre hechos de negocio, evidencia derivada e inferencias
probabilísticas. Esta frontera permite utilizar la misma disciplina de identidad, ámbito,
versionado, reconstrucción y trazabilidad tanto en minería descriptiva como en modelos
predictivos, recuperación aumentada y agentes.

\section{De la operación al conocimiento}

Se distinguen cinco niveles semánticos:

\begin{enumerate}[leftmargin=*]
  \item \textbf{Hecho canónico:} estado validado por un caso de uso y sujeto a
  invariantes del dominio.
  \item \textbf{Evento:} observación temporal de un cambio o interacción, con tiempo,
  ámbito y procedencia.
  \item \textbf{Derivación analítica:} transformación determinista o controlada, como
  una característica, agregado o fragmento.
  \item \textbf{Inferencia:} salida probabilística de un modelo, acompañada de versión,
  evidencia y medidas de incertidumbre disponibles.
  \item \textbf{Decisión:} acción aceptada por una persona o caso de uso autorizado.
\end{enumerate}

Una inferencia no asciende automáticamente a hecho. Si una clasificación, recomendación
o respuesta debe modificar el negocio, se formula como propuesta y atraviesa validación,
autorización e idempotencia. Esta regla evita que una salida plausible pero incorrecta
sobrescriba la fuente canónica.

\section{Capacidades analíticas}

\begin{table}[H]
\centering
\caption{Adaptación de COBRIA a modalidades analíticas.}
\label{tab:capacidades-analiticas}
\begin{tabularx}{\textwidth}{p{3.2cm}p{4cm}X}
\toprule
Modalidad & Proyección principal & Requisitos COBRIA \\
\midrule
Descriptiva & Agregado o cubo & Definición, ventana, ámbito y reconciliación.\\
Minería de datos & Conjunto de datos versionado & Semilla, filtros, linaje y calidad.\\
Series temporales & Secuencia por entidad & Tiempo de evento, ingestión y marca de corte.\\
Aprendizaje supervisado & Conjunto de características y etiquetas & Corte temporal y prevención de fuga.\\
Texto y multimedia & Fragmentos y descriptores & Resumen criptográfico, derechos, modelo extractor y versión.\\
Búsqueda semántica & Índice de embeddings & Modelo compatible, permisos y reconstrucción.\\
RAG & Corpus recuperable & Evidencia, vigencia y autorización previa.\\
Agentes & Herramientas tipadas & Capacidades mínimas, aprobación y auditoría.\\
\bottomrule
\end{tabularx}
\end{table}

\section{Componentes de la extensión}

La extensión de inteligencia se compone de un Constructor de conjuntos de datos, un Registro de características, un
Almacén de artefactos, un Registro de modelos, un Pasarela de recuperación, un Motor de políticas y un Evidencia
Registro. Estos componentes se apoyan en Núcleo, Metadatos, Proyecciones y Gobierno; no
acceden a rutas físicas mediante convenciones privadas.

\begin{figure}[H]
\centering
\begin{tikzpicture}[
  node distance=.7cm and .8cm,
  box/.style={draw,align=center,minimum width=2.4cm,minimum height=.82cm},
  derived/.style={box,fill=cobriaLight},
  line/.style={-{Latex[length=2mm]},thick}
]
\node[box] (core) {Núcleo y\\eventos};
\node[derived,right=of core] (datosNodo) {Constructor de\\conjuntos de datos};
\node[derived,right=of datosNodo] (características) {Registro de\\características};
\node[derived,below=of características] (models) {Registro de\\modelos};
\node[derived,left=of models] (retrieval) {Pasarela de\\recuperación};
\node[box,left=of retrieval] (usecase) {Caso de uso\\autorizado};
\node[box,below=of datosNodo] (policy) {Motor de políticas y Registro de evidencias};
\draw[line] (core) -- (datosNodo);
\draw[line] (datosNodo) -- (características);
\draw[line] (características) -- (models);
\draw[line] (models) -- (retrieval);
\draw[line] (retrieval) -- (usecase);
\draw[line,dashed] (policy) -- (datosNodo);
\draw[line,dashed] (policy) -- (models);
\draw[line,dashed] (policy) -- (retrieval);
\end{tikzpicture}
\caption{Componentes de la extensión analítica e inteligente de COBRIA.}
\label{fig:extension-inteligencia}
\end{figure}

\section{Principios de adaptación}

La extensión conserva ocho principios: autoridad canónica, derivación declarada,
reconstruibilidad proporcional, aislamiento por ámbito, corrección temporal, mínimo dato
necesario, evidencia enlazable y supervisión proporcional al riesgo. La
reconstruibilidad es proporcional porque no toda salida probabilística puede reproducirse
bit a bit; en esos casos se exige reproducibilidad del proceso, captura de parámetros y
medición de variación.

\chapter{Contratos de conjuntos de datos y linaje}

\section{Definición formal}

Un contrato de conjunto de datos se representa como:

\begin{equation}
\mathcal{D}=\langle id,v,Q,F,W,S,L,P,V,C\rangle,
\end{equation}

donde $id$ es la identidad; $v$, la versión; $Q$, la selección de fuentes; $F$, las
transformaciones; $W$, la política temporal o marca de corte; $S$, el esquema de salida;
$L$, el conjunto de etiquetas; $P$, las políticas; $V$, las validaciones; y $C$, las
condiciones de publicación. Una ejecución produce un manifiesto $M$ con entradas,
salidas, resúmenes criptográficos, tiempos, código y entorno.

El conjunto de datos se considera reconstruible si existen fuentes permitidas suficientes, las
transformaciones y dependencias están fijadas, el contrato puede ejecutarse y el
resultado puede verificarse según una equivalencia declarada. Para datos tabulares
deterministas puede exigirse igualdad exacta; para conversiones multimedia o librerías
no deterministas se requiere tolerancia documentada.

\section{Contrato mínimo}

\begin{longtable}{>{\bfseries}p{3.2cm}p{5.1cm}p{6cm}}
\caption{Contenido mínimo de un contrato de conjunto de datos.}\label{tab:contrato-conjunto de datos}\\
\toprule
Elemento & Pregunta que responde & Evidencia conservada \\
\midrule
\endfirsthead
\toprule
Elemento & Pregunta que responde & Evidencia conservada \\
\midrule
\endhead
Propósito & ¿Para qué análisis puede utilizarse? & Tarea, población y usos excluidos.\\
Fuentes & ¿De dónde procede cada campo? & Tipo canónico, evento, versión y ámbito.\\
Selección & ¿Qué se incluye o excluye? & Predicados, ventanas y razones.\\
Transformación & ¿Cómo se obtiene cada variable? & Función, parámetros y código.\\
Tiempo & ¿Qué era conocido en el corte? & tiempo del evento, tiempo de ingestión y marca de corte.\\
Esquema & ¿Qué significa cada columna? & Tipo, unidad, dominio y nulabilidad.\\
Calidad & ¿Qué condiciones debe cumplir? & Reglas, umbrales y resultados.\\
Privacidad & ¿Qué datos están permitidos? & Base de acceso, minimización y retención.\\
Publicación & ¿Cuándo es utilizable? & Aprobación, hash, ubicación y estado.\\
Retiro & ¿Cuándo deja de utilizarse? & Fecha, motivo y dependencias notificadas.\\
\bottomrule
\end{longtable}

Este contrato complementa prácticas de documentación como las fichas de datos, que
buscan explicitar motivación, composición, proceso de recolección, usos y mantenimiento
\parencite{gebru2021datasheets}. COBRIA añade correspondencia ejecutable con entidades,
proyecciones y ámbitos del sistema.

\section{Grafo de linaje}

El linaje se modela como un grafo dirigido acíclico $G_L=(V_L,E_L)$. Los nodos pueden
ser fuentes, ejecuciones, conjuntos de datos, características, modelos, evaluaciones y
predicciones. Una arista expresa \emph{derivado de}, \emph{entrenado con},
\emph{evaluado con} o \emph{producido por}. Cada nodo posee identidad y versión; una etiqueta textual
sin referencia estable no constituye linaje suficiente.

Para una predicción $y$, la consulta de procedencia debe poder recorrer:

\begin{equation}
y\rightarrow \text{modelo}\rightarrow \text{Conjunto de características}\rightarrow
\text{conjunto de datos}\rightarrow \text{fuente canónica o evento}.
\end{equation}

El grafo no almacena necesariamente los datos completos. Puede conservar resúmenes criptográficos,
manifiestos y referencias protegidas. Los permisos se vuelven a comprobar al navegar;
poseer el ID de un nodo no concede acceso a su contenido.

\section{Algoritmo de materialización}

\begin{verbatim}
BuildDataset(contractId, version, ámbito, cutoff):
  contract <- registry.require(contractId, version)
  policy.require("conjunto de datos.build", ámbito, contract.purpose)
  snapshot <- resolveSources(contract.sources, ámbito, cutoff)
  run <- ledger.start(contract, snapshot, actor, environment)
  for batch in snapshot.readBatches():
      rows <- contract.transform(batch)
      contract.validate(rows)
      artifact.write(run.stagingVersion, rows)
      run.checkpoint(batch.cursor, rows.hash)
  summary <- verify(run.stagingVersion, contract)
  if not summary.accepted: quarantine(run, summary)
  else publishAtomically(run, summary)
  return run.manifiesto
\end{verbatim}

La publicación se realiza después de validar, no mientras el conjunto de datos está incompleto.
Una ejecución fallida queda en cuarentena con sus métricas y no reemplaza la versión
vigente.

\section{Calidad y observabilidad}

Las reglas se agrupan en esquema, completitud, dominio, unicidad, integridad referencial,
distribución, tiempo y privacidad. Un cambio estadístico no siempre es un defecto; puede
reflejar un cambio real. Por eso las alertas de deriva generan revisión y no corrección
automática del canónico.

\chapter{Almacén de características, series temporales y minería de datos}

\section{Características como proyecciones}

Una característica se define mediante:

\begin{equation}
\phi=\langle nombre,v,entidad,ambito,tipo,unidad,g,W,frescura,responsable\rangle.
\end{equation}

La función $g$ transforma fuentes autorizadas; $W$ define la ventana temporal. La
identidad de entidad permite unir valores sin depender de campos ambiguos. El registro
describe significado y versión; los valores pueden almacenarse fuera de línea, en línea o en
ambos medios.

Los Almacenes de características buscan coordinar cálculo, almacenamiento y servicio de
características entre entrenamiento e inferencia \parencite{orr2021featurestores,
li2023featurestore}. COBRIA los incorpora como una especialización de sus proyecciones:
la definición vive en metadatos, la fuente conserva autoridad y las materializaciones
pueden reconstruirse.

\section{Equivalencia fuera de línea--en línea}

Sea $\phi_o(e,t)$ la característica calculada fuera de línea y $\phi_n(e,t)$ la servida en línea.
Se exige una relación de equivalencia:

\begin{equation}
\phi_o(e,t)\simeq_{\epsilon}\phi_n(e,t),
\end{equation}

donde $\epsilon=0$ para valores deterministas exactos y puede ser una tolerancia
justificada para cómputo numérico. Las diferencias se investigan por contrato, versión,
marca de corte, ventana y código. Compartir nombre no demuestra equivalencia.

\section{Corrección punto-en-tiempo}

Para una observación de entrenamiento en $t_c$, COBRIA solo admite una fuente $x$ si su
tiempo de disponibilidad satisface $t_a(x)\leq t_c$. El tiempo del evento puede ser
anterior y el de ingestión posterior; usar únicamente event time puede introducir
información que todavía no estaba disponible.

\begin{equation}
\mathcal{D}(t_c)=\{x\mid \operatorname{ambito}(x)=s\land t_a(x)\leq t_c\land
\operatorname{politica}(x)=\mathrm{permitido}\}.
\end{equation}

Esta condición reduce fuga temporal. Las etiquetas usan una ventana futura separada y
no participan en la construcción de características del mismo corte.

\begin{figure}[H]\centering
\begin{tikzpicture}\begin{axis}[ybar,bar width=22pt,width=.78\textwidth,height=6.5cm,ymin=0,ymax=1.08,ylabel={Valor},symbolic x coords={Precisión,Exhaustividad,F1},xtick=data,ymajorgrids=true,nodes near coords]
\addplot[fill=gray!45,draw=black] coordinates {(Precisión,1.0) (Exhaustividad,1.0) (F1,1.0)};\end{axis}\end{tikzpicture}
\caption{Contratos sintéticos de recuperación en 30 repeticiones. No constituyen evaluación humana de un modelo generativo.}\label{fig:contratos-ia}\end{figure}


\section{Eventos tardíos y marcas de avance temporal}

Los eventos se clasifican como a tiempo, tardíos pero aceptables, o fuera de la ventana
de corrección. La política declara cuánto se espera, cuándo se recalcula una proyección
y cuándo se publica una corrección. Un marca de corte no significa que nunca aparecerá otro
evento; expresa una decisión operacional sobre completitud suficiente.

\section{Minería descriptiva y descubrimiento}

COBRIA facilita segmentación, reglas de asociación, detección de anomalías y análisis de
secuencias porque cada observación conserva entidad, ámbito, tiempo y procedencia. El
proceso exploratorio puede descubrir nuevas variables o patrones, pero sus resultados
se registran como hipótesis analíticas. Una agrupación no redefine automáticamente una
categoría del negocio y una anomalía no implica fraude o error.

Para análisis no supervisado se versionan preprocesamiento, escalamiento, imputación,
selección de variables, métrica de distancia, algoritmo y semilla. Esto permite separar
la inestabilidad del algoritmo de cambios en los datos de entrada.

\section{Catálogo analítico}
\enlargethispage{3\baselineskip}

El catálogo enlaza entidades, eventos, características, conjuntos de datos, modelos y responsables.
Debe responder quién puede usar un artefacto, con qué propósito, de qué versión depende,
cuándo fue actualizado, qué calidad alcanzó y qué consumidores se verán afectados por
un cambio. No es solo una lista de archivos: refleja relaciones ejecutables del grafo de
linaje.

\chapter{Multimodalidad, embeddings y búsqueda semántica}

\section{Objeto multimedia canónico}

El canónico de un recurso multimedia no contiene necesariamente los bytes. Conserva
identidad, ámbito, propietario, tipo MIME validado, tamaño, hash, ubicación lógica,
tiempos, consentimiento o base de uso, retención y estado. Las miniaturas,
transcripciones, descriptores y embeddings son derivados.

Separar contenido y metadatos evita transportar objetos grandes en consultas comunes y
permite políticas distintas. Una URL temporal no es identidad permanente; el Artefacto
Pasarela resuelve la referencia solo después de autorización.

\section{Pipeline multimodal}

\begin{figure}[H]
\centering
\begin{tikzpicture}[
  node distance=.75cm and .8cm,
  box/.style={draw,align=center,minimum width=2.25cm,minimum height=.82cm},
  derived/.style={box,fill=cobriaLight},
  line/.style={-{Latex[length=2mm]},thick}
]
\node[box] (asset) {Recurso y\\metadatos};
\node[derived,right=of asset] (extract) {Extractor\\versionado};
\node[derived,right=of extract] (chunks) {Fragmentos /\\descriptores};
\node[derived,below=of chunks] (embed) {Modelo de\\embedding};
\node[derived,left=of embed] (index) {Índice\\semántico};
\node[box,left=of index] (retrieve) {Recuperación\\autorizada};
\draw[line] (asset) -- (extract);
\draw[line] (extract) -- (chunks);
\draw[line] (chunks) -- (embed);
\draw[line] (embed) -- (index);
\draw[line] (index) -- (retrieve);
\end{tikzpicture}
\caption{Pipeline multimodal reconstruible.}
\label{fig:pipeline-multimodal}
\end{figure}

Cada etapa genera un manifiesto. Para una transcripción se conserva idioma, modelo,
parámetros y segmentos temporales; para imagen, regiones o descriptores; para video,
fotogramas o intervalos; para series, frecuencia y método de remuestreo. Los errores de
extracción no modifican el recurso fuente.

\section{Contrato de embedding}

Un embedding se representa como:

\begin{equation}
\begin{aligned}
z=\langle &sourceId,sourceVersion,chunkId,modelId,\\
           &modelVersion,d,normalization,vector,policy\rangle.
\end{aligned}
\end{equation}

La dimensión $d$, normalización y modelo determinan compatibilidad. Dos vectores con la
misma dimensión no son necesariamente comparables. Una migración crea un índice nuevo,
recalcula y valida recuperación antes del cambio de versión.

Los modelos basados en atención permiten representar dependencias contextuales y son la
base de numerosos modelos modernos de lenguaje y modalidades relacionadas
\parencite{vaswani2017attention}. COBRIA no depende de una arquitectura de modelo
específica; registra qué modelo produjo cada representación.

\section{Fragmentación}

El contrato de fragmentación declara unidad, solapamiento, límites semánticos, tamaño
máximo y tratamiento de tablas o marcas temporales. Un fragmento mantiene referencia a
la fuente y desplazamientos suficientes para mostrar evidencia. Si cambia el documento, los
fragmentos anteriores quedan asociados a su versión y se generan otros nuevos.

Fragmentos demasiado pequeños pierden contexto; demasiado grandes disminuyen precisión
de recuperación y aumentan costo de instrucción. La elección se evalúa por corpus y tarea,
no mediante una constante universal.

\section{Recuperación híbrida}

La búsqueda puede combinar filtros estructurados, texto y similitud vectorial. Primero
se limita por ámbito, permiso, vigencia y tipo; después se puntúa. Aplicar autorización
solo al final permite que un documento prohibido influya en ranking o contexto.

\begin{equation}
score(x,q)=\alpha s_{vector}+\beta s_{text}+\gamma s_{metadata},
\end{equation}

con pesos versionados. El resultado incluye score por componente, modelo, consulta
normalizada y referencias. El score ordena candidatos; no prueba verdad.

\section{Derecho a retirar y reconstrucción}

Cuando un recurso se retira, se invalidan fragmentos, embeddings, cachés y conjuntos de datos que
lo contienen según la política. El grafo de linaje permite localizar dependencias. Un
modelo ya entrenado plantea un problema diferente: puede requerir evaluación, nuevo
entrenamiento o documentación de la limitación, según obligación y riesgo. Borrar una
fila del índice no demuestra que toda influencia desapareció.

\chapter{Recuperación aumentada y agentes gobernados}

\section{RAG como composición de proyecciones}

La recuperación aumentada combina una memoria paramétrica con evidencia recuperada para
producir respuestas condicionadas por fuentes externas \parencite{lewis2020rag}. En
COBRIA, el corpus, los fragmentos y el índice son proyecciones; la respuesta es una
inferencia; las entidades y documentos autorizados conservan autoridad.

Una ejecución se formaliza como:

\begin{equation}
r=\operatorname{Generar}(M_v,p,\operatorname{Recuperar}(I_v,q,ambito,politica),\theta),
\end{equation}

donde $M_v$ es el modelo, $p$ la plantilla, $I_v$ el índice, $q$ la consulta y $\theta$
los parámetros. El registro conserva todas estas versiones junto con fragmentos
recuperados, scores, filtros, tiempos y resultado.

\section{Flujo seguro de recuperación}

\begin{enumerate}[leftmargin=*]
  \item autenticar actor y resolver ámbito;
  \item clasificar propósito y sensibilidad de la solicitud;
  \item construir filtros obligatorios antes de recuperar;
  \item recuperar candidatos y verificar vigencia;
  \item limitar y ordenar evidencia;
  \item generar con plantilla y modelo versionados;
  \item verificar formato, citas y políticas de salida;
  \item registrar evidencia y devolver respuesta con incertidumbre apropiada.
\end{enumerate}

No se envía al modelo más contexto del necesario. Las instrucciones presentes dentro de
documentos recuperados se tratan como contenido, no como autoridad para cambiar
políticas o ejecutar herramientas.

\section{Contrato de herramienta para agentes}

Una herramienta se declara como:

\begin{equation}
\begin{aligned}
\mathcal{T}=\langle &nombre,v,entrada,salida,capacidad,ambito,riesgo,\\
                     &confirmacion,idempotencia,compensacion\rangle.
\end{aligned}
\end{equation}

El agente no recibe credenciales de almacenamiento ni acceso general a Repositorio. Usa
herramientas estrechas que llaman casos de uso. El contrato especifica esquemas,
precondiciones, efectos, permisos, necesidad de confirmación y forma de compensar cuando
sea posible.

\begin{table}[H]
\centering
\caption{Niveles de autonomía propuestos.}
\label{tab:niveles-autonomia}
\begin{tabularx}{\textwidth}{p{2cm}p{4cm}X}
\toprule
Nivel & Capacidad & Control requerido \\
\midrule
A0 & Responder sin herramientas & Evidencia y filtrado de salida.\\
A1 & Consultar datos de bajo riesgo & Ámbito, mínimo privilegio y auditoría.\\
A2 & Preparar una propuesta & Vista previa y validación previa.\\
A3 & Ejecutar acción reversible & Confirmación, idempotencia y compensación.\\
A4 & Acción de alto impacto & Aprobación humana explícita y segregación.\\
\bottomrule
\end{tabularx}
\end{table}

El nivel se asigna por herramienta y contexto, no por agente completo. Una misma sesión
puede consultar un catálogo en A1 y requerir A4 para modificar una obligación financiera.

\section{Memoria de agentes}

La memoria de conversación, las preferencias inferidas y los resúmenes son proyecciones
con retención y propósito. No deben convertirse silenciosamente en perfil canónico. Una
preferencia confirmada puede persistirse mediante un caso de uso; una inferida conserva
marca de procedencia y expiración.

La memoria se segmenta por persona, organización y contexto. Recuperar memoria de otro ámbito
es una violación aunque el texto parezca inofensivo. Las instrucciones de borrado deben
propagarse a índices y cachés.

\section{Evidencia y verificabilidad}

El Registro de evidencias registra solicitud, actor, ámbito, política, modelo, instrucción versionado,
herramientas, entradas resumidas o protegidas, fuentes, salida, validaciones, decisión y
efectos. No debe almacenar secretos ni duplicar indiscriminadamente datos personales.
Los resúmenes criptográficos permiten verificar integridad sin revelar contenido.

Una respuesta con citas no es automáticamente correcta. Se evalúan relevancia de la
evidencia, soporte de afirmaciones, cobertura, vigencia y fidelidad de atribución. Las
pruebas incluyen preguntas sin respuesta para comprobar abstención.

\chapter{Gobierno, seguridad y gestión del riesgo}

\section{Clasificación de artefactos}

Cada artefacto se clasifica por autoridad, sensibilidad, propósito, retención,
reconstruibilidad y criticidad. La combinación determina controles. Un embedding puede
parecer opaco, pero puede revelar similitud o información sensible; hereda al menos la
clasificación de su fuente salvo evaluación que justifique otra cosa.

\begin{table}[H]
\centering
\caption{Gobierno según tipo de artefacto.}
\label{tab:gobierno-artefactos}
\begin{tabularx}{\textwidth}{p{3cm}p{3.2cm}X}
\toprule
Artefacto & Autoridad & Control destacado \\
\midrule
Canónico & Primaria & Invariantes, autorización y revisión.\\
Conjunto de datos & Derivada & Propósito, linaje, calidad y retención.\\
Característica & Derivada & Corrección temporal y equivalencia.\\
Representación vectorial & Derivada & Compatibilidad, permisos y retiro.\\
Modelo & Inferencial & Evaluación, alcance y documentación.\\
Predicción & Inferencial & Evidencia, vigencia e incertidumbre.\\
Acción de agente & Propuesta o efecto & Capacidad, confirmación y auditoría.\\
\bottomrule
\end{tabularx}
\end{table}

\section{Registro y documentación de modelos}

El Registro de modelos conserva identidad, versión, tarea, propietario, código, dependencias,
conjunto de datos, Conjunto de características, hiperparámetros, métricas, población evaluada, limitaciones, fecha,
estado y política de retiro. Las fichas de modelo ofrecen una estructura para comunicar
uso previsto, rendimiento y limitaciones \parencite{mitchell2019modelcards}.

Un modelo no pasa a producción solo por superar una métrica global. Se revisan segmentos
relevantes, calibración, robustez, privacidad, seguridad, costo y compatibilidad con el
caso de uso. La aprobación enlaza exactamente los artefactos evaluados.

\section{Separación de funciones}

COBRIA distingue productor de datos, propietario del contrato, desarrollador del modelo,
revisor de riesgo, aprobador y operador. En contextos pequeños una persona puede asumir
varias funciones, pero las decisiones quedan explícitas. Las acciones de alto impacto no
se autoaprueban desde el mismo proceso que generó la propuesta.

\section{Amenazas específicas}

El análisis incluye contaminación de datos, fuga entre organizaciones, inyección de instrucciones,
recuperación de contenido prohibido, extracción de información, dependencia de modelos,
deriva, abuso de herramientas, resultados fabricados y automatización excesiva. Para IA
generativa se consideran además confabulación, integridad de la información, propiedad
intelectual y riesgos de la interacción humano--máquina, en consonancia con el perfil de
IA generativa del AI RMF \parencite{nist2024genai,nist2023airmf}.

\section{Controles por frontera}

\begin{enumerate}[leftmargin=*]
  \item \textbf{Ingesta:} validación, análisis de programas maliciosos, procedencia y clasificación.
  \item \textbf{Transformación:} ejecución aislada, dependencias fijadas y resúmenes criptográficos.
  \item \textbf{Almacenamiento:} cifrado, segregación, mínimo privilegio y retención.
  \item \textbf{Recuperación:} filtros previos, límites y verificación de vigencia.
  \item \textbf{Inferencia:} modelos aprobados, plantillas versionadas y cuotas.
  \item \textbf{Acción:} herramientas tipadas, confirmación e idempotencia.
  \item \textbf{Observación:} registros minimizados, alertas y revisión de incidentes.
\end{enumerate}

\section{Ciclo de vida}

Los estados propuestos para conjuntos de datos y modelos son borrador, validación, aprobado,
activo, degradado, retirándose y retirado. Una transición requiere evidencia y actor.
Un artefacto degradado puede continuar disponible bajo restricciones o volver a una
versión segura; la política define la respuesta.

Los sistemas de aprendizaje acumulan deuda más allá del código: dependencias de datos,
configuración, consumidores ocultos y cambios de comportamiento pueden elevar el costo
de mantenimiento \parencite{sculley2015debt}. El grafo de linaje y los contratos de
COBRIA hacen visibles varias de estas dependencias, aunque no las eliminan.

\section{Monitoreo y respuesta}

El monitoreo separa salud del servicio, calidad de datos, desempeño del modelo, equidad,
seguridad y costo. Una alerta debe enlazar versión, población y ventana. La respuesta
puede detener publicación, volver a una versión, reducir autonomía, desactivar una
herramienta, reconstruir un derivado o solicitar revisión humana.

\chapter{Evaluación de adaptabilidad analítica e inteligente}

\section{Preguntas de evaluación}

La evaluación de la Parte IV busca responder:

\begin{enumerate}[leftmargin=*]
  \item ¿Puede el mismo modelo de entidad, ámbito y proyección representar datos
  empresariales, eventos temporales y contenido multimodal?
  \item ¿Qué proporción de un conjunto de datos y sus características posee linaje completo?
  \item ¿Puede reproducirse una versión sin acceder a datos fuera del corte temporal?
  \item ¿Cuánto difieren características fuera de línea y en línea?
  \item ¿La recuperación respeta autorización y aporta evidencia suficiente?
  \item ¿Los agentes ejecutan únicamente capacidades declaradas y confirmadas?
  \item ¿Qué costo adicional introducen gobierno, versionado y reconstrucción?
\end{enumerate}

\section{Escenarios anonimizados}

El Sistema A se representará mediante eventos, secuencias, métricas, texto y recursos
multimedia sintéticos. Permitirá evaluar ventanas, extracción, embeddings, recuperación
y cadenas de procesamiento de IA. El Sistema B utilizará entidades empresariales, catálogos,
configuraciones, documentos y operaciones multiinquilino sintéticas. Permitirá evaluar
Conjunto de característicass, aislamiento, RAG autorizado y herramientas de agentes.

La comparación no exige las mismas proyecciones en ambos sistemas. La adaptabilidad se
apoya si los principios y contratos permanecen, mientras las políticas y
transformaciones se especializan por dominio.

\section{Experimentos propuestos}

\begin{longtable}{>{\bfseries}p{1.4cm}p{4.7cm}p{8cm}}
\caption{Experimentos para analítica e inteligencia artificial.}\label{tab:experimentos-ia}\\
\toprule
ID & Experimento & Indicadores \\
\midrule
\endfirsthead
\toprule
ID & Experimento & Indicadores \\
\midrule
\endhead
AI-1 & Reconstruir conjunto de datos versionado & Resumen criptográfico, cobertura, tiempo, diferencias y linaje.\\
AI-2 & Introducir eventos tardíos & Registros corregidos, marca de corte y estabilidad.\\
AI-3 & Comparar fuera de línea y en línea & Error absoluto, tasa de divergencia y frescura.\\
AI-4 & Migrar modelo de embedding & Cobertura, coexistencia, recuperación y costo.\\
AI-5 & Evaluar RAG por ámbito & Precisión@k, evidencia, fugas y abstención.\\
AI-6 & Inyectar instrucciones en corpus & Ataques bloqueados y efectos no autorizados.\\
AI-7 & Ejecutar herramientas de agente & Aceptación, rechazo, confirmación e idempotencia.\\
AI-8 & Retirar una fuente & Dependencias detectadas y derivados invalidados.\\
AI-9 & Cambiar distribución de entrada & Deriva detectado, tiempo de respuesta y degradación.\\
AI-10 & Reproducir entrenamiento & Métricas, variación y artefactos faltantes.\\
\bottomrule
\end{longtable}

\section{Métricas de linaje y reproducibilidad}

La cobertura de linaje se define como:

\begin{equation}
L_c=\frac{\text{campos o artefactos con ruta completa a fuente}}
{\text{campos o artefactos evaluados}}.
\end{equation}

La reproducibilidad exacta utiliza proporción de registros idénticos; la aproximada
declara tolerancia por variable. También se mide completitud del manifiesto, dependencias
resueltas y tiempo hasta reconstrucción. Un $L_c$ alto no garantiza calidad de la fuente,
pero permite localizarla y responsabilizar su transformación.

\section{Métricas de recuperación}

Se utilizarán Precisión@k, Exhaustividad@k cuando exista conjunto relevante, rango recíproco,
proporción de afirmaciones respaldadas, corrección de citas, vigencia, abstención y tasa
de violaciones de ámbito. Las evaluaciones automáticas se complementan con revisión
humana ciega sobre una muestra y acuerdo entre evaluadores.

La respuesta final se descompone para evitar que fluidez o estilo oculten fallos de
recuperación. Se comparan búsqueda estructurada, vectorial e híbrida con el mismo corpus
y política.

\section{Métricas de agentes}

Se registran tareas completadas, pasos, herramientas llamadas, parámetros inválidos,
acciones rechazadas, confirmaciones solicitadas, efectos duplicados, compensaciones,
tiempo, costo y severidad del peor efecto. El éxito funcional no compensa una violación
de autorización.

\section{Criterios de adaptabilidad}

COBRIA se considera adaptable a una modalidad cuando: la fuente se expresa con identidad
y ámbito; la derivación posee contrato; el tiempo relevante puede representarse; existe
una política de acceso; el artefacto puede reconstruirse o su límite se documenta; y la
salida se integra mediante un caso de uso sin eludir autoridad. Si una modalidad exige
romper estas condiciones, se registra como límite de la propuesta.

\section{Resultados negativos previstos}

La investigación buscará activamente escenarios adversos: conjuntos de datos pequeños donde el
gobierno sea desproporcionado, embeddings cuya reconstrucción sea costosa, modelos no
deterministas con variación relevante, información histórica insuficiente, recuperación
que no mejore una búsqueda estructurada y tareas donde un agente añada riesgo sin valor.
Estos casos delimitan COBRIA y fortalecen la validez del manuscrito.

\section{Síntesis de la Parte IV}

Esta parte formalizó la adaptabilidad de COBRIA desde datos operacionales hasta minería
e inteligencia artificial. Los conjuntos de datos y características se definieron como
proyecciones versionadas; el linaje como un grafo navegable; los recursos multimodales,
fragmentos y embeddings como artefactos derivados; y RAG como una composición gobernada
de recuperación e inferencia.

También se estableció que un agente opera mediante herramientas tipadas y no mediante
acceso directo al almacenamiento. Las salidas probabilísticas conservan condición de
inferencia o propuesta hasta que un caso de uso las valida. Gobierno, seguridad,
documentación y monitoreo acompañan el ciclo completo, con controles proporcionales a
sensibilidad y autonomía.

Finalmente, se definieron diez experimentos que permiten evaluar linaje, corrección
temporal, equivalencia fuera de línea--en línea, recuperación, resistencia a instrucciones
maliciosas, herramientas de agentes y retiro de datos. Con ello queda preparada la
siguiente parte del documento: implementar una referencia neutral y describir, sin
identificarlos, cómo los dos sistemas materializan y adaptan estos principios.


\part{Implementación de referencia y casos de estudio}
\chapter{Método de análisis de las implementaciones}

Esta parte estudia cómo los principios de COBRIA aparecen en dos sistemas existentes sin
convertir el manuscrito en documentación de productos. Los casos se mantienen anonimizados
como Sistema A y Sistema B. El primero está orientado a eventos, secuencias, contenido
multimodal y componentes de inteligencia; el segundo es una plataforma empresarial
multiinquilino con un dominio amplio, operaciones críticas y una cantidad elevada de
entidades.

El análisis se realizó sobre una instantánea del código disponible el 20 de agosto de
2026. Los conteos representan archivos estructurales encontrados en esa fecha y no
usuarios, transacciones, cobertura funcional ni calidad. Los nombres de clases,
organizaciones, personas, rutas físicas, credenciales y reglas comerciales particulares
no se publican.

\section{Preguntas del estudio de implementación}

La observación busca responder cinco preguntas:

\begin{enumerate}[leftmargin=*]
  \item ¿Cómo se distribuyen modelo, acceso a datos, servicios, casos de uso, esquemas y
  catálogos?
  \item ¿Qué elementos corresponden a patrones conocidos y cuáles constituyen la
  integración COBRIA?
  \item ¿Dónde aparecen entidades canónicas y proyecciones o índices derivados?
  \item ¿Qué tan consistente es la aplicación de la estructura dentro de cada caso?
  \item ¿Qué cambios serían necesarios para implementar COBRIA de manera verificable y
  no solo nominal?
\end{enumerate}

\section{Fuentes de evidencia}

Se utilizaron archivos de objetos, modelos base, Daters, conectores, servicios, casos de
uso, esquemas, registries, catálogos, documentación técnica y pruebas existentes. La
presencia de un archivo indica intención estructural; solo el código ejecutable y las
pruebas pueden apoyar una afirmación de comportamiento.

\begin{table}[H]
\centering
\caption{Niveles de evidencia utilizados.}
\label{tab:niveles-evidencia-implementacion}
\begin{tabularx}{\textwidth}{p{2cm}p{4cm}X}
\toprule
Nivel & Evidencia & Interpretación permitida \\
\midrule
E0 & Nombre o documento & Intención; no prueba implementación.\\
E1 & Clase o contrato & Capacidad modelada.\\
E2 & Flujo conectado & Capacidad integrada en código.\\
E3 & Prueba automatizada & Comportamiento verificado en entorno controlado.\\
E4 & Telemetría reproducible & Comportamiento observado bajo carga declarada.\\
\bottomrule
\end{tabularx}
\end{table}

Este criterio evita atribuir funcionalidad completa a una clase porque su nombre
contenga términos como inteligencia, sincronización o auditoría. También impide usar el
número de archivos como sustituto de mantenibilidad.

\section{Procedimiento de inventario}

El procedimiento fue: identificar directorios de capas; clasificar archivos por función;
examinar bases y ejemplos representativos; buscar relaciones entre objeto y Dater;
revisar conectores y rutas; localizar índices, cachés, metadatos y servicios; y contrastar
el hallazgo con la documentación. Los resultados se sintetizaron en categorías neutrales.

Los conteos se calcularon mediante un criterio reproducible: archivos JavaScript bajo
los directorios de objetos, Daters, servicios, casos de uso, schemas y catálogos o datos
de referencia. Archivos generados y manuales pueden incrementar estas cifras; por eso se
reportan como tamaño estructural aproximado.

\section{Neutralidad y ética}

No se copiaron datos operacionales ni secretos. Los ejemplos de esta parte son
pseudocódigo o estructuras reducidas. Las descripciones de dominio se mantienen al nivel
necesario para explicar variación arquitectónica. Un lector puede reproducir la
implementación de referencia sin conocer la identidad de los sistemas originales.

\chapter{Implementación neutral de referencia}

\section{Objetivo}

La implementación neutral no busca reemplazar las aplicaciones estudiadas. Su función es
mostrar la unidad mínima de COBRIA sin depender de un marco de trabajo o proveedor. Debe ser
pequeña, ejecutable y suficientemente explícita para probar identidad, ámbito, mapeo,
proyección, reconstrucción y autorización.

\section{Estructura propuesta}

\begin{verbatim}
src/
  core/
    entities/          entidades canónicas
    value-objects/     valores inmutables
    policies/          invariantes de dominio
  application/
    use-cases/         operaciones autorizadas
    ports/             contratos de persistencia y eventos
  data/
    mappers/           dominio <-> representación física
    repositories/      acceso por identidad y consultas declaradas
    projections/       contratos, writers y rebuilders
    connectors/        adaptadores del proveedor
  metadata/
    schemas/           tipos y validación
    catalogs/          vocabularios y referencias
    relationships/     vínculos semánticos
  analytics/
    conjuntos de datos/          contratos y manifiestos
    características/          definiciones fuera de línea/en línea
    lineage/           grafo de procedencia
  intelligence/
    retrieval/         fragmentos, embeddings y búsqueda
    models/            registro y evaluación
    tools/             capacidades gobernadas para agentes
  governance/
    policies/          acceso, retención y clasificación
    evidence/          auditoría y evidencia
\end{verbatim}

La estructura expresa responsabilidades, no una obligación de crear un archivo por cada
elemento. Un sistema pequeño puede agrupar módulos siempre que mantenga las fronteras.

\section{Entidad canónica}

Una implementación mínima contiene identidad, ámbito, revisión, versión de esquema,
estado y atributos. El constructor valida forma básica; las transiciones con significado
de negocio se realizan mediante métodos o políticas.

\begin{verbatim}
class CanonicalEntity {
  constructor({ id, ámbito, revision = 0, schemaVersion = 1,
                status = "active", attributes = {} }) {
    requireId(id); requireScope(ámbito);
    this.id = id;
    this.ámbito = ámbito;
    this.revision = revision;
    this.schemaVersion = schemaVersion;
    this.status = status;
    this.attributes = structuredClone(attributes);
  }
}
\end{verbatim}

El objeto no conoce URL, SDK, tabla, colección ni caché. Esta independencia corresponde
al principio de separar dominio y persistencia y se relaciona con Mapeador de datos y
Repositorio \parencite{fowler2002poeaa}.

\section{Mapeador y Repositorio}

El Mapeador define correspondencia reversible entre nombres de dominio y campos físicos.
El Repositorio ofrece operaciones en lenguaje del dominio y exige ámbito en todas ellas.

\begin{verbatim}
class EntityMapper {
  toDomain(id, ámbito, crudo) { /* mapeo, valores predeterminados y versión */ }
  toRaw(entity)            { /* representación persistible */ }
}

class EntityRepositorio {
  get(ámbito, id)                 { /* lectura por clave */ }
  save(ámbito, entity, revision)  { /* control optimista */ }
  listBy(ámbito, query, cursor)   { /* consulta permitida */ }
}
\end{verbatim}

El Repositorio no debe devolver el SDK del proveedor ni permitir que la presentación
forme rutas. Un Conector traduce sus operaciones a la API concreta.

\section{Dater como variante de acceso a datos}

Los dos casos utilizan el término local \emph{Dater} para clases que convierten objetos,
ejecutan CRUD, manejan caché y, en algunos módulos, mantienen índices. El nombre no
designa un patrón reconocido distinto: funcionalmente combina responsabilidades de Data
Acceso Objeto, Repositorio, Mapeador de datos e Mapa de identidad. COBRIA puede conservar el término
por compatibilidad, pero recomienda declarar cuál de esas responsabilidades cumple cada
clase.

\begin{table}[H]
\centering
\caption{Correspondencia del Dater con patrones consolidados.}
\label{tab:dater-patrones}
\begin{tabularx}{\textwidth}{p{3.2cm}p{4cm}X}
\toprule
Responsabilidad observada & Patrón relacionado & Recomendación \\
\midrule
CRUD del proveedor & Data Acceso Objeto & Mantener detrás de un contrato.\\
Obtener entidad por identidad & Repositorio & Usar vocabulario del dominio.\\
Convertir crudo y objeto & Mapeador de datos & Separar mapeo cuando crece.\\
Conservar instancia por ID & Mapa de identidad & Definir ámbito e invalidación.\\
Listas con TTL & Caché & Medir frescura y fallos.\\
Índices derivados & Vista materializada & Versionar y reconstruir.\\
\bottomrule
\end{tabularx}
\end{table}

\section{Contrato de proyección ejecutable}

\begin{verbatim}
const ProjectionContract = {
  id: "entity-by-status",
  version: 2,
  sourceType: "Entidad",
  dependencies: [],
  consistency: "eventual",
  transform(entity) { return [[entity.status, entity.id, true]]; },
  verify(source, projection) { /* cobertura y huérfanos */ },
  rebuild({ ámbito, cursor })  { /* lotes y checkpoints */ }
};
\end{verbatim}

El contrato reúne propiedades que en las implementaciones aparecen distribuidas entre
rutas, métodos de Dater, procesos de trabajo y convenciones. Su versión permite crear una proyección
paralela y efectuar cambio de versión.

\section{Caso de uso y transacción lógica}

Un caso de uso coordina autorización, carga, invariantes, persistencia, cola de salida y resultado.

\begin{verbatim}
UpdateEntity(command, actor):
  policy.require(actor, "entity.update", command.ámbito)
  current <- repository.require(command.ámbito, command.id)
  updated <- domain.apply(current, command.change)
  repository.save(command.ámbito, updated, current.revision)
  cola de salida.append(EntityChanged(updated.id, updated.revision))
  evidence.record(command, actor, updated)
  return presenter.summary(updated)
\end{verbatim}

Si el motor permite actualización multipath atómica, canónico, índices críticos y
cola de salida pueden compartir una frontera. Si no, se documenta la compensación y se emplean
procesos de trabajo idempotentes.

\section{Pruebas mínimas}

La referencia requiere pruebas de ida y vuelta del Mapeador, aislamiento de ámbito,
revisión optimista, idempotencia, reconstrucción desde cero, reconstrucción incremental,
detección de huérfanos, cambio de versión y rechazo de escritura directa en una
proyección. Para analítica agrega corte temporal, linaje y equivalencia fuera de línea--en línea.

\chapter{Caso de estudio A: sistema orientado a eventos y contenido}

\section{Caracterización estructural}

En la fecha de corte, el inventario del Sistema A encontró aproximadamente 88 archivos
de objetos de dominio, 91 Daters, 182 servicios, siete casos de uso explícitos, dos
schemas centrales y 102 archivos de catálogos o datos de referencia. Estas cantidades
describen una base heterogénea: módulos antiguos conviven con subsistemas recientes más
formalizados.

El dominio combina participantes, grupos, sesiones, eventos, acciones evaluadas,
planificación, salud, inventario, archivos multimedia, facturación y artefactos de IA.
La variedad permite observar si una estructura común se adapta a información temporal,
conectada y multimodal.

\begin{table}[H]
\centering
\caption{Inventario neutral del Sistema A.}
\label{tab:inventario-sistema-a}
\begin{tabularx}{\textwidth}{p{3.5cm}p{2.2cm}X}
\toprule
Categoría & Conteo & Observación \\
\midrule
Objetos de dominio & 88 & Estilos de modelado de distintas generaciones.\\
Daters & 91 & CRUD, caché, serialización e índices.\\
Servicios & 182 & Operación, analítica, multimedia e IA.\\
Casos de uso explícitos & 7 & Concentrados en flujos recientes.\\
Esquemas centrales & 2 & La validación también aparece distribuida.\\
Catálogos y referencias & 102 & Vocabularios, configuración y datos compartidos.\\
\bottomrule
\end{tabularx}
\end{table}

\section{Objetos y serialización}

Los objetos observados usan clases JavaScript con constructores, getters, setters y, en
varios módulos, métodos de serialización. Algunos emplean campos privados y copias
defensivas para estructuras anidadas; otros exponen propiedades públicas y conservan
alias para compatibilidad. Esta coexistencia demuestra adaptabilidad, pero también riesgo
de semántica desigual.

En los módulos de inteligencia, los objetos representan versiones de modelos, métricas,
experimentos, consentimientos, tareas, auditorías, conocimiento, retroalimentación y solicitudes
de acción. Esto se alinea con la Parte IV porque el modelo y la decisión no se reducen a
un valor suelto: poseen identidad, estado, procedencia y tiempos.

\section{Daters e índices derivados}

Los Daters del Sistema A acceden al conector de datos, convierten cargas útiles, mantienen
Mapa locales y actualizan índices por modelo, tipo, estado u otras dimensiones. Algunos
construyen una actualización multipath que persiste el registro principal y varias rutas
derivadas. Este es un ejemplo concreto del intercambio de COBRIA: una escritura lógica
produce varias escrituras físicas para acelerar lecturas.

\begin{figure}[H]
\centering
\begin{tikzpicture}[
  node distance=.8cm and .85cm,
  box/.style={draw,align=center,minimum width=2.3cm,minimum height=.82cm},
  derived/.style={box,fill=cobriaLight},
  line/.style={-{Latex[length=2mm]},thick}
]
\node[box] (object) {Objeto de\\dominio};
\node[box,right=of object] (dater) {Dater};
\node[box,right=of dater] (connector) {Conector\\NoSQL};
\node[derived,below=of connector] (base) {Ruta\\principal};
\node[derived,left=of base] (indexes) {Índices por\\dimensión};
\node[derived,left=of indexes] (cache) {Mapa e índices\\locales};
\draw[line] (object) -- (dater);
\draw[line] (dater) -- (connector);
\draw[line] (connector) -- (base);
\draw[line] (connector) -- (indexes);
\draw[line] (dater) -- (cache);
\end{tikzpicture}
\caption{Flujo de datos observado y abstraído en el Sistema A.}
\label{fig:flujo-sistema-a}
\end{figure}

La ruta principal se aproxima al canónico y los índices a proyecciones, pero esa
clasificación debe verificarse por entidad. Un carga útil duplicado puede contener campos
que otro módulo modifica directamente; en tal caso existe autoridad compartida y no una
proyección COBRIA válida.

\section{Servicios, casos de uso e IA}

La capa de servicios incluye coordinación de eventos, análisis, reportes, sincronización,
multimedia, salud, facturación e inteligencia. Los casos de uso explícitos aparecen en
operaciones que requieren transiciones y permisos. En otros módulos, un servicio o Dater
todavía concentra la coordinación.

La capacidad de IA es estructuralmente relevante: existen artefactos para versiones,
experimentos, métricas, preferencias, consentimiento, conocimiento, auditoría, retroalimentación
y acciones propuestas. Esto apoya la adaptabilidad conceptual de COBRIA, pero no prueba
exactitud de modelos ni seguridad de agentes. Esas capacidades deben alcanzar E3 o E4
mediante los experimentos de la Parte IV.

\section{Fortalezas observadas}

\begin{itemize}[leftmargin=*]
  \item separación visible entre presentación, lógica y datos;
  \item objetos especializados para eventos, multimedia e inteligencia;
  \item Daters con cachés e índices adecuados para consultas frecuentes;
  \item actualizaciones multipath en operaciones con derivados;
  \item servicios para analítica y ciclo de vida de artefactos;
  \item adopción progresiva de casos de uso en flujos con reglas.
\end{itemize}

\section{Deuda y oportunidades}

Los principales riesgos son estilos de objetos heterogéneos, validación distribuida,
responsabilidades amplias en ciertos Daters, índices sin contrato uniforme, rutas
conocidas por múltiples módulos y cobertura desigual de casos de uso. La optimización
local mediante Mapa puede perder ámbito o frescura si la clave e invalidación no se
declaran.

La migración recomendada no reescribe todos los módulos. Se priorizan entidades con
datos sensibles, múltiples índices, IA o alto volumen. Para cada una se registra fuente
canónica, contrato de proyección, versión, propietario, reconstrucción y pruebas. Los
alias de compatibilidad pueden permanecer en el Mapeador mientras se estabiliza el modelo.

\section{Evolución longitudinal posterior del Sistema A}

Una segunda observación del repositorio, posterior al inventario inicial, permitió
analizar cambio arquitectónico y no solo estructura estática. La evolución no se usa
como un nuevo prueba comparativa de rendimiento; constituye evidencia de aplicabilidad de los
principios de COBRIA en una migración real y gradual.

\begin{table}[H]
\centering
\caption{Evidencia longitudinal neutral del Sistema A.}
\label{tab:evolucion-sistema-a}
\begin{tabularx}{\textwidth}{p{4.1cm}p{2.3cm}X}
\toprule
Propiedad observada & Evidencia & Interpretación COBRIA \\
\midrule
Frontera Presentación--Data & 0 importaciones directos & La dependencia se hace verificable
por CI, no solo convencional.\\
Compatibilidad temporal & 134 consumidores inventariados & La deuda se congela y solo
puede disminuir.\\
Vocabulario físico central & 36 nodos semánticos & Los nombres compactos quedan detrás
de un registro de rutas.\\
Modelos de lectura explícitos & 4 repositorios de vistas & Listados y resúmenes declaran
almacén derivado separado.\\
Pruebas arquitectónicas y de escala & 16 artefactos relevantes & Se verifican rutas,
concurrencia, vistas, presupuestos y fronteras.\\
\bottomrule
\end{tabularx}
\end{table}

Las nuevas vistas de listas y resúmenes incluyen \texttt{complete},
\texttt{updatedAt}, TTL y operaciones de reemplazo, invalidación o reconstrucción.
Una vista fría agrupa solicitudes concurrentes en una sola reconstrucción; las pruebas
comprueban que veinte solicitudes comparten un único warmup. La autorización se ejecuta
en cada solicitud contra la política canónica, aun cuando el resultado materializado
esté fresco. Este detalle separa optimización de autoridad.

Las operaciones de dominio aplican actualizaciones multipath para canónico, índices,
auditoría y notificación; utilizan versiones, transiciones explícitas, idempotencia y
controles concurrentes para invariantes. Además, el contexto utilizado por telemetría y
FinOps excluye valores no validados provenientes del cliente. Estas decisiones amplían
COBRIA con dos reglas: \emph{la optimización puede fallar sin falsificar el canónico} y
\emph{la atribución operacional debe proceder del ámbito autorizado}.

La migración del interfaz aporta otra lección. Una frontera temporal puede conservar
134 consumidores heredado sin permitir nuevos consumidores. El manifiesto y el límite
automático convierten el patrón estrangulador en una transición medible. COBRIA debe, por
tanto, exigir a cada puente: propietario, inventario, fecha o condición de retiro,
presupuesto que no crezca y pruebas de delegación.

\section{Implicación para la validez del caso}

La observación longitudinal fortalece H8 porque el mismo sistema evolucionó hacia rutas
semánticas, servicios de aplicación, políticas separadas, proyecciones materializadas y
controles automatizados arquitectónicos sin una reescritura total. No demuestra causalidad: los cambios
fueron realizados por participantes conocedores de la propuesta y no se compararon
contra un equipo de control. Sí proporciona evidencia concreta de que COBRIA puede
funcionar como arquitectura evolutiva y no únicamente como taxonomía retrospectiva.

\chapter{Caso de estudio B: sistema empresarial multiinquilino}

\section{Caracterización estructural}

El Sistema B posee una escala estructural considerablemente mayor. En la fecha de corte
se encontraron aproximadamente 680 objetos de dominio, 680 Daters, 1,624 servicios, 18
casos de uso explícitos, 686 schemas y 514 archivos de catálogos o datos de referencia.
La correspondencia cercana entre objetos, Daters y schemas revela una estrategia
sistemática, parte de ella apoyada por generación o convenciones repetibles.

\begin{table}[H]
\centering
\caption{Inventario neutral del Sistema B.}
\label{tab:inventario-sistema-b}
\begin{tabularx}{\textwidth}{p{3.5cm}p{2.2cm}X}
\toprule
Categoría & Conteo & Observación \\
\midrule
Objetos de dominio & 680 & Cobertura de numerosos subdominios.\\
Daters & 680 & Correspondencia casi uno a uno con objetos.\\
Servicios & 1,624 & Políticas, procesos de trabajo, contratos y orquestadores.\\
Casos de uso explícitos & 18 & Operaciones críticas seleccionadas.\\
Esquemas & 686 & Validación estructural extensa.\\
Catálogos y referencias & 514 & Vocabularios y configuración por dominio.\\
\bottomrule
\end{tabularx}
\end{table}

Estos números no permiten concluir que todo módulo esté integrado o probado. Sí ofrecen
evidencia de una arquitectura dirigida por estructura y metadatos, adecuada para estudiar
escalabilidad organizacional.

\section{Modelo base y mapeo físico}

Un modelo base representativo inicializa identidad y valores predeterminados, restringe campos
desconocidos, serializa a JSON y mantiene un mapa explícito entre campo físico y campo de
dominio. El método inverso reconstruye una instancia desde datos crudo. Esta forma
materializa Mapeador de datos de manera directa.

La ventaja es que una entidad puede conservar vocabulario estable aunque cambien claves
físicas. El riesgo aparece cuando valores predeterminados silenciosos ocultan ausencia o versiones
antiguas. COBRIA recomienda distinguir campo ausente, nulo y vacío y registrar la versión
del esquema utilizada en la reconstrucción.

\section{Dater base}

El Dater base observado recibe ruta, modelo, conector y TTL. Implementa lectura por ID,
alta, actualización, borrado, lista, paginación por cursor, consulta por hijo, carga
múltiple, escucha en tiempo real, caché de entidades y caché de listas. Convierte entre
crudo y objeto mediante el modelo.

\begin{figure}[H]
\centering
\begin{tikzpicture}[
  node distance=.75cm and .8cm,
  box/.style={draw,align=center,minimum width=2.25cm,minimum height=.82cm},
  derived/.style={box,fill=cobriaLight},
  line/.style={-{Latex[length=2mm]},thick}
]
\node[box] (service) {Servicio /\\caso de uso};
\node[box,right=of service] (dater) {Dater base\\especializado};
\node[box,right=of dater] (mapper) {Modelo y\\mapa de campos};
\node[box,below=of mapper] (connector) {Conector};
\node[derived,left=of connector] (store) {Almacenamiento\\multiinquilino};
\node[derived,left=of store] (cache) {Caché con\\TTL};
\draw[line] (service) -- (dater);
\draw[line] (dater) -- (mapper);
\draw[line] (dater) -- (connector);
\draw[line] (connector) -- (store);
\draw[line] (dater) -- (cache);
\end{tikzpicture}
\caption{Abstracción del acceso común observado en el Sistema B.}
\label{fig:flujo-sistema-b}
\end{figure}

Esta base reduce duplicación y ofrece una API uniforme. Sin embargo, el argumento de
ruta debe incorporar ámbito de forma no opcional. Una ruta base global o una clave de
caché basada solo en ID podría romper aislamiento aunque el resto de la aplicación
conozca el organización.

\section{Metadatos, relaciones y catálogos}

El Sistema B contiene registries para políticas de campo, relaciones, clasificaciones de
referencias y correspondencia con catálogos. Los campos pueden declararse configurables,
externos o asociados con vocabularios. Esta capa convierte reglas que normalmente
quedarían dispersas en datos inspeccionables.

La arquitectura dirigida por metadatos favorece generación de objetos, schemas, Daters,
documentación y controles de interfaz. Para que la generación sea confiable debe existir
una fuente primaria del metamodelo, validación del generador y detección de divergencia.
Editar artefactos generados manualmente sin retroalimentar el metamodelo crea deriva.

\section{Servicios y operaciones críticas}

Los servicios abarcan consulta, disponibilidad, inventario, finanzas, comunicaciones,
integraciones, privacidad, sincronización, reportes y operaciones multiinquilino. Existen
contratos de repositorio, procesos de trabajo, políticas, state machines y orquestadores. Los casos de
uso explícitos se concentran en operaciones donde estados, permisos y efectos derivados
requieren una frontera clara.

Los contratos operacionales muestran ámbito de organización y unidad, claves de
idempotencia, autenticación, recálculo en servidor, rate limits y errores semánticos.
Estos mecanismos coinciden con COBRIA Núcleo y Gobierno. La persistencia de cola de salida,
operaciones idempotentes, versiones de esquema, bloqueos y trabajos programados ofrece
puntos naturales para reconstrucción y observabilidad.

\section{Fortalezas observadas}

\begin{itemize}[leftmargin=*]
  \item correspondencia sistemática entre objeto, Dater y esquema;
  \item modelo base con mapeo entre dominio y representación física;
  \item Dater base con caché, paginación y tiempo real;
  \item metadatos explícitos para campos, relaciones y catálogos;
  \item contratos de repositorio y casos de uso en operaciones críticas;
  \item idempotencia, cola de salida, procesos de trabajo y controles multiinquilino modelados;
  \item suficiente variedad de dominios para evaluar extensibilidad.
\end{itemize}

\section{Deuda y oportunidades}

La escala produce riesgos de proliferación: una clase por estructura puede aumentar
costo de navegación, generación y pruebas. Una correspondencia uno a uno no garantiza
que cada estructura posea identidad o comportamiento suficiente para ser entidad. Los
1,624 servicios pueden incluir políticas pequeñas y artefactos generados; deben
agruparse por responsabilidad y dependencia antes de inferir complejidad.

COBRIA recomienda clasificar cada tipo como entidad, value object, evento, comando,
proyección o catálogo. También propone registrar responsabilidad de proyecciones, reconstrucción
y presupuesto. La base común debe exigir ámbito, revisión y versión; los servicios deben
depender de contratos, no de rutas ni SDKs.

\chapter{Comparación entre casos y derivación de COBRIA}

\section{Similitudes estructurales}

Ambos sistemas separan presentación, lógica y datos; representan el dominio mediante
objetos; encapsulan persistencia en Daters; utilizan servicios; conservan catálogos y
datos compartidos; y desnormalizan para responder consultas. Ambos poseen conectores a
un almacenamiento NoSQL y cachés locales. Estas similitudes motivaron la búsqueda de un
vocabulario común.

\begin{longtable}{>{\bfseries}p{3.2cm}p{5.6cm}p{5.6cm}}
\caption{Comparación neutral de los casos.}\label{tab:comparacion-casos}\\
\toprule
Dimensión & Sistema A & Sistema B \\
\midrule
\endfirsthead
\toprule
Dimensión & Sistema A & Sistema B \\
\midrule
\endhead
Dominio & Eventos, secuencias, multimedia e inteligencia. & Operaciones empresariales
multiinquilino de gran amplitud.\\
Objetos & Heterogéneos y especializados. & Sistemáticos, con modelo base frecuente.\\
Daters & Especializados; varios mantienen índices propios. & Base reutilizable y amplia
correspondencia por entidad.\\
Esquemas & Parciales o distribuidos. & Cobertura estructural extensa.\\
Metadatos & Catálogos y configuraciones por módulo. & Registries de campos, relaciones y
catálogos.\\
Casos de uso & Adopción focalizada en flujos recientes. & Adopción focalizada en
operaciones críticas.\\
Proyecciones & Índices locales/remotos y artefactos analíticos. & Índices, reportes,
cachés, cola de salida y snapshots.\\
IA & Componentes explícitos de versión, auditoría y retroalimentación. & Preparación estructural,
automatización y datos empresariales.\\
Riesgo principal & Heterogeneidad y contratos implícitos. & Proliferación y divergencia
entre artefactos.\\
\bottomrule
\end{longtable}

\section{Diferencias que prueban adaptabilidad}

El Sistema A necesita representar tiempo fino, acciones, secuencias y archivos. El
Sistema B necesita ámbito jerárquico, catálogos configurables, operaciones idempotentes y
muchas relaciones. COBRIA no intenta igualar sus modelos. La adaptación ocurre porque
ambos pueden expresar identidad, ámbito, versión, mapeo y proyección, pero emplean
contratos diferentes.

En A, una proyección puede indexar eventos por participante, contexto o ventana y un
pipeline puede generar clips o características. En B, una proyección puede resolver
disponibilidad, saldos, reportes o búsquedas por organización. El mecanismo es común; la
semántica pertenece al dominio.

\section{Elementos heredados y contribución}

Los objetos, DAO/Repositorio, Mapeador de datos, Capa de servicios, casos de uso, caché, catálogos,
vistas materializadas, cola de salida y metadatos son conceptos existentes. El Dater es un nombre
local para una combinación de esos patrones, no una nueva categoría científica.

La contribución COBRIA se ubica en la integración de estas ideas alrededor de:

\begin{enumerate}[leftmargin=*]
  \item una fuente canónica con identidad, ámbito, revisión y versión;
  \item una separación explícita entre dominio y forma física;
  \item proyecciones con contrato, responsabilidad, consistencia y reconstrucción;
  \item metadatos que conectan modelo, esquema, relaciones y catálogos;
  \item continuidad de linaje desde operación hasta conjunto de datos, modelo y decisión;
  \item una frontera que impide a derivados o agentes adquirir autoridad por accidente.
\end{enumerate}

\section{Patrones de diseño identificados}

\begin{table}[H]
\centering
\caption{Patrones observados y papel dentro de COBRIA.}
\label{tab:patrones-observados}
\begin{tabularx}{\textwidth}{p{3.4cm}p{4.2cm}X}
\toprule
Patrón & Evidencia típica & Papel COBRIA \\
\midrule
Repositorio/DAO & Daters y contratos de repositorio & Aislar persistencia.\\
Mapeador de datos & fromRaw, toRaw y mapas de campos & Separar semántica y forma física.\\
Mapa de identidad & Mapas por identidad & Evitar instancias duplicadas en sesión.\\
Capa de servicios & Servicios de dominio y aplicación & Coordinar operaciones.\\
Caso de uso & Flujos explícitos y autorizados & Definir frontera transaccional.\\
Vista materializada & Índices, resúmenes y snapshots & Optimizar lectura con reconstrucción.\\
CQRS parcial & Modelos de consulta especializados & Separar necesidades de lectura.\\
Cola de salida & Eventos pendientes y procesos de trabajo & Reducir ventana de pérdida.\\
Registro & Metadatos y catálogos & Centralizar contratos inspeccionables.\\
Estrategia y política & Validadores y políticas intercambiables & Variar reglas sin rutas físicas.\\
Adaptador/Pasarela & Conectores y proveedores externos & Encapsular infraestructura.\\
Máquina de estados & Transiciones de operaciones críticas & Restringir ciclos de vida.\\
\bottomrule
\end{tabularx}
\end{table}

\section{Modelo de madurez}

\begin{table}[H]
\centering
\caption{Niveles de madurez COBRIA para una entidad.}
\label{tab:madurez-cobria}
\begin{tabularx}{\textwidth}{p{1.5cm}p{3.2cm}X}
\toprule
Nivel & Nombre & Evidencia \\
\midrule
M0 & Acceso directo & interfaz de usuario o servicio conoce rutas y cargas útiles.\\
M1 & Encapsulado & Objeto y Dater/Repositorio básico.\\
M2 & Canónico & Identidad, ámbito, mapeo y esquema explícitos.\\
M3 & Reconstruible & Proyecciones versionadas con rebuild y verificación.\\
M4 & Gobernado & Linaje, políticas, observabilidad y recuperación.\\
M5 & Adaptativo & Optimización basada en evidencia con aprobación.\\
\bottomrule
\end{tabularx}
\end{table}

La madurez se evalúa por entidad o flujo, no por sistema completo. Ambos casos pueden
contener simultáneamente módulos M1 y M4. Esta granularidad permite una migración basada
en riesgo.

\chapter{Estrategia de adopción y migración}

\section{Adopción incremental}

La estrategia recomendada sigue siete pasos:

\begin{enumerate}[leftmargin=*]
  \item seleccionar una entidad con dolor verificable;
  \item declarar identidad, ámbito, autoridad e invariantes;
  \item aislar Mapeador y Repositorio o documentar el Dater existente;
  \item inventariar copias, índices, cachés y consumidores;
  \item crear contratos de proyección y verificación;
  \item instrumentar lecturas, escrituras, anomalías y costo;
  \item reconstruir en paralelo, comparar y efectuar cambio de versión reversible.
\end{enumerate}

No se recomienda renombrar masivamente los Daters. La migración comienza por contratos y
propiedades verificables. El nombre puede cambiar cuando aporte claridad y no rompa
consumidores.

\section{estrangulador para datos}

Una fachada compatible intercepta operaciones antiguas y las dirige al nuevo Repositorio.
Durante la transición puede comparar lecturas o emitir eventos para nuevas proyecciones.
La escritura dual sin verificación aumenta riesgo; se prefiere una fuente autoritativa y
un derivado reconstruible.

\begin{figure}[H]
\centering
\begin{tikzpicture}[
  node distance=.85cm and .9cm,
  box/.style={draw,align=center,minimum width=2.35cm,minimum height=.82cm},
  derived/.style={box,fill=cobriaLight},
  line/.style={-{Latex[length=2mm]},thick}
]
\node[box] (heredado) {Consumidor\\existente};
\node[box,right=of heredado] (facade) {Fachada\\compatible};
\node[derived,right=of facade] (usecase) {Nuevo caso\\de uso};
\node[derived,below=of usecase] (repo) {Repositorio y\\Mapeador};
\node[box,left=of repo] (canonical) {Canónico};
\node[derived,left=of canonical] (projection) {Proyección\\versionada};
\draw[line] (heredado) -- (facade);
\draw[line] (facade) -- (usecase);
\draw[line] (usecase) -- (repo);
\draw[line] (repo) -- (canonical);
\draw[line] (canonical) -- (projection);
\end{tikzpicture}
\caption{Migración incremental mediante una fachada compatible.}
\label{fig:migracion-strangler}
\end{figure}

\section{Priorización por riesgo y valor}

Se asigna una puntuación basada en sensibilidad, frecuencia, amplificación, número de
copias, incidentes, dificultad de reconstrucción y consumidores. Las entidades con alto
riesgo y alto uso se atienden primero. Un catálogo estático y estable puede permanecer
en M1 o M2 si no existe beneficio neto de mayor formalidad.

\section{Generación dirigida por metadatos}

En el Sistema B, la escala hace atractiva la generación. El metamodelo puede producir
esqueleto de entidad, esquema, Mapeador, Dater/Repositorio, documentación y pruebas de
contrato. La generación debe ser determinista; el repositorio rechaza divergencias entre
salida generada y archivos versionados.

El generador no decide invariantes de negocio ni crea proyecciones por cada campo. Esas
decisiones requieren demanda de consulta, responsable y presupuesto.

\section{Compatibilidad y versionado}

Los cambios se clasifican como aditivos, deprecatorios o incompatibles. Un Mapeador puede
leer varias versiones y escribir solo la vigente. La migración se ejecuta por lotes y
con checkpoint. Las proyecciones incluyen versión en ruta o clave para permitir
coexistencia y reversión.

\section{Presupuesto de deuda y retiro de puentes}

Toda compatibilidad temporal se registra en un manifiesto reproducible. Sea \(L_t\) el
conjunto de consumidores heredado permitidos en la versión \(t\). La regla de integración
continua exige:

\[
L_{t+1}\subseteq L_t.
\]

Un consumidor nuevo requiere migración al contrato vigente, no ampliación silenciosa
del manifiesto. El puente solo reexporta o adapta; no recibe nuevas reglas de negocio. Su
eliminación se realiza por dominio, con pruebas de comportamiento y reducción del
presupuesto después de cada bloque estable.

\section{Criterio de finalización}

Una entidad se considera migrada cuando: toda escritura autorizada atraviesa casos de
uso definidos; la fuente canónica es única; el ámbito es obligatorio; el ida y vuelta del
Mapeador está probado; cada derivado posee contrato y responsable; una reconstrucción completa
produce el resultado esperado; la telemetría permite medir costo y anomalías; y los
controles automatizados impiden reintroducir rutas físicas, dependencias invertidas o consumidores heredado.

\chapter{Reproducibilidad, limitaciones y síntesis de la Parte V}

\section{Paquete reproducible previsto}

El manuscrito debe acompañarse de un paquete que no contenga los proyectos privados. El
paquete incluirá implementación neutral, generadores sintéticos, contratos de ejemplo,
scripts de inventario, pruebas, configuración del prueba comparativa y plantillas de resultados.
Cada versión publicada fijará versiones y resúmenes criptográficos.

\begin{table}[H]
\centering
\caption{Contenido del paquete reproducible.}
\label{tab:paquete-reproducible}
\begin{tabularx}{\textwidth}{p{3.6cm}X}
\toprule
Artefacto & Contenido \\
\midrule
Referencia & Núcleo, Repositorio, Mapeador, Gestor de proyecciones y Gobierno.\\
Datos sintéticos & Perfiles A y B con semilla y escalas configurables.\\
Contratos & Entidades, proyecciones, conjuntos de datos y herramientas.\\
Migraciones & Total, incremental, cambio de versión y reversión.\\
Pruebas & Aislamiento, idempotencia, ida y vuelta y reconstrucción.\\
Prueba comparativa & Cargas, instrumentación y análisis estadístico.\\
Documentación & Diccionario, decisiones y guía de réplica.\\
\bottomrule
\end{tabularx}
\end{table}

\section{Limitaciones del análisis estático}

El inventario no demuestra qué archivos se ejecutan en producción, qué rutas contienen
datos, cuántos módulos tienen consumidores ni qué cobertura alcanzan las pruebas. Los
conteos pueden incluir archivos generados, contratos pequeños o capacidades parciales.
La documentación puede estar desactualizada. Estas limitaciones se resolverán en la
Parte VI mediante ejecución controlada, trazas y experimentos.

\section{Riesgo de sesgo del investigador}

El investigador es también autor o conocedor de los sistemas, lo que facilita acceso y
comprensión, pero puede favorecer interpretaciones positivas. Se mitigará publicando
criterios previos, scripts, resultados negativos, niveles de evidencia y una
implementación neutral. Cuando sea posible, otro evaluador revisará la clasificación de
una muestra de módulos.

\section{Validez de la generalización}

Los casos comparten lenguaje y almacenamiento NoSQL, por lo que no prueban adaptación a
todos los ecosistemas. La contribución generalizable es el mecanismo y sus condiciones,
no los nombres de carpetas ni las cantidades. Una réplica futura debe incluir otro
lenguaje, otro motor documental y un equipo independiente.

\section{Síntesis de la Parte V}

Esta parte conectó la formalización de COBRIA con evidencia estructural de dos sistemas
reales y distintos. El Sistema A mostró objetos heterogéneos, Daters con índices y una
amplia superficie de eventos, multimedia e inteligencia. El Sistema B mostró una
correspondencia masiva entre objetos, Daters y schemas, un modelo base con mapeo físico,
metadatos y operaciones multiinquilino más sistemáticas.

El análisis confirmó que ``Objeto--Dater--Servicio--Catálogo'' no constituye por sí solo
un patrón nuevo. Es una composición local de Modelo de dominio, DAO/Repositorio, Mapeador de datos,
Capa de servicios, Registro, Caché y vistas materializadas. COBRIA adquiere interés de
investigación al convertir esa composición en una arquitectura explícita, formal y
evaluada alrededor de autoridad canónica, proyecciones reconstruibles, ámbito y linaje
hacia IA.

También se definió una implementación neutral, un modelo de madurez y una estrategia de
migración incremental. Los conteos estructurales sirven para caracterizar, no para
celebrar tamaño. La siguiente parte ejecutará el protocolo: medirá rendimiento,
reconstrucción, aislamiento, evolución, linaje y adaptabilidad, y presentará tanto
resultados favorables como negativos.


\part{Evaluación experimental, análisis y validez}
\chapter{Diseño de evaluación exclusivamente NoSQL}

\section{Delimitación tecnológica}

Esta edición evalúa COBRIA únicamente dentro de sistemas NoSQL documentales. La unidad
de persistencia es el documento u objeto agregado; las consultas se expresan mediante
claves, filtros, índices del motor, colecciones derivadas y flujos de cambios. No se
transfieren cifras obtenidas con motores de otro paradigma. Esta decisión reduce el
volumen de evidencia cuantitativa disponible, pero alinea la inferencia con el objeto
real de investigación.

La evaluación distingue tres niveles. El nivel estructural verifica contratos en código
y configuración. El nivel ejecutable usa emuladores y entornos NoSQL aislados. El nivel
empírico mide latencia, costo, frescura y recuperación en condiciones reproducibles.
Una afirmación no asciende de nivel sin la evidencia correspondiente.

\section{Casos documentales anonimizados}

Se estudian tres aplicaciones con dominios distintos: eventos deportivos y multimedia,
operación hotelera y operación comercial. Sus nombres se omiten. Las tres emplean
documentos, colecciones, autenticación, funciones y almacenamiento administrado. Se
buscan contratos repetidos sin revelar organizaciones, personas ni credenciales.

\begin{table}[H]
\centering
\caption{Perfiles NoSQL utilizados en la evaluación.}
\begin{tabularx}{\textwidth}{p{2.4cm}X X}
\toprule
Perfil & Carga dominante & Riesgo arquitectónico observado\\
\midrule
N1 & eventos, resultados y multimedia & derivados obsoletos y reconstrucción costosa\\
N2 & reservas, disponibilidad y catálogos & concurrencia, tiempo y aislamiento\\
N3 & inventario, ventas y cierres & agregados operacionales y trazabilidad\\
\bottomrule
\end{tabularx}
\end{table}

\chapter{Hipótesis, variables y protocolo}

\section{Hipótesis activas}

\begin{enumerate}[leftmargin=*]
  \item[HN1.] Una proyección declarada por contrato puede reconstruirse desde documentos
  canónicos sin convertirse en autoridad paralela.
  \item[HN2.] El ámbito obligatorio impide recuperar documentos o fragmentos de otra
  organización antes de calcular similitud o construir contexto.
  \item[HN3.] Revisión, idempotencia y publicación por versión reducen duplicados,
  regresiones y estados parcialmente visibles.
  \item[HN4.] Un conjunto analítico versionado puede reproducirse desde fuentes,
  catálogos y transformaciones declaradas.
  \item[HN5.] Una proyección solo se justifica si mejora un objetivo de lectura y su
  mantenimiento permanece dentro de un presupuesto explícito.
\end{enumerate}

HN1--HN4 admiten verificación funcional. HN5 necesita mediciones específicas por motor,
región, índice, escala y mezcla de operaciones. Esta edición no publica un umbral
universal de lecturas por escritura.

\section{Variables y unidades}

Las variables independientes son motor documental, perfil, escala, número de ámbitos,
forma del documento, política de consistencia, estrategia de reconstrucción y nivel de
concurrencia. Las dependientes incluyen p50, p95 y p99; documentos y bytes leídos;
operaciones facturables; frescura; duplicados; divergencia; duración de reconstrucción;
violaciones de ámbito; y exactitud del conjunto analítico.

La corrida completa es la unidad inferencial. Cada manifiesto registra semilla, versión
de código, configuración, región, índices, calentamiento, orden de cargas y criterio de
descarte. Las enmiendas se publican antes de observar el resultado que podrían favorecer.

\chapter{Experimentos NoSQL}

\section{P1: lectura selectiva documental}

P1 compara la consulta canónica indexada con una proyección de cobertura sobre el mismo
conjunto lógico. Ambas configuraciones usan el mismo motor NoSQL y los mismos documentos
de entrada. La proyección gana únicamente si reduce el percentil objetivo o el costo
facturable sin alterar resultados, ámbito ni frescura más allá del límite acordado.

\section{P2: amplificación de escritura}

P2 actualiza documentos canónicos y registra operaciones adicionales, tamaño de eventos,
invocaciones, reintentos y tiempo hasta convergencia. El resultado es un vector de costo,
no una sola cifra. Una reducción de lectura no compensa automáticamente mayor complejidad.

\section{R1: reconstrucción total e incremental}

R1 elimina una versión derivada en un entorno controlado, reconstruye una candidata,
verifica equivalencia y publica mediante un cambio de versión. La variante incremental
procesa revisiones posteriores a un punto de control. Ambas rutas deben ser idempotentes,
conservar procedencia y producir el mismo resultado lógico.

\section{S1: aislamiento multiinquilino}

S1 genera identidades, documentos y consultas para múltiples ámbitos. Incluye ausencia
de ámbito, ámbito manipulado, claves de caché incompletas, búsqueda semántica y tareas
administrativas. La aceptación exige rechazo temprano y cero documentos cruzados. El
filtro posterior a la recuperación no satisface el contrato.

\section{A1: analítica e inteligencia artificial}

A1 reconstruye un conjunto fechado, publica variables versionadas y ejecuta recuperación
autorizada. Cada fragmento conserva fuente, ámbito, revisión, versión de transformación
y política. Se miden precisión, exhaustividad, abstención, deriva y fugas. La evaluación
de recuperación no se presenta como evaluación de un modelo generativo.

\chapter{Evidencia y resultados responsables}

\section{Evidencia disponible}

Los tres perfiles exhiben objetos canónicos, capas de acceso documental, servicios,
catálogos y representaciones derivadas. La recurrencia apoya la relevancia práctica del
problema, no la novedad de cada mecanismo. La contribución de COBRIA está en convertir
esa recurrencia en contratos comparables: autoridad, ámbito, revisión, tiempo, versión,
equivalencia y reconstrucción.

Las pruebas disponibles cubren reglas de acceso, matrices de roles, aislamiento, ciclos
de vida, idempotencia y reconstrucción. También existe evidencia longitudinal de menor
acceso directo desde presentación hacia la capa documental. Estos resultados prueban
rutas concretas, pero no productividad humana ni superioridad general.

\section{Resultados COBRIA Core 1.0}

La primera réplica congelada se ejecutó sobre dos motores documentales embebidos,
PouchDB y LokiJS, con escalas de 100, 1,000 y 5,000 documentos y treinta corridas por
celda. En total se conservaron 180 corridas. Todas las celdas superaron conjuntamente
P1, P2, R1, S1 y A1: equivalencia lógica completa, reconstrucción repetible,
amplificación de dos escrituras por cambio proyectado, rechazo de ámbito ausente y cero
documentos cruzados.

En PouchDB las medianas de lectura canónica fueron 1,186; 13,589 y 69,291 ms, frente a
0,440; 4,719 y 23,131 ms para la proyección. Las reconstrucciones medianas requirieron
8,376; 89,301 y 500,070 ms. En LokiJS las lecturas canónicas fueron 0,145; 1,263 y
7,026 ms, frente a 0,052; 0,409 y 2,220 ms; R1 requirió 0,663; 6,452 y 46,011 ms.

Estos resultados apoyan conformidad funcional y muestran una reducción de lectura en
este arnés embebido. No representan redes, regiones, concurrencia administrada, cuotas
ni facturación. HN1--HN4 quedan apoyadas para esta implementación de referencia; HN5
recibe apoyo local condicionado y permanece abierta para servicios administrados.

\chapter{Amenazas a la validez y síntesis}

Los casos comparten ecosistema y autoría. Los emuladores no reproducen cuotas,
particiones, latencia regional ni facturación administrada. Los datos anonimizados y
sintéticos pueden subrepresentar deriva y errores históricos. La inspección no sustituye
ejecución; las pruebas automáticas no sustituyen revisión humana; y una implementación
correcta no demuestra comprensión por todos los equipos.

La evidencia permite afirmar que COBRIA Core 1.0 es formalizable, ejecutable y conforme
en dos motores NoSQL documentales embebidos. No permite afirmar una mejora universal de
rendimiento ni un umbral económico transferible. El siguiente nivel de evidencia debe
provenir de servicios documentales administrados y una réplica externa.


\part{Síntesis, adopción y conclusiones}
\chapter{Síntesis integral de COBRIA}

COBRIA organiza decisiones que aparecen cuando un sistema NoSQL conserva documentos
operacionales y necesita búsquedas, resúmenes, conjuntos analíticos y herramientas de
inteligencia artificial. Existe una fuente canónica por hecho de dominio y toda
representación derivada declara cómo se obtiene, verifica, delimita y reconstruye.

\section{Respuesta a la investigación}

La estructura Objeto--Dater--Servicio--Catálogo no constituye por sí sola una teoría
nueva. COBRIA aporta una composición explícita para documentos NoSQL: identidad y
revisión en el canónico; transformación y procedencia en el Dater; reglas y autorización
en servicios; semántica versionada en catálogos; y reconstrucción en cada proyección.

La reducción de lectura depende de consulta, índice, forma documental y distribución.
Una proyección se adopta después de compararla con la consulta canónica indexada. La
amplificación se registra mediante operaciones, bytes, eventos, reintentos,
almacenamiento y trabajo operacional.

La adaptabilidad hacia minería de datos e inteligencia artificial procede del linaje:
conjuntos, características, fragmentos e inferencias conservan fuente, tiempo, versión y
ámbito. Esto permite reconstruir un corte, detectar deriva y evitar que recuperación o
agentes crucen fronteras de autorización.

\chapter{Guía de adopción}

Considérese COBRIA cuando existen múltiples consumidores, documentos heterogéneos,
aislamiento organizacional, evolución frecuente, derivados costosos de reparar o
productos de datos reproducibles. Simplifíquese cuando una colección indexada satisface
el objetivo o el costo de gobierno supera el riesgo.

\begin{enumerate}[leftmargin=*]
  \item identificar la escritura autoritativa y cerrar rutas paralelas;
  \item hacer obligatorio el ámbito en acceso, caché, eventos y recuperación;
  \item añadir revisión, tiempo y versión de esquema;
  \item medir la consulta antes de crear una proyección;
  \item declarar transformación, equivalencia y reconstrucción;
  \item publicar por versión y practicar reparación;
  \item extender procedencia hacia analítica e inteligencia artificial;
  \item retirar derivados sin beneficio o responsable.
\end{enumerate}

\chapter{Contribuciones y límites}

La contribución conceptual es el dato reconstruible como contrato. La formal es la tupla
que une fuente, transformación, ámbito, versión, equivalencia y reparación. La
arquitectónica es una frontera coherente desde documentos hasta productos de IA. La
práctica es un catálogo, listas de comprobación y protocolo reproducible.

La contribución empírica conserva dos estratos. El primero contiene 180 corridas locales
históricas sobre PouchDB y LokiJS. El segundo contiene 270 corridas distribuidas
históricas sobre MongoDB 8.0, CouchDB 3.4 y OpenSearch 2.19.3. Todas superaron P1, P2,
R1, S1 y A1 bajo sus oráculos originales y no registraron fugas de ámbito. Como una
auditoría posterior fortaleció el arnés, estas cifras describen evidencia histórica y no
certificación vigente. La repetición preliminar corregida completó casos pequeños en
MongoDB y CouchDB; OpenSearch 3.2.0 quedó pendiente por incompatibilidad de plataforma.

Falta repetir las celdas distribuidas completas con el arnés corregido; ejecutar
OpenSearch 3.2.0 en ARM nativo o infraestructura compatible; realizar una réplica
independiente; completar el doble cribado bibliográfico; y ejecutar el estudio humano de
transferencia. P1, P2, R1, S1 y A1 forman el conjunto de criterios de conformidad, pero
la certificación independiente permanece pendiente.

\chapter{Conclusión}

COBRIA hace visibles seis preguntas: qué documento manda, a qué ámbito pertenece, qué
versión representa, qué copia puede perderse, cómo se reconstruye y cuánto cuesta
conservarla. Al extenderlas a conjuntos de datos, recuperación e inferencias, ofrece un
puente entre ingeniería NoSQL, minería de datos e inteligencia artificial.

La propuesta es suficientemente concreta para implementarse y refutarse. La réplica
embebida apoya sus contratos nucleares, mientras la evaluación administrada y humana
sigue abierta. Esa frontera define con honestidad qué se sabe y qué debe venir después.


\part{Manual de réplica y operación NoSQL}
\chapter{Propósito del manual de réplica}

Esta parte traduce las hipótesis de COBRIA a un procedimiento ejecutable en motores
NoSQL documentales. Su función no es imponer un proveedor, sino impedir que cada réplica
redefina silenciosamente la unidad de datos, el índice, el ámbito o el criterio de
equivalencia. Una comparación solo es interpretable cuando las configuraciones responden
la misma pregunta lógica y registran las diferencias físicas relevantes.

\section{Resultado esperado}

Al terminar una réplica, un tercero debe poder reconstruir cuatro productos: el conjunto
de documentos de entrada, las configuraciones comparadas, el registro de cada corrida y
el informe que enlaza cifras con observaciones originales. También debe poder comprobar
qué se excluyó, por qué se excluyó y si una enmienda ocurrió antes o después de conocer
el resultado afectado.

\section{Principios de neutralidad}

\begin{enumerate}[leftmargin=*]
  \item la semántica del dominio permanece igual entre motores;
  \item cada adaptador usa capacidades documentales normales del motor evaluado;
  \item los índices se declaran y congelan antes de medir;
  \item las proyecciones contienen únicamente los campos necesarios para la consulta;
  \item costo, frescura y reparación se informan junto con la latencia;
  \item una configuración que pierde permanece en el conjunto publicado.
\end{enumerate}

\chapter{Selección y descripción de motores documentales}

\section{Criterios de inclusión}

Un motor es elegible cuando almacena documentos u objetos agregados, permite claves y
filtros con ámbito, ofrece índices declarables, admite operaciones de actualización y
posee una vía reproducible local o administrada. La réplica debe incluir al menos dos
motores con modelos operacionales distintos para evitar que una capacidad particular se
confunda con una propiedad general de COBRIA.

\section{Ficha del motor}

Cada adaptador publica versión, modalidad de despliegue, región, consistencia observable,
límites de documento, estrategia de partición, índices, paginación, transacciones
documentales disponibles, flujos de cambios, cuotas y unidad de facturación. Si una
capacidad no existe, el protocolo registra la adaptación en vez de simular que todos los
motores son idénticos.

\begin{table}[H]\centering
\caption{Campos mínimos de la ficha del motor NoSQL.}
\begin{tabularx}{\textwidth}{p{4cm}X}\toprule
Campo & Pregunta de control\\\midrule
Identidad de versión & ¿Qué versión exacta o fecha del servicio se evaluó?\\
Distribución & ¿Cómo se asignan documentos a particiones o regiones?\\
Consistencia & ¿Qué garantía observa la consulta medida?\\
Índices & ¿Qué campos, orden y ámbito cubre cada índice?\\
Costo & ¿Qué operaciones, bytes y almacenamiento se facturan?\\
Recuperación & ¿Cómo se exportan, reconstruyen y verifican documentos?\\
\bottomrule\end{tabularx}\end{table}

\chapter{Modelo documental neutral}

\section{Documento canónico}

El documento de referencia contiene identidad estable, ámbito, revisión, versión de
esquema, tiempo de validez, tiempo de registro, estado y carga variable. Los campos
opcionales se distribuyen según tasas congeladas. Las claves no incluyen información
personal y se derivan de la semilla del generador.

\section{Proyección de cobertura}

La proyección conserva identidad, ámbito, revisión, campos de filtro, orden y salida de
la consulta objetivo. No contiene reglas de negocio ni acepta escrituras autoritativas.
Cada elemento enlaza documento fuente y versión de transformación. Una variante más
ancha se permite como análisis de sensibilidad, pero se publica como configuración
distinta.

\section{Catálogos y evolución}

Los datos incluyen versiones antiguas, valores desconocidos y sustituciones de catálogo.
La transformación debe migrar o aislar cada caso según una política declarada. El
generador conserva la verdad esperada para que la limpieza pueda evaluarse sin inspección
manual de todas las filas lógicas.

\chapter{Generación de cargas semirrealistas}

\section{Perfiles}

El perfil N1 representa eventos y multimedia; N2, reservas y disponibilidad; N3,
inventario, ventas y cierres. Cada perfil modifica selectividad, tamaño, frecuencia de
actualización, eventos tardíos y campos opcionales. Las diferencias deben afectar la
forma de la carga sin revelar nombres ni reglas de los sistemas que la inspiraron.

\section{Anomalías controladas}

\begin{table}[H]\centering
\caption{Anomalías que debe producir el generador.}
\begin{tabularx}{\textwidth}{p{3.4cm}p{3cm}X}\toprule
Anomalía & Parámetro & Resultado esperado\\\midrule
Campo ausente & tasa por perfil & valor predeterminado o exclusión declarada\\
Documento antiguo & versión de esquema & migración determinista\\
Evento tardío & retraso máximo & corte temporal reproducible\\
Entrega duplicada & identidad lógica & un solo efecto\\
Catálogo desconocido & tasa de novedad & cuarentena observable\\
Revisión obsoleta & distancia de revisión & rechazo sin regresión\\
\bottomrule\end{tabularx}\end{table}

La semilla gobierna identidades, orden y anomalías. Cambiarla produce otra muestra, no
otra definición del experimento. Toda tasa se registra en el manifiesto y se incluye en
el identificador de configuración.

\chapter{Protocolo P1: lectura documental}

\section{Configuraciones}

N0 consulta documentos canónicos usando el índice nativo apropiado. N1 consulta una
proyección de cobertura reconstruible. N2, opcional, evalúa una proyección más ancha para
estimar el costo de duplicar campos. No se utiliza una consulta deliberadamente sin
índice como adversario principal, porque la pregunta científica es si la proyección
mejora una implementación NoSQL razonable.

\section{Ejecución}

Cada corrida prepara un conjunto nuevo o restaura una instantánea congelada, verifica
conteos lógicos, calienta según protocolo y rota el orden de configuraciones. Se registran
latencias por consulta, documentos examinados cuando el motor lo permite, bytes de
respuesta, operaciones facturables y marca de frescura. La corrida se resume mediante
p50, p95 y p99.

\section{Interpretación}

Una diferencia pequeña de latencia puede ser irrelevante si la proyección duplica datos
sensibles, aumenta costo o incumple frescura. El informe presenta un vector de efectos.
La recomendación se formula por región de carga y nunca como superioridad universal.

\chapter{Protocolo P2: escritura y convergencia}

\section{Camino de actualización}

P2 modifica el documento canónico con revisión esperada. El mecanismo produce o detecta
el cambio, actualiza el derivado de forma idempotente y publica la nueva marca de
frescura. Se miden tiempo de confirmación canónica y tiempo de convergencia por separado.

\section{Amplificación}

La amplificación incluye documentos escritos, eventos, reintentos, bytes, invocaciones,
lecturas auxiliares, almacenamiento y tiempo operacional. Una implementación que oculta
trabajo en una función o canal de cambios no posee amplificación cero: debe atribuir ese
costo a la configuración.

\section{Frontera de decisión}

Sea $\Delta Q$ el ahorro por lectura y $\Delta W$ el costo adicional por escritura. La
frontera temporal simple exige $Q\Delta Q>W\Delta W$. El informe añade costo monetario,
almacenamiento, frescura y riesgo como dimensiones separadas. Si $\Delta Q\leq0$, la
proyección no se justifica por esa consulta.

\chapter{Protocolos de reconstrucción y publicación}

\section{Reconstrucción total}

R1 crea una versión candidata desde documentos canónicos, catálogos y transformación
congelada. La candidata no es visible para consumidores hasta superar conteo, hash por
partición, muestreo de contenido e invariantes de ámbito. La publicación cambia un alias
o registro de versión y conserva la anterior durante una ventana reversible.

\section{Reconstrucción incremental}

La ruta incremental consume cambios posteriores a un punto de control. Cada lote registra
inicio, fin, documentos procesados, duplicados y revisión máxima. Una reconstrucción
total periódica actúa como verificador independiente. Si ambas rutas divergen, la
incremental se considera defectuosa aunque la consulta parezca correcta.

\section{Interrupción segura}

El arnés interrumpe preparación, escritura parcial, verificación y cambio de versión. En
ningún caso una candidata incompleta debe reemplazar la versión activa. La reanudación
puede continuar o reiniciar, pero siempre produce el mismo resultado lógico.

\chapter{Aislamiento, autorización y recuperación}

\section{Matriz adversarial}

La prueba combina identidad válida, identidad ausente, ámbito manipulado, rol insuficiente,
clave de caché incompleta, tarea administrativa y recuperación semántica. Los candidatos
se restringen antes de ordenar o construir contexto. Un resultado ocultado después de
recuperarlo ya constituye una violación del contrato.

\section{Pruebas negativas}

Cada ruta positiva exige una negativa equivalente. Se verifican identificadores de otro
ámbito, referencias cruzadas, eventos reinyectados, paginación y búsqueda por campos
opcionales. Los registros de auditoría evitan contenido sensible, pero conservan motivo,
política, actor, ámbito y correlación.

\chapter{Productos para minería de datos e inteligencia artificial}

\section{Conjunto reproducible}

El conjunto declara población, corte temporal, fuente, versiones, limpieza, catálogos,
exclusiones, semilla y división. Su identificador cambia cuando una dependencia cambia.
La réplica reconstruye el conjunto desde cero y compara manifiesto, conteos, distribución
y huellas de contenido.

\section{Recuperación con evidencia}

Cada fragmento enlaza documento, revisión, ámbito, transformación y autorización. Se
miden precisión, exhaustividad, abstención, fugas y citas válidas. Las instrucciones
contenidas en documentos se tratan como datos, no como órdenes. La calidad lingüística
de una respuesta requiere un estudio separado.

\section{Herramientas de agentes}

Una herramienta declara capacidad, esquema de entrada, validación, idempotencia, efecto,
confirmación y auditoría. Las pruebas incluyen repetición, entrada parcial, cambio de
ámbito y cancelación. Una acción de alto riesgo sin confirmación falla el protocolo aun
si el resultado funcional parece correcto.

\chapter{Observabilidad y costos}

\section{Indicadores}

El tablero mínimo contiene latencia, error, frescura, atraso, duplicados, rechazos de
revisión, reconstrucciones, divergencia, operaciones por solicitud, almacenamiento y
costo por ámbito. Los indicadores se segmentan por versión de proyección para detectar
mezclas durante migración.

\section{Presupuesto de proyecciones}

Cada derivado consume un presupuesto compuesto por almacenamiento, escritura, eventos,
monitoreo y carga cognitiva. El responsable revisa uso y beneficio en una fecha fija.
Una proyección sin consultas, sin verificador o con reconstrucción rota entra en retiro.

\section{Criterio de retiro}

Retirar significa detener consumidores, verificar alternativas, eliminar el productor,
esperar la ventana de seguridad y suprimir datos derivados. La fuente canónica permanece
intacta. El proceso se prueba como cualquier migración y conserva un registro de decisión.

\chapter{Análisis estadístico y resultados negativos}

\section{Unidad inferencial}

La unidad es la corrida completa. Las consultas internas estiman un percentil de esa
corrida y no se presentan como observaciones independientes. Las comparaciones usan
pares generados con el mismo conjunto y semilla. Se informan mediana, intervalo
remuestreado, diferencia, razón y dirección de cada par.

\section{Multiplicidad y sensibilidad}

Las familias de hipótesis se congelan antes de medir y se ajustan por comparaciones
múltiples. El análisis de sensibilidad cambia una dimensión por vez: escala, tamaño de
documento, selectividad, consistencia o ancho de proyección. Las exploraciones se marcan
como tales y no reemplazan el análisis confirmatorio.

\section{Publicación de derrotas}

Una celda donde el documento canónico indexado iguala o supera a la proyección es un
resultado útil. Delimita dónde COBRIA debe simplificarse. El paquete conserva corridas
válidas aunque contradigan HN5 y explica cualquier descarte mediante una regla previa.

\chapter{Paquete abierto y revisión independiente}

\section{Contenido del depósito}

El depósito incluye protocolo, enmiendas, generador, adaptadores, manifiestos, eventos
crudos, análisis, tablas, figuras, entorno, licencias y declaración de disponibilidad.
Un comando de verificación comprueba huellas y vuelve a producir las salidas publicadas.

\section{Revisión técnica}

La persona revisora recibe una copia anonimizada, rúbrica y matriz de trazabilidad. Debe
examinar equivalencia de configuraciones, unidad inferencial, exclusiones, afirmaciones y
amenazas. Sus observaciones se responden una por una sin borrar desacuerdos sustantivos.

\section{Criterio de cierre}

La fase cuantitativa se considera completa cuando dos motores NoSQL terminan todas las
celdas congeladas, el paquete reproduce cifras y una revisión independiente no detecta
una falla que invalide la inferencia principal. Antes de ese punto, los resultados se
describen como piloto o evidencia parcial.

\chapter{Cuadernos de campo reproducibles}

\section{Bitácora de preparación}

Antes de ejecutar una celda, la bitácora registra identificador, fecha, responsable,
motor NoSQL, versión, configuración efectiva, semilla, conjunto, estado de índices y
huellas de los artefactos. El registro no depende de memoria personal: se produce desde
el mismo entorno que ejecuta la prueba y se conserva junto al resultado crudo. Si una
configuración cambia, nace una celda nueva; nunca se sobrescribe la evidencia anterior.

La preparación verifica además espacio disponible, sincronización temporal, ausencia de
carga ajena y aislamiento de credenciales. Una lectura breve de calentamiento confirma
que el adaptador puede resolver documento, ámbito y revisión. Este control evita gastar
una corrida completa para descubrir un error elemental de montaje.

\section{Bitácora de ejecución}

Cada fase deja marcas de inicio y final, conteos aceptados y rechazados, memoria máxima,
operaciones por segundo, percentiles de latencia y códigos de error. Los mensajes libres
se complementan con campos estructurados para permitir análisis automático. El operador
puede anotar un incidente, pero no clasificarlo retroactivamente sin conservar la razón y
la regla de exclusión utilizada.

\section{Bitácora de cierre}

Al terminar se verifica que los conteos coincidan con el manifiesto, que no existan
ámbitos cruzados, que las colas estén drenadas y que las huellas puedan recalcularse.
Una corrida válida se sella como inmutable. Una corrida inválida también se conserva,
etiquetada con su causa, porque la tasa y naturaleza de los fallos forman parte de la
evaluación de operabilidad.

\chapter{Migraciones y evolución del contrato}

\section{Cambio aditivo}

El cambio preferido agrega un campo opcional con semántica, valor ausente explícito y
versión de contrato. Lectores nuevos aceptan documentos antiguos; lectores antiguos
ignoran el campo nuevo cuando esa conducta es segura. La prueba mínima cubre ambas
direcciones y demuestra que la ausencia no se confunde con cero, cadena vacía o falso.

\section{Cambio incompatible}

Renombrar, dividir o reinterpretar un campo exige lectura doble, escritura controlada y
una ventana de convivencia. La migración se expresa como función determinista e
idempotente. Antes de retirar la versión previa se reconstruyen proyecciones, se comparan
resultados por ámbito y se comprueba que ningún consumidor siga solicitando el contrato
anterior.

\section{Reversión}

Toda migración declara el punto hasta el cual puede revertirse sin pérdida. Si el cambio
destruye información, se conserva una copia temporal gobernada y con caducidad. Revertir
no significa restaurar archivos a ciegas: significa volver a una combinación coherente
de contrato, documentos, catálogos, índices, proyecciones y consumidores.

\chapter{Rúbricas de conformidad COBRIA}

\section{Rúbrica del documento canónico}

La conformidad se valora de cero a dos por criterio: identidad estable, ámbito explícito,
revisión monotónica, procedencia, contrato versionado y validación en frontera. Cero
indica ausencia; uno, presencia parcial o no verificable; dos, evidencia automática. Un
promedio alto no compensa un cero en aislamiento o autoridad, que son condiciones de
seguridad y no preferencias de estilo.

\section{Rúbrica del derivado}

Un derivado demuestra propietario, propósito de consulta, versión de proyección,
idempotencia, frescura observable, reconstrucción total y procedimiento de retiro. Se
evalúa también si puede confundirse con la fuente autorizada. La etiqueta ``caché'' no
exime ningún criterio: si persiste y alimenta decisiones, debe poder explicarse y rehacerse.

\section{Rúbrica del uso analítico}

El conjunto analítico enlaza cada fila o fragmento con documento, revisión, ámbito y
transformación. La partición de entrenamiento evita contaminación temporal y cruce de
entidades. Una salida de inteligencia artificial conserva citas comprobables y se
abstiene cuando no reúne evidencia autorizada suficiente. La métrica de calidad se
publica junto con la de fuga, no en sustitución de ella.

\chapter{Escenarios de aceptación integral}

\section{Escenario A: lectura y actualización concurrentes}

Se inicia una lectura sobre una revisión conocida mientras llega una actualización del
mismo documento. El sistema puede devolver la revisión anterior o la nueva según el
contrato declarado, pero nunca una mezcla de campos. El evento posterior incluye la
nueva revisión, la proyección converge y una repetición no duplica efectos.

\section{Escenario B: reconstrucción bajo cambio de esquema}

Se elimina un derivado, se activa una nueva versión de proyección y se recorre la fuente
canónica. El resultado se compara con una construcción incremental limpia usando huellas,
conteos y consultas de aceptación. Toda diferencia queda localizada por documento y
transformación; una coincidencia meramente global no basta.

\section{Escenario C: recuperación con inteligencia artificial}

Una consulta solicita información disponible en dos ámbitos, pero la identidad solo
autoriza uno. El recuperador filtra antes de clasificar, conserva procedencia y entrega
únicamente fragmentos permitidos. Si esos fragmentos no sustentan una respuesta, el
sistema se abstiene. La prueba falla ante una cita inexistente, una revisión obsoleta no
declarada o cualquier exposición del ámbito prohibido.

\section{Escenario D: degradación y recuperación}

Se interrumpe temporalmente el consumidor que mantiene una proyección. La fuente acepta
escrituras dentro de su contrato, el atraso se vuelve visible y las lecturas degradadas
se etiquetan. Al reanudar, el consumidor procesa desde el último punto confirmado,
converge sin duplicados y demuestra equivalencia con una reconstrucción independiente.


\IfFileExists{partes/parte-viii-anexos-reproducibilidad.tex}{%
  \part{Anexos técnicos y reproducibilidad}
  \chapter{Glosario operacional}

Este glosario fija el uso de términos dentro del manuscrito. No pretende reemplazar
definiciones generales de ingeniería de software.

\begin{longtable}{>{\bfseries}p{3.6cm}p{10.6cm}}
\caption{Glosario COBRIA.}\label{tab:glosario-cobria}\\
\toprule
Término & Definición operacional \\
\midrule
\endfirsthead
\toprule
Término & Definición operacional \\
\midrule
\endhead
Autoridad & Capacidad reconocida para establecer un hecho de negocio.\\
Canónico & Representación autoritativa de una entidad dentro de un ámbito y versión.\\
Catálogo & Vocabulario o conjunto de referencias permitidas y gobernadas.\\
Conector & Adaptador que encapsula la API de almacenamiento o proveedor.\\
Contrato & Especificación versionada de entradas, salidas, políticas y verificaciones.\\
Dater & Nombre local de una clase que combina acceso, mapeo, Repositorio, caché o índices.\\
Conjunto de datos & Proyección analítica versionada con selección, transformación y manifiesto.\\
Derivado & Artefacto calculado desde una fuente que no adquiere autoridad automáticamente.\\
Entidad & Objeto con identidad, ámbito, atributos, ciclo de vida, versión y revisión.\\
Registro de evidencias & Registro minimizado que enlaza solicitud, evidencia, decisión y efecto.\\
Conjunto de características & Conjunto versionado de características con tiempo y procedencia.\\
Idempotencia & Propiedad por la cual repetir una operación no multiplica su efecto.\\
Inferencia & Salida probabilística de un modelo, distinta de un hecho confirmado.\\
Linaje & Grafo que conecta artefactos con fuentes, transformaciones y versiones.\\
Mapeador & Componente que traduce entre dominio y representación física.\\
Metadato & Dato que declara estructura, relación, política, versión o significado.\\
Cola de salida & Registro persistente de una intención de publicación asociada a un cambio.\\
Proyección & Representación derivada optimizada para una lectura, proceso o análisis.\\
Gestor de proyecciones & Componente que mantiene, verifica, reconstruye y versiona proyecciones.\\
Reconstruibilidad & Capacidad de regenerar un derivado desde fuentes y contratos suficientes.\\
Repositorio & Interfaz de acceso a entidades expresada en lenguaje del dominio.\\
Revisión & Contador o token utilizado para detectar actualización concurrente obsoleta.\\
Ámbito & Frontera de pertenencia y autorización, como organización, organización o contexto.\\
Servicio & Componente que expresa coordinación o capacidad de dominio/aplicación.\\
Caso de uso & Frontera autorizada que coordina una intención y sus efectos.\\
Marca de corte & Punto temporal o cursor declarado para procesar cambios o cortes.\\
\bottomrule
\end{longtable}

\chapter{Plantillas de contratos}
\begingroup
\small

\section{Contrato de entidad}

\begin{verbatim}
entity:
  name: ExampleEntity
  version: 1
  identity: id
  ámbito: [organizaciónId]
  revision: revision
  lifecycle: [draft, active, archived]
  fields:
    - name: title
      type: string
      required: true
      sensitivity: internal
  invariants:
    - "title must not be empty"
  retention: policy-entity-default
  responsable: domain-team
\end{verbatim}

\section{Contrato de Mapeador}

\begin{verbatim}
mapper:
  entity: ExampleEntity
  version: 2
  readsPhysicalVersions: [1, 2]
  writesPhysicalVersion: 2
  fields:
    id: key
    organizaciónId: organización_id
    title: display_name
    revision: rev
  roundTrip:
    equality: semantic
    ignored: [storage_updated_at]
\end{verbatim}

\section{Contrato de proyección}

\begin{verbatim}
projection:
  id: example-by-status
  version: 3
  sources: [ExampleEntity@1]
  key: [organizaciónId, status, id]
  consistency: eventual
  freshnessSloSeconds: 30
  deterministic: true
  idempotent: true
  authoritative: false
  rebuild:
    mode: [full, incremental]
    batchSize: 500
    checkpoint: required
    cambio de versión: version-alias
  verify:
    coverage: 1.0
    orphanCount: 0
  responsable: query-team
\end{verbatim}

\section{Contrato de conjunto de datos}

\begin{verbatim}
conjunto de datos:
  id: example-training-set
  version: 4
  purpose: approved-task
  ámbito: organización
  sources: [ExampleEntity@1, ExampleEvent@2]
  cutoff: required
  transformationsCommit: abc123
  esquema: example-conjunto de datos-esquema@4
  quality:
    completenessMin: 0.98
    duplicateMax: 0
  privacy:
    allowedFields: [category, eventTime, value]
    retentionDays: 90
  lineage: required
  responsable: equipo-analitica
\end{verbatim}

\section{Contrato de característica}

\begin{verbatim}
feature:
  name: rolling_count_30d
  version: 2
  entity: ExampleEntity
  type: integer
  window: 30d
  availabilityTime: ingestionTime
  toleranciaFueraEnLinea: 0
  freshnessSeconds: 300
  sourceDataset: example-events@2
  responsable: ml-platform
\end{verbatim}

\section{Contrato de herramienta de agente}

\begin{verbatim}
tool:
  name: propose_entity_change
  version: 1
  capability: entity.change.propose
  inputSchema: proposal-input@1
  outputSchema: proposal-output@1
  ámbito: inherited-and-verified
  risk: medium
  confirmation: required-before-effect
  idempotencyKey: required
  reversible: true
  compensation: cancel-proposal
  audit: evidence-ledger
\end{verbatim}

\endgroup
\chapter{Listas de comprobación de ingeniería y gobierno}

\section{Lista de comprobación de entidad canónica}

\begin{longtable}{p{1cm}p{12.8cm}}
\caption{Verificación de una entidad canónica.}\label{tab:check-entidad}\\
\toprule
OK & Criterio \\
\midrule
\endfirsthead
\toprule
OK & Criterio \\
\midrule
\endhead
$\square$ & Posee identidad estable y definida.\\
$\square$ & El ámbito es obligatorio y forma parte de toda clave de acceso.\\
$\square$ & Distingue atributos, relaciones y derivados.\\
$\square$ & Declara versión de esquema y estrategia de lectura histórica.\\
$\square$ & Usa revisión o precondición para conflictos.\\
$\square$ & Sus invariantes se ejecutan dentro de casos de uso.\\
$\square$ & No conoce rutas ni SDK del proveedor.\\
$\square$ & Tiene responsable, clasificación y retención.\\
$\square$ & Su Mapeador pasa ida y vuelta semántico.\\
$\square$ & Existe una sola autoridad para cada hecho.\\
\bottomrule
\end{longtable}

\section{Lista de comprobación de proyección}

\begin{longtable}{p{1cm}p{12.8cm}}
\caption{Verificación de una proyección reconstruible.}\label{tab:check-proyeccion}\\
\toprule
OK & Criterio \\
\midrule
\endfirsthead
\toprule
OK & Criterio \\
\midrule
\endhead
$\square$ & Declara problema de lectura y consumidores.\\
$\square$ & Identifica fuentes y versiones suficientes.\\
$\square$ & No acepta escrituras autoritativas.\\
$\square$ & Define ámbito, clave, función y versión.\\
$\square$ & Declara consistencia y frescura.\\
$\square$ & Es idempotente o documenta por qué no puede serlo.\\
$\square$ & Posee rebuild total; incremental si está justificado.\\
$\square$ & Verifica cobertura, huérfanos y diferencias.\\
$\square$ & Tiene presupuesto, responsable, observabilidad y retiro.\\
$\square$ & Puede hacer cambio de versión o reversión sin mezclar versiones.\\
\bottomrule
\end{longtable}

\section{Lista de comprobación multiinquilino}

\begin{longtable}{p{1cm}p{12.8cm}}
\caption{Verificación de aislamiento por ámbito.}\label{tab:check-organización}\\
\toprule
OK & Criterio \\
\midrule
\endfirsthead
\toprule
OK & Criterio \\
\midrule
\endhead
$\square$ & API, caso de uso y Repositorio requieren ámbito.\\
$\square$ & Rutas e índices incluyen ámbito antes de identidad.\\
$\square$ & Claves de caché y de idempotencia incluyen ámbito.\\
$\square$ & Procesos de trabajo vuelven a autorizar mensajes.\\
$\square$ & Exports, backups y conjuntos de datos conservan separación.\\
$\square$ & Reconstrucciones requieren alcance explícito.\\
$\square$ & Telemetría no expone cargas útiles de otros organizaciones.\\
$\square$ & IDs iguales entre organizaciones forman parte de las pruebas.\\
$\square$ & Las pruebas negativas abarcan lectura y escritura.\\
$\square$ & Una fuga se trata como fallo crítico, no como tasa promedio.\\
\bottomrule
\end{longtable}

\section{Lista de comprobación de IA y agentes}

\begin{longtable}{p{1cm}p{12.8cm}}
\caption{Verificación de componentes inteligentes.}\label{tab:check-ia}\\
\toprule
OK & Criterio \\
\midrule
\endfirsthead
\toprule
OK & Criterio \\
\midrule
\endhead
$\square$ & Conjunto de datos y Conjunto de características poseen propósito, versión y linaje.\\
$\square$ & El corte temporal evita utilizar información futura.\\
$\square$ & Modelo y instrucción están versionados.\\
$\square$ & Recuperación filtra ámbito y permisos antes de generar.\\
$\square$ & La evidencia enlaza fragmento y fuente vigente.\\
$\square$ & El contenido recuperado no puede redefinir políticas.\\
$\square$ & Herramientas tienen capacidades mínimas y schemas.\\
$\square$ & Acciones sensibles solicitan confirmación.\\
$\square$ & Las operaciones mutantes son idempotentes.\\
$\square$ & Una inferencia no se persiste como hecho sin caso de uso.\\
$\square$ & Existen monitoreo, retiro y respuesta a incidentes.\\
\bottomrule
\end{longtable}

\chapter{Diccionario de métricas}

\begin{longtable}{>{\bfseries}p{3.2cm}p{5.2cm}p{5.2cm}}
\caption{Definición de métricas del estudio.}\label{tab:diccionario-metricas}\\
\toprule
Métrica & Cálculo o unidad & Precaución \\
\midrule
\endfirsthead
\toprule
Métrica & Cálculo o unidad & Precaución \\
\midrule
\endhead
Latencia p50 & Mediana en ms por corrida & No sustituir por media.\\
Latencia p95/p99 & Percentil por corrida & Requiere muestra suficiente.\\
Rendimiento efectivo & Operaciones exitosas por segundo & Reportar errores y concurrencia.\\
$A_r$ & Bytes leídos / bytes útiles & Definir respuesta vacía.\\
$A_w$ & Escrituras físicas / escritura lógica & Separar réplicas internas.\\
$A_s$ & Almacenamiento total / canónico & Incluir índices y staging.\\
Frescura & Confirmación a visibilidad correcta & Usar reloj y frontera declarados.\\
Cobertura & Esperados presentes / esperados & No detecta registros extra.\\
Huérfanos & Derivados sin fuente válida & Reportar conteo absoluto.\\
Exactitud de reconstrucción & $1-(faltantes+extras+diferentes)/esperados$ & Desglosar errores.\\
Cobertura de linaje & Artefactos con ruta completa / evaluados & No implica calidad.\\
Deriva & Distancia o prueba por variable & Cambio no siempre es defecto.\\
Precisión@k & Relevantes recuperados / $k$ & Depende del juicio de relevancia.\\
Exhaustividad@k & Relevantes recuperados / relevantes & Requiere conjunto completo.\\
Abstención & Sin respuesta cuando no hay evidencia & Evaluar falsos rechazos.\\
Violaciones ámbito & Accesos o efectos cruzados & Objetivo obligatorio: cero.\\
Esfuerzo & Tiempo activo y artefactos modificados & Separar aprendizaje.\\
Tamaño de efecto & Diferencia estandarizada o razón & Informar intervalo.\\
\bottomrule
\end{longtable}

\chapter{Esquema del paquete reproducible}

\section{Estructura de directorios}

\begin{verbatim}
cobria-replication/
  README.md
  LICENSE
  CITATION.cff
  environment/
    versions.bloqueo
    hardware-template.json
  reference/
    src/
    tests/
  contracts/
    entities/
    projections/
    conjuntos de datos/
    tools/
  generators/
    profile-r/
    profile-a/
    profile-b/
  experiments/
    performance/
    reconstruction/
    isolation/
    evolution/
    intelligence/
  analysis/
    notebooks/
    scripts/
  results/
    crudo/
    processed/
    figures/
    manifests/
\end{verbatim}

\section{Manifiesto de corrida}

\begin{verbatim}
{
  "runId": "uuid",
  "startedAt": "ISO-8601",
  "revisión": "git-sha",
  "profile": "R|A|B",
  "configuration": "B0|B1|C1|C2|C3|C4",
  "seed": 42001,
  "N": 100000,
  "organizaciones": 10,
  "workload": "W2",
  "concurrency": 32,
  "projectionCount": 3,
  "environmentHash": "sha256",
  "conjunto de datosHash": "sha256",
  "result": "completed|failed|excluded",
  "exclusionReason": null
}
\end{verbatim}

\section{Formato de evento de medición}

\begin{verbatim}
{
  "runId": "uuid",
  "operationId": "uuid",
  "operation": "query-by-status",
  "organizaciónAlias": "t-03",
  "startedNs": 0,
  "durationNs": 0,
  "reads": 0,
  "writes": 0,
  "bytesRead": 0,
  "bytesWritten": 0,
  "retries": 0,
  "resultCount": 0,
  "status": "ok|business-reject|error|timeout"
}
\end{verbatim}

\section{Reglas de publicación}

Los resultados crudo son inmutables. El procesamiento escribe otra carpeta y conserva
versión del guion. Las tablas del manuscrito se generan, no se editan manualmente. Una
exclusión permanece en crudo y aparece en el reporte de flujo. Los secretos se suministran
fuera del paquete y nunca se imprimen.

\section{Lista de comprobación de réplica}

\begin{enumerate}[leftmargin=*]
  \item verificar hash y licencia del paquete;
  \item instalar versiones fijadas;
  \item registrar hardware y límites;
  \item generar perfiles con semillas publicadas;
  \item ejecutar pruebas de corrección antes del prueba comparativa;
  \item correr configuraciones en orden aleatorio;
  \item conservar registros y manifiestos;
  \item ejecutar análisis sin modificar crudo;
  \item comparar tablas y documentar divergencias;
  \item publicar entorno y resultados negativos.
\end{enumerate}

\chapter{Plantillas de registro y publicación}

\section{Ficha de experimento}

\begin{table}[H]
\centering
\caption{Plantilla de ficha experimental.}
\label{tab:ficha-experimento}
\begin{tabularx}{\textwidth}{p{3.8cm}X}
\toprule
Campo & Registro \\
\midrule
ID y versión & \rule{8cm}{0.4pt}\\
Hipótesis & \rule{8cm}{0.4pt}\\
Variable independiente & \rule{8cm}{0.4pt}\\
Variables dependientes & \rule{8cm}{0.4pt}\\
Controles & \rule{8cm}{0.4pt}\\
Conjunto de datos y hash & \rule{8cm}{0.4pt}\\
Configuraciones & \rule{8cm}{0.4pt}\\
Repeticiones & \rule{8cm}{0.4pt}\\
Exclusiones previas & \rule{8cm}{0.4pt}\\
Regla de decisión & \rule{8cm}{0.4pt}\\
Revisión & \rule{8cm}{0.4pt}\\
\bottomrule
\end{tabularx}
\end{table}

\section{Ficha de resultado}

\begin{table}[H]
\centering
\caption{Plantilla de resultado.}
\label{tab:ficha-resultado}
\begin{tabularx}{\textwidth}{p{3.8cm}X}
\toprule
Campo & Registro \\
\midrule
Corridas incluidas & \rule{8cm}{0.4pt}\\
Corridas excluidas & \rule{8cm}{0.4pt}\\
Resumen descriptivo & \rule{8cm}{0.4pt}\\
Tamaño de efecto & \rule{8cm}{0.4pt}\\
Intervalo & \rule{8cm}{0.4pt}\\
Resultado de hipótesis & Apoyada / no apoyada / inconclusa\\
Costo adverso & \rule{8cm}{0.4pt}\\
Amenaza principal & \rule{8cm}{0.4pt}\\
Artefactos & \rule{8cm}{0.4pt}\\
\bottomrule
\end{tabularx}
\end{table}

\section{Ficha de incidente experimental}

Un incidente registra identificador de corrida, instante, etapa, síntoma, impacto, datos afectados, causa
conocida, decisión de exclusión, evidencia y acción preventiva. Un incidente de
aislamiento se reporta aunque la corrida no sea válida para rendimiento.

\section{Declaración de disponibilidad}

La publicación deberá indicar qué artefactos son abiertos, cuáles se sustituyeron por
datos sintéticos y cuáles no pueden compartirse. La imposibilidad de publicar código
privado no impide compartir implementación neutral, generadores, contratos y resultados.

\section{Declaración de autoría}

\ifnum\blindreview=1
La identidad del investigador y autor se omite en esta copia para revisión ciega. Las contribuciones
\else
El investigador y autor de la propuesta es Billy Jak Cordero Porras. Las contribuciones
\fi
de revisores, colaboradores y herramientas se declararán según las políticas de la
institución y de la sede de publicación. La autoría implica responsabilidad sobre
afirmaciones, datos, análisis y correcciones.

\chapter{Matriz de trazabilidad del manuscrito}

\begin{longtable}{>{\bfseries}p{1.4cm}p{4.2cm}p{4.6cm}p{3.6cm}}
\caption{Trazabilidad entre preguntas, artefactos y evaluación.}\label{tab:trazabilidad-manuscrito}\\
\toprule
RQ & Artefacto COBRIA & Experimento & Resultado requerido \\
\midrule
\endfirsthead
\toprule
RQ & Artefacto COBRIA & Experimento & Resultado requerido \\
\midrule
\endhead
RQ1 & Modelo de complejidad & P1--P5, R1--R2 & Curvas y amplificación.\\
RQ2 & Proyecciones de consulta & P1 & Latencia, lecturas y bytes.\\
RQ3 & Presupuesto de proyección & P2--P3 & $A_w$, $A_s$ y equilibrio.\\
RQ4 & Gestor de proyecciones & R1--R5 & Exactitud, tiempo y recuperación.\\
RQ5 & Metadatos y generador & Evolución y generator test & Esfuerzo y divergencia.\\
RQ6 & Ámbito y políticas & Pruebas negativas & Cero fugas y efectos.\\
RQ7 & Perfiles A y B & Mapeo y ejecución comparada & Principios sin ruptura.\\
RQ8 & Conjunto de datos/Feature contracts & AI-1--AI-2 & Linaje y equivalencia.\\
RQ9 & Retrieval, tools y ledger & AI-3--AI-7 & Evidencia y acciones seguras.\\
RQ10 & Lienzo y costo & Todos, incluidos negativos & Región de no adopción.\\
\bottomrule
\end{longtable}

\section{Trazabilidad de partes}

\begin{table}[H]
\centering
\caption{Función de cada parte del documento maestro.}
\label{tab:trazabilidad-partes}
\begin{tabularx}{\textwidth}{p{2cm}p{4.3cm}X}
\toprule
Parte & Función & Producto \\
\midrule
I & Plantear problema y preguntas & Hipótesis y alcance.\\
II & Formalizar artefacto & Modelo, contratos y algoritmos.\\
III & Razonar sobre costo & Complejidad y optimización.\\
IV & Extender a datos e IA & Linaje, RAG, agentes y gobierno.\\
V & Conectar con implementaciones & Casos anónimos y referencia.\\
VI & Diseñar evaluación & Experimentos y análisis.\\
VII & Sintetizar y concluir & Guía, contribuciones y límites.\\
VIII & Facilitar réplica & Plantillas, checklists y diccionario.\\
\bottomrule
\end{tabularx}
\end{table}

\section{Criterio de completitud para una versión publicable}

El manuscrito estará listo para someterse como estudio empírico cuando: la referencia y
los generadores sean públicos; P1, P2, R1, R2, aislamiento, evolución, AI-1 y AI-3 tengan
corridas reproducibles; las hipótesis se actualicen sin eliminar resultados negativos;
otro evaluador revise una muestra; y tablas y figuras provengan de scripts.

Si se publica antes de esos experimentos, debe presentarse como propuesta arquitectónica
y protocolo de investigación, no como validación de rendimiento. Esta distinción debe
aparecer en título, resumen, conclusiones y carta de presentación.

\section{Cierre de los anexos}

Los anexos transforman COBRIA de una descripción conceptual en un artefacto inspeccionable.
Los contratos hacen explícitas sus propiedades; los checklists permiten auditar una
implementación; el diccionario fija métricas; y el paquete reproducible define cómo
convertir futuras ejecuciones en evidencia.

\enlargethispage{2\baselineskip}
La utilidad de estas plantillas depende de su aplicación crítica. Marcar todas las casillas
sin pruebas no aumenta madurez. Cada afirmación debe conservar una ruta hacia código,
configuración, datos, ejecución o revisión que otra persona pueda examinar.
%
}{%
  \part{Nota sobre anexos históricos}
  \chapter{Material no incluido en esta edición}
  El anexo histórico independiente no forma parte del paquete actual. Los procedimientos
  vigentes de réplica y operación se conservan en la parte anterior; no se reconstruyen
  anexos eliminados ni resultados ausentes mediante inferencias editoriales.
}

\backmatter
\printbibliography[heading=bibintoc,title={Referencias}]
\end{document}
